% concatenated from aos-paper.tex
\documentclass[aos,preprint]{imsart}

%% Packages
\RequirePackage{amsthm,amsmath,amsfonts,amssymb}
\RequirePackage[numbers,sort&compress]{natbib}
%\RequirePackage[authoryear]{natbib}%% uncomment this for author-year citations
\RequirePackage[colorlinks,citecolor=blue,urlcolor=blue]{hyperref}
\RequirePackage{graphicx}
\usepackage{lipsum}
\usepackage{amsfonts}
\usepackage{graphicx}
\usepackage{epstopdf}
\usepackage{ragged2e}
% \usepackage{algorithmic}


\usepackage{mathtools}
\usepackage{bm}
\usepackage[mathscr]{euscript}
\usepackage[T1]{fontenc}
\usepackage{enumitem}
\usepackage{xcolor}
\usepackage{multirow}
\usepackage{booktabs}
\usepackage{microtype}
\usepackage{fix-cm}
\usepackage{colortbl}
\usepackage[utf8]{inputenc}
\usepackage{tikz}
\usepackage{pbox}
\usepackage{fourier} 
\usepackage{array}
\usepackage{makecell}

\usepackage{algorithm}
% \usepackage{algpseudocode}

\usepackage{bbding}

\startlocaldefs
%%%%%%%%%%%%%%%%%%%%%%%%%%%%%%%%%%%%%%%%%%%%%%
%%                                          %%
%% Uncomment next line to change            %%
%% the type of equation numbering           %%
%%                                          %%
%%%%%%%%%%%%%%%%%%%%%%%%%%%%%%%%%%%%%%%%%%%%%%
%\numberwithin{equation}{section}
%%%%%%%%%%%%%%%%%%%%%%%%%%%%%%%%%%%%%%%%%%%%%%
%%                                          %%
%% For Axiom, Claim, Corollary, Hypothesis, %%
%% Lemma, Theorem, Proposition              %%
%% use \theoremstyle{plain}                 %%
%%                                          %%
%%%%%%%%%%%%%%%%%%%%%%%%%%%%%%%%%%%%%%%%%%%%%%
\theoremstyle{plain}
\newtheorem{axiom}{Axiom}
\newtheorem{claim}[axiom]{Claim}
\newtheorem{theorem}{Theorem}[section]
\newtheorem{lemma}[theorem]{Lemma}
\DeclareMathOperator{\pr}{\mathbb P}
\DeclareMathOperator{\E}{\mathbb E}
\DeclareMathOperator{\Var}{Var}
\DeclareMathOperator{\Cov}{Cov}
\DeclareMathOperator{\Corr}{Corr}
\DeclareMathOperator{\MSE}{MSE}
\DeclareMathOperator{\ind}{\mathbb I}

\newcommand{\dif}{\mathop{}\!\mathrm{d}}
\newcommand{\Real}{\mathbb R}
\newcommand{\Int}{\mathbb Z}
\newcommand{\NatInt}{\mathbb N}
\newcommand{\SFM}{\mathsf M}
\newcommand{\SFY}{\mathsf Y}
\newcommand{\SFZ}{\mathsf Z}
\newcommand{\CalO}{\mathcal O}
\newcommand{\CalT}{\mathcal T}
\newcommand{\CalH}{\mathcal H}
\newcommand{\CalX}{\mathcal X}
\newcommand{\CalP}{\mathcal P}
\newcommand{\CalU}{\mathcal U}
\newcommand{\CalK}{\mathcal K}
\newcommand{\BX}{\bold X}
\newcommand{\BI}{\bold I}
\newcommand{\Bx}{\bold x}
\newcommand{\Bs}{\bold s}
\newcommand{\BFX}{\bold X}
\newcommand{\BFx}{\bold x}
\newcommand{\BB}{\bold B}
\newcommand{\By}{\bold y}
\newcommand{\BU}{\bold U}
\newcommand{\Bu}{\bold u}
\newcommand{\BK}{\bold K}
\newcommand{\im}{\text{i}}
\newcommand{\BPhi}{\boldsymbol{\Phi}}
\newcommand{\CalF}{\mathcal{F}}
\newcommand{\CalW}{\mathcal{W}}
\newcommand{\Bomega}{\boldsymbol{\omega}}
\newcommand{\Balpha}{\boldsymbol{\alpha}}
\newcommand{\Bbeta}{\boldsymbol{\beta}}
\newcommand{\QED}{\hfill $\square$}
\newcommand{\TBC}{\textcolor{red}{To be completed...}}
\newcommand{\cmtS}[1]{{\color{red}{(Simon: #1)}}}
\newcommand{\cmtL}[1]{{\color{orange}{(Leon: #1)}}}

\renewcommand\theadalign{bc}
\renewcommand\theadfont{\bfseries}
\renewcommand\theadgape{\Gape[4pt]}
\renewcommand\cellgape{\Gape[4pt]}

\newtheorem{defn}{Definition}
\newtheorem{remark}{Remark}
\newtheorem{lem}{Lemma}
\newtheorem{corollary}{Corollary}
\newtheorem{prop}{Proposition}

%%%%%%%%%%%%%%%%%%%%%%%%%%%%%%%%%%%%%%%%%%%%%%
%%                                          %%
%% For Assumption, Definition, Example,     %%
%% Notation, Property, Remark, Fact         %%
%% use \theoremstyle{remark}                %%
%%                                          %%
%%%%%%%%%%%%%%%%%%%%%%%%%%%%%%%%%%%%%%%%%%%%%%
\theoremstyle{remark}
\newtheorem{definition}[theorem]{Definition}
\newtheorem*{example}{Example}
\newtheorem*{fact}{Fact}
%%%%%%%%%%%%%%%%%%%%%%%%%%%%%%%%%%%%%%%%%%%%%%
%% Please put your definitions here:        %%
%%%%%%%%%%%%%%%%%%%%%%%%%%%%%%%%%%%%%%%%%%%%%%

\endlocaldefs

\begin{document}

\begin{frontmatter}
\title{The BdryMat\'ern GP: Reliable incorporation of boundary information on irregular domains for Gaussian process modeling}
%\title{A sample article title with some additional note\thanksref{t1}}
\runtitle{BdryMat\'ern GP}
%\thankstext{T1}{A sample additional note to the title.}

\begin{aug}
%%%%%%%%%%%%%%%%%%%%%%%%%%%%%%%%%%%%%%%%%%%%%%%
%% Only one address is permitted per author. %%
%% Only division, organization and e-mail is %%
%% included in the address.                  %%
%% Additional information can be included in %%
%% the Acknowledgments section if necessary. %%
%% ORCID can be inserted by command:         %%
%% \orcid{0000-0000-0000-0000}               %%
%%%%%%%%%%%%%%%%%%%%%%%%%%%%%%%%%%%%%%%%%%%%%%%
\author[A]{\fnms{Liang}~\snm{Ding}\ead[label=e1]{liang\_ding@fudan.edu.cn}},
\author[B]{\fnms{Simon}~\snm{Mak}\ead[label=e2]{sm769@duke.edu}}
\and
\author[C]{\fnms{C. F. Jeff}~\snm{Wu}\ead[label=e3]{jeffwu@cuhk.edu.cn}}
%%%%%%%%%%%%%%%%%%%%%%%%%%%%%%%%%%%%%%%%%%%%%%
%% Addresses                                %%
%%%%%%%%%%%%%%%%%%%%%%%%%%%%%%%%%%%%%%%%%%%%%%
\address[A]{School of Data Science,
Fudan University\printead[presep={ ,\ }]{e1}}

\address[B]{Department of Statistical Science,
Duke University\printead[presep={,\ }]{e2}}
\address[C]{School of Data Science,
The Chinese University of Hong Kong (Shenzhen)\printead[presep={,\ }]{e3}}
\end{aug}

\begin{abstract}
Gaussian processes (GPs) are broadly used as surrogate models for expensive computer simulators of complex phenomena. However, a key bottleneck is that its training data are generated from this expensive simulator and thus can be highly limited. A promising solution is to supplement the learning model with boundary information from scientific knowledge. However, despite recent work on boundary-integrated GPs, such models largely cannot accommodate boundary information on irregular (i.e., non-hypercube) domains, and do not provide sample path smoothness control or approximation error analysis, both of which are important for reliable surrogate modeling. We thus propose a novel BdryMat\'ern GP modeling framework, which can reliably integrate Dirichlet, Neumann and Robin boundaries on an irregular connected domain with a boundary set that is twice-differentiable almost everywhere. Our model leverages a new BdryMat\'ern covariance kernel derived in path integral form via a stochastic partial differential equation formulation. Similar to the GP with Mat\'ern kernel, we prove that sample paths from the BdryMat\'ern GP satisfy the desired boundaries with smoothness control on its derivatives. We further present an efficient approximation procedure for the BdryMat\'ern kernel using finite element modeling with rigorous error analysis. Finally, we demonstrate the effectiveness of the BdryMat\'ern GP in a suite of numerical experiments on incorporating broad boundaries on irregular domains.

% In many problems, additional boundary information (i.e., the behavior of the phenomena along input boundaries) is known beforehand, either from governing physics or scientific knowledge. In this work, we propose a framework of GP models based on stochastic Partial Differential Equations (SPDEs) with Dirichlet and Robin boundary conditions. For general SPDEs, we show that their covariances can be represented in a path integral form. Based on this path integral form, we provide a regression approximation, called regression kernel, along with the necessary criteria to accurately approximate the SPDE covariances. We also show that for covariances in tensor form, closed-form solutions are available. Numerical experiments are conducted to illustrate that incorporating boundary information can significantly improve tperformance of the model.
\end{abstract}

\begin{keyword}[class=MSC]
\kwd[Primary ]{62M40}
\kwd[; secondary ]{62G08}
\end{keyword}

\begin{keyword}
\kwd{Bayesian Nonparametrics}
\kwd{Computer Experiments}
\kwd{Gaussian Processes}
\kwd{Partial Differential Equations}
\kwd{Uncertainty Quantification}
\end{keyword}

\end{frontmatter}
%%%%%%%%%%%%%%%%%%%%%%%%%%%%%%%%%%%%%%%%%%%%%%
%% Please use \tableofcontents for articles %%
%% with 50 pages and more                   %%
%%%%%%%%%%%%%%%%%%%%%%%%%%%%%%%%%%%%%%%%%%%%%%
%\tableofcontents

\section{Introduction}

With groundbreaking advances in scientific modeling and computing, complex phenomena (e.g., rocket propulsion and particle collisions) can now be reliably simulated via computer code, which typically solve a sophisticated system of partial differential equations (PDEs). Such ``computer experiments'' have had demonstrated success in advancing broad scientific and engineering applications, including particle physics \citep{sengupta2025hybrid,liyanage2022efficient}, aerospace engineering \citep{mak2018efficient,miller2024expected} and robotics \citep{choi2021use,liu2025quip}. A key bottleneck of such virtual experiments, however, is that each run demands a large amount of computing power. For example, a single simulation of a heavy-ion particle collision and its subsequent subatomic interactions can require over thousands of CPU hours \citep{ji2024conglomerate}. Given a limited computational budget, this makes the exploration of the simulation response surface over an input space $\mathcal{X}$ a highly challenging task.

A proven solution is probabilistic surrogate modeling \citep{gramacy2020surrogates}. The idea is simple but effective. First, one runs the expensive simulator (denoted $f$) on a set of $n$ designed input points on $\mathcal{X} \subseteq \mathbb{R}^d$, where $\mathcal{X}$ is the compact domain of its $d$ input parameters. Second, using such data, one trains a probabilistic model that can predict with uncertainty the simulator output $f(\mathbf{x}_{\rm new})$ at an untested input $\mathbf{x}_{\rm new} \in \mathcal{X}$, thus bypassing an expensive simulation run. Gaussian processes (GPs; \citep{gramacy2020surrogates}) and their extensions are widely used as surrogate models as they provide a closed-form predictive distribution on $f(\cdot)$, which facilitates timely downstream scientific decision-making, e.g., optimization \citep{chen2024hierarchical} and inverse problems \citep{ehlers2024bayesian}. In practice, one faces a further obstacle of limited sample sizes (i.e., small $n$) for surrogate training, due to the costly nature of the expensive computer simulator. For complex applications, e.g., in aerospace engineering \citep{narayanan2024misfire}, this can restrict the experimenter to only tens of runs over the domain $\mathcal{X}$, which results in mediocre surrogate performance.

To address this, recent work has explored the integration of known ``physics'' from scientific knowledge within surrogate models; this field of ``physics-integrated machine learning'' \citep{willard2020integrating} is a rising and promising area. For GP surrogates, this includes the incorporation of known shape constraints \citep{WangBerger16,Golchi15}, mechanistic equations \citep{Wheeler14}, manifold embeddings \citep{li2023additive,seshadri2019dimension} and boundary information \citep{tan2018gaussian,vernon2019known,jackson2023efficient,ding2019bdrygp}; we focus on the last topic in this paper. Here, boundary information refers to known information on $f$ along the boundaries of its domain $\mathcal{X}$. Such knowledge may be elicited via a careful analysis of the underlying PDE system or via scientific intuition. Take, e.g., the bending of a metal beam under a fixed force, and consider its displacement $f(\mathbf{x})$ as either beam inverse thickness ($x_1$) and beam length ($x_2$) varies. As thickness grows large (i.e., $x_1 \rightarrow 0$) or length becomes small (i.e., $x_2 \rightarrow 0$), intuition suggests that its displacement $f(\mathbf{x})$ should approach zero; such boundary information can be integrated for surrogate training. Similar information can be elicited in the simulation of viscous flows \cite{white2006viscous}, which are used in climatology and high energy physics. Such flows are dictated by the complex Navier-Stokes equations \cite{temam2001navier}, but at the limits of inputs (e.g., zero viscosity or fluid incompressibility), these equations may admit closed-form solutions \cite{kiehn2001some,humphrey2016introduction}. The integration of this boundary information can greatly improve surrogate modeling with limited simulation runs.

There is a growing literature on integrating boundary information within GPs. Here, Dirichlet, Neumann and Robin boundaries (see \cite{zwillinger2021handbook}) refer to knowledge of $f$, its derivative, or a linear combination of both along the boundaries of $\mathcal{X}$, respectively. \cite{tan2018gaussian, li2022improving} proposed a boundary-modified GP, which modifies the mean and covariance structure of a stationary GP to incorporate Dirichlet boundaries on discrete points. \cite{vernon2019known} and \cite{jackson2023efficient} adopted a projection-based approach for integrating Dirichlet boundaries on axis-symmetric and/or axis-parallel boundaries on a hypercube domain $\mathcal{X}$. \cite{gulian2022gaussian,solin2019know,solin2020hilbert} advocated for a low-rank kernel approximation strategy, where its underlying basis functions explicitly satisfy the prescribed Dirichlet boundary conditions. Finally, \cite{dalton2024boundary} makes use of so-called ``approximate distance functions'' (ADFs) for constructing kernels that can integrate Dirichlet, Robin or Neumann boundaries on non-hypercube domains.

Such models, however, have several limitations. First, aside from \cite{dalton2024boundary}, they consider only the integration of Dirichlet boundaries and not Neumann or Robin boundaries, which can often be elicited in applications \cite{dalton2024boundary}. Second, existing models do not provide a characterization of smoothness (i.e., differentiability) for its resulting sample paths. As such, they do not permit \textit{smoothness control} of the fitted probabilistic model, which is important for effective surrogate modeling \cite{santner2003}. Finally, the input space $\mathcal{X}$ can be ``irregular'', i.e., it takes the form of a non-hypercube domain; see Figure \ref{fig:task} (left) for an example. Such irregular domains are broadly encountered in physical science applications, including glaciology \citep{pratola2017design} and fluid dynamics \cite{mak2018efficient,dalton2024boundary}. Incorporating boundary information on irregular domains requires a careful approximation approach with rigorous \textit{error analysis} for reliable surrogate modeling, which existing methods do not provide. For modern applications that feature rich boundary information on irregular input spaces, there is thus a need for a theoretically grounded GP framework that reliably integrates such information with smoothness and approximation error control for effective surrogate modeling.

The proposed BdryMat\'ern GP targets such a framework (see Figure \ref{fig:task} for a visualization). Recall that a key appeal of the Mat\'ern GP, i.e., a GP with (isotropic) Mat\'ern kernel, is its control of sample path differentiability via a smoothness parameter $\nu$ \citep{stein1999}. To extend this for the boundary-integrated setting, we begin with the stochastic partial differential equation (SPDE) representation of a Mat\'ern GP \citep{LindgrenRueLindstrom11}, and impose the known boundary information on domain $\mathcal{X}$ as a constraint on this SPDE system. Provided $\mathcal{X}$ is connected with a boundary set that is twice-differentiable almost everywhere, we can then solve this system to derive the proposed BdryMat\'ern covariance kernel in path integral form using the well-known Feynman-Kac formula \citep{ito1957fundamental,stroock1971diffusion,papanicolaou1990probabilistic} for elliptical PDEs. With this new BdryMat\'ern kernel, we prove that sample paths from the BdryMat\'ern GP enjoy similar smoothness control as the Mat\'ern GP while satisfying the desired boundaries. We further present an efficient approach for approximating the BdryMat\'ern kernel on an irregular domain $\mathcal{X}$ via finite element modeling (FEM; \cite{huebner2001finite}) with rigorous error analysis. We finally demonstrate in a suite of numerical experiments the effectiveness of the BdryMat\'ern GP in reliably integrating broad boundary information on irregular domains.

% This code numerically solves a system of governing equations representing the underlying science of the problem. Due to the time-intensive nature of these numerical simulations \cite{yeh2018common}, Gaussian Processes (GPs; \cite{sacks1989}) are often used as surrogate models to emulate the expensive computer code. Let \(\mathbf{x} \in \mathcal{X}\) be a vector of \(d\) code inputs residing in a general compact domain, and let \(f(\mathbf{x})\) be the corresponding code output. The idea is to adopt a GP prior for \(f(\cdot)\), then use the posterior process given data to infer code output at an unobserved input. GP emulators are now widely used to study a broad range of scientific and engineering problems, such as rocket engines \cite{Mak18}, universe expansions \cite{kaufman2011efficient}, and high energy physics \cite{goh2013prediction}.


% In many applications, there is additional knowledge on the phenomenon than simply computer code output, and incorporating such knowledge can improve GP predictive performance. This ``physics-integrated'' GP modeling has garnered much attention in recent years \cite{Wheeler14,Golchi15,WangBerger16}. 
% % This notion of ``physics-integrated'' GP modeling has garnered much attention in the past few years. One flavor is the monotone GPs studied in \cite{Golchi15}, \cite{WangBerger16} and \cite{Matt17}, which capture known monotone behavior in the black-box computer code. \cite{Wheeler14} proposed a mechanistic hierarchical GP model which incorporates information from a linear ordinary differential equation. Recently, \cite{MattnLi18} proposed a GP model for parabolic PDEs via the principle of superposition.
% We consider here three types of information called \textit{Dirichlet boundaries},  \textit{Neumann boundaries}, and \textit{Robin boundaries}\cite{bazilevs2007weak}, which specifies the values of $f$ and its derivative along certain input boundaries. Boundary conditions are often available from governing physics or from simple physical considerations \cite{Matt18}. One example is the simulation of viscous flows \cite{white2006viscous}, widely used in climatology and high energy physics. Such flows are dictated by the complex Navier-Stokes equations \cite{temam2001navier}, and can be very time-consuming to simulate. At the limits of certain variables (e.g., zero viscosity or fluid incompressibility), the Navier-Stokes equations can be greatly simplified for efficient, even closed-form, solutions \cite{kiehn2001some,humphrey2016introduction}. Incorporating this boundary information within the GP can allow for improved predictive performance.

% \cmtS{this needs a more thorough literature review: there has been recent work on boundary-integrated GPs in the literature (see work by Ian Vernon and others). also needs discussion on prior work for physics-integrated GPs. i can add a couple of paragraphs on this.} Despite its promise, the integration of GPs with boundary information is largely unexplored in the literature, with the only reference being a recent paper by \cite{Matt18}. In this paper, a flexible \textit{Boundary Modified Gaussian Process} (BMGP) is proposed, which can integrate a broad range of boundaries by modifying the mean and variance structure of a stationary GP. However, unlike commonly used GPs with Matérn and Gaussian kernels, the GPs with boundary information constructed in \cite{Matt18}, though flexible, lack interpretability, which is also a key ingredient in uncertainty quantification. In this study, we extend the SPDE representation of integer-order Matérn GPs on \(\mathbb{R}^2\) from \cite{LindgrenRueLindstrom11} to compact domains with irregular boundary conditions and general Matérn GPs. Additionally, we propose efficient methods for computing the GPs' posteriors.

This paper is organized as follows: Section \ref{sec:GP_review} provides an overview of GPs and existing work on boundary-integrated GPs. Section \ref{sec:Bdrymat_SPDE} derives the new BdryMat\'ern kernel in path integral form, and proves smoothness control on sample paths from its corresponding BdryMat\'ern GP. Section \ref{sec:kernel_approx} outlines a finite-element-method approach for kernel approximation with supporting error analysis. Section \ref{sec:tensor_mat} presents a tensor form of the BdryMat\'ern GP that provides closed-form kernels on the unit hypercube domain. Section \ref{sec:numeric} demonstrates the effectiveness of the BdryMat\'ern GP in numerical experiments, and Section \ref{sec:conclusion} concludes the paper.

\begin{figure}[!t]
\begin{center}
\includegraphics[width=\textwidth]
{figures/workflow}
\end{center}
\caption{Visualizing the proposed BdryMat\'ern GP for reliable incorporation of boundary information for probabilistic surrogate modeling. Figures in this workflow adapted from \cite{dalton2024boundary}, \cite{mak2018efficient} and \cite{ji2024graphical}.}
\label{fig:task}
\end{figure}

\section{Background \& Motivation}
\label{sec:GP_review}
We first provide in this section a brief overview on GP modeling, the isotropic Mat\'ern kernel and its representation as an SPDE. We then review existing GP models that integrate boundaries, and their potential limitations.

\subsection{GP Modeling and the Mat\'ern Kernel}
Let $\Bx \in \mathcal{X}\subseteq\Real^d$ be an input vector lying on a compact domain $\mathcal{X}$, and let $f(\Bx)$ denote its black-box output from the expensive computer simulator. We adopt the following Gaussian process model on $f$:
\begin{equation}
    f(\cdot) \sim \text{GP}\{\mu(\Bx),k(\Bx,\Bx')\}.
\end{equation}
Here, $\mu(\Bx)$ is its mean function, and $k(\Bx,\Bx')$ is its covariance function or kernel. With little prior information on $f$, $\mu(\Bx)$ is typically specified as a constant $\mu$ (to be estimated from data), and $k(\Bx,\Bx')$ is selected to ensure a desirable degree of smoothness for GP sample paths (e.g., the Mat\'ern kernel later). Given known boundaries on $f$, however, both $\mu(\Bx)$ and $k(\Bx,\Bx')$ should be carefully specified so that its sample paths satisfy such boundaries. We will later denote this \textit{boundary-integrated} mean and covariance function as $\mu_{\mathcal{B}}(\mathbf{x})$ and $k_{\mathcal{B}}(\Bx,\Bx')$, respectively.

% Let $\Bx \in \mathcal{X}\subseteq\Real^d$ be an input vector on domain $\mathcal{X}$, with $f(\Bx)$ denoting its corresponding computer code output. In Gaussian process emulation \cite{sacks1989,santner2003}, $f$ is assumed to be a realization of a Gaussian process with mean function $\mu: \mathcal{X} \to \Real$, and covariance function $k: \mathcal{X} \times \mathcal{X} \to \Real$. Further details on GP modeling can be found in \cite{Adler81}.

Next, suppose the computer simulator is run at a set of $n$ designed input points $\BX =\{\Bx_1,\cdots,\Bx_n\}\subset \mathcal{X}$, yielding observations $f(\BX)=[f(\Bx_1),\cdots,f(\Bx_n)]^T$. In what follows, we presume that the simulator $f$ is \textit{deterministic}, meaning the data $f(\BX)$ are observed without noise. This is often assumed in the computer experiments literature \citep{santner2003}, when the simulator solves a deterministic PDE system and returns the same solution $f(\Bx)$ given the same inputs $\Bx$. One can easily accommodate for Gaussian simulation noise via the use of a nugget term \citep{peng2014choice} in the below predictive equations.

Conditional on data $f(\BX)$, the posterior predictive distribution of $f(\Bx_{\rm new})$ at an unevaluated input point $\Bx_{\rm new}$ takes the form of a Gaussian distribution $[f(\mathbf{x}_{\rm new})|f(\mathbf{X})] \sim \mathcal{N}\{\hat{f}_n(\Bx_{\rm new}),k_n(\Bx_{\rm new},\Bx_{\rm new})\}$, with posterior mean and variance given by:
\begin{align}
\begin{split}
    \hat{f}_n(\Bx)&=\mu(\Bx)+\mathbf{k}(\Bx,\BX)\mathbf{K}^{-1}(\BX,\BX)\left[f(\BX)-\mu(\BX)\right],\\
    k_n(\Bx,\Bx')&=k(\Bx,\Bx')-\mathbf{k}(\Bx,\BX)\mathbf{K}^{-1}(\BX,\BX)\mathbf{k}(\BX,\Bx').
    \end{split}
    \label{eq:kriging}
\end{align}
Here, $\mathbf{k}(\BX,\Bx_{\rm new}) = [k(\Bx_{\rm new},\Bx_1),\cdots,k(\Bx_{\rm new},\Bx_n)]^T$ is the covariance vector between the design points $\BX$ and a new point $\Bx_{\rm new}$, $k(\Bx,\BX) = k(\BX,\Bx)^T$, and $\mathbf{K}(\BX,\BX) = {{[k(\Bx_i},\Bx_j)]_{i=1}^n}_{j=1}^n$ is the covariance matrix over design points. Equation \eqref{eq:kriging} provides the means for ``emulating'' the expensive simulator $f$ with probabilistic uncertainty quantification over the domain $\mathcal{X}$.

% The posterior mean $\hat{f}_n$ is typically used as a predictor (or emulator) for unknown  output $f$, since it is optimal under quadratic and absolute error loss \cite{santner2003}. The posterior variance $k_n(\Bx,\Bx)$ then quantifies the uncertainty of the predictor $\hat{f}_n(\Bx)$ at  $\Bx$.
% We will denote this predictor as $\hat{f}_n(\cdot) = \hat{\mu}_n(\cdot)$. 

For GP modeling, a popular covariance function choice is the (isotropic) Mat\'ern kernel \citep{stein1999}, given by:
\begin{equation}
    \label{eq:matern}
    k_{\nu}(\BFx,\BFx')=\frac{\sigma^2}{2^{\nu-1}\Gamma(\nu)}(\kappa\|\BFx-\BFx'\|)^\nu K_\nu(\kappa\|\BFx-\BFx'\|).
\end{equation}
Here, $\nu > 0$ is its smoothness parameter, $\kappa >0$ is its inverse length-scale parameter, $\sigma^2 > 0$ is its variance parameter, and $K_\nu$ is the modified Bessel function of the second kind \citep{abramowitz1965handbook}. A key appeal for the Mat\'ern kernel with smoothness $\nu$ is that its sample path smoothness can be controlled by $\nu$; in particular, its sample paths are \textup{($\lceil\nu\rceil$-1)}-differentiable almost surely \citep{da2023sample}. Such smoothness control is desirable in broad scientific applications, e.g., climatology \citep{wiens2020modeling} and geostatistics \citep{pardo2008geostatistics}. It is particularly important for our motivating surrogate modeling application, where the black-box function $f$ (which often represents the solution of a PDE system) may have known smoothness, i.e., differentiability, properties \citep{Evans15}. Integrating this smoothness information via a careful selection of $\nu$ in the Mat\'ern kernel permits interpretable surrogate modeling with improved predictive performance \citep{stein1999}.

% the sample paths of a GP with Ma\'tern-$\nu$ kernel is almost-$\nu$ H\"older \cite{da2023sample}

Consider now $\mathcal{W}_{\nu}(\Bx)$, the zero-mean Mat\'ern GP with smoothness parameter $\nu$, inverse length-scale parameter $\kappa$, and process variance $\sigma^2 = \Gamma(\nu)[\Gamma(\alpha)(4\pi)^{d/2}\kappa^{2\nu}\tau^2]^{-1}$, where $\alpha=\nu+d/2$ and $\tau^2 > 0$ is given. (For brevity, all parameters except for $\nu$ are suppressed in the notation $\mathcal{W}_{\nu}(\Bx)$.) One can show (see, e.g., \cite{whittle1954stationary,whittle1963stochastic,LindgrenRueLindstrom11}) that $\mathcal{W}_{\nu}(\Bx)$ is the unique stationary solution to the following SPDE system:
\begin{equation}
    \label{eq:SPDE}
    \tau (\kappa^2-\bigtriangleup)^{\alpha/2}f(\BFx)=\mathcal{W}(\Bx),\quad \BFx\in\Real^d.
\end{equation}
Here, $\CalW(\Bx)$ is a Gaussian white noise process with unit spectral, and $\bigtriangleup=\sum_{j=1}^d\partial_{x_j^2}$ is the Laplacian operator. This SPDE representation is leveraged in various contexts for scalable computation and/or approximation, e.g., within Integrated Nested Laplace Approximations (INLA; \cite{rue2009approximate}). Furthermore, from \eqref{eq:SPDE}, it follows that $\mathcal{W}_{\nu+2}(\Bx)$ is a solution to the SPDE:
\begin{equation}
    (\kappa^2-\bigtriangleup)f(\BFx)=\mathcal{W}_\nu(\Bx),\quad \BFx\in\Real^d.
    \label{eq:SPDE2}
\end{equation}
We will use a modification of this representation later to derive the BdryMat\'ern GP model.

% and  $(\kappa^2-\bigtriangleup)^{\nu+d/2}$ is the fractional Laplacian operator defined via Fourier transform:
% \[\mathcal{F}[(\kappa^2-\bigtriangleup)^{\alpha/2} \phi](\Bomega)=(\kappa^2+\|\Bomega\|^2)^{\alpha/2}\CalF[\phi](\Bomega)\]
% for any function $\phi$ with well-defined Fourier transform.

% \cite{LindgrenRueLindstrom11} showed that GPs satisfying \eqref{eq:SPDE} can be approximated by \emph{Markov random fields} on lattices or triangulations in \(\mathbb{R}^2\) with sparse precision matrices. However, this method does not extend to GPs with boundary conditions, as boundary effects disrupt the stationary assumption crucial to the approach.

%\cmtS{i think we need a subsection introducing the different types of boundaries we aim to incorporate. along with motivating applications where such boundaries are known and available. otherwise it is unclear why our method \& theory are important. i can add this in the next few days.}\\
%\cmtL{Jeff suggests moving the discussion on the importance of the method to Section 1, as illustrated in Figure \ref{fig:task}.}


\subsection{Existing Work on Boundary-Integrated GPs}

Next, we briefly review existing literature on boundary-integrated GPs, and highlight their potential limitations. This literature can be broadly categorized into projection-based, eigenfunction-based and ADF-based approaches; we discuss each of these categories below.

% When boundary conditions are imposed on the SPDE in \eqref{eq:SPDE}, the resulting kernel function becomes non-stationary and typically has no closed-form expression. This complication arises because stationarity—central to most standard GP constructions—is no longer valid once the solution must satisfy prescribed conditions at the domain boundaries. Consequently, current methods for constructing GP models with boundary conditions largely focus on desiging the kernel functions so that they inherently respect these boundary constraints, often by embedding the PDE’s structural or continuity requirements into the kernel itself.  We review three widely used approaches, which we refer to as the \emph{projection approach, eigenfunction approach}, and \emph{Green's function approach}, respectively.

\textbf{Projection-based approaches:} One class of boundary-integrated GPs (see \cite{vernon2019known,jackson2023efficient}) makes use of careful projections onto known boundaries on the input space $\mathcal{X}$. Suppose $\mathcal{X} = [0,1]^d$ is the unit hypercube, and further suppose we know the Dirichlet boundaries of $f$ along the $(d-1)$-dimensional hyperplane $\mathcal{B} = \{\Bx: x_1 = 0\}$. Let $f(\cdot) \sim \text{GP}\{\mu,k(\Bx,\Bx')\}$, where the covariance kernel takes the product form $k(\mathbf{x},\mathbf{x}') = \sigma^2 \prod_{l=1}^d k_l(x_l,x_l')$. Consider the prediction of $f(\mathbf{x})$ at a new point $\mathbf{x}$, conditional on evaluations of $f$ on a finite set of points $\mathbf{B} = \{\mathbf{b}_1, \cdots, \mathbf{b}_m,\Bx_{\mathcal{B}}\} \subset \mathcal{B}$ on the boundary, where $\Bx_{\mathcal{B}}$ is the projection of $\Bx$ onto $\mathcal{B}$. Then \cite{vernon2019known} shows that the posterior predictive distribution $f(\mathbf{x})|f(\mathbf{B})$ takes the form $\mathcal{N}\{\mu_{\mathcal{B}}(\mathbf{x}),k_{\mathcal{B}}(\Bx,\Bx)\}$, where:
\begin{align}
\begin{split}
\mu_{\mathcal{B}}(\mathbf{x}) &= \mu + k_1(x_1)[f(\Bx_{\mathcal{B}}) - \mu],\\
k_{\mathcal{B}}(\Bx,\Bx')&= \sigma^2 [k_1(x_1-x_1')-k_1(x_1)k_1(x_1')] \prod_{l=2}^d k_l(x_l,x_l').
\label{eq:proj}
\end{split}
\end{align}

Note that the above predictive equations depend only on the projected point $\Bx_{\mathcal{B}}$ and not the other boundary points $\mathbf{b}_1, \cdots, \mathbf{b}_m$. Thus, the use of the projected point $\Bx_{\mathcal{B}}$ in the boundary-integrated mean function and an appropriately modified covariance function is sufficient for capturing boundary information in this setting. One can then directly use $\mu_{\mathcal{B}}$ and $k_{\mathcal{B}}$ within the GP equations \eqref{eq:kriging} for probabilistic surrogate modeling. This projection-based framework can be extended (see \cite{jackson2023efficient}) to accommodate multiple hyperplane boundaries of varying dimensions on the unit hypercube domain $\mathcal{X} = [0,1]^d$.

Despite its elegance, there are several limitations for such approaches. First, it is unclear whether projection-based methods extend for the integration of Neumann and Robin boundaries, which are present in many problems \citep{dalton2024boundary}. Second, the above framework may not extend for the case of irregular input spaces $\mathcal{X}$, which arises in complex applications. We will address both limitations via the proposed BdryMat\'ern GP.

% \cite{vernon2019known,jackson2023efficient} discretize the boundary using sample points \(\{\mathbf{c}_i\}\) on the boundary and design non-stationary kernel functions of product form or additive form
% \[
% k(\mathbf{x}, \mathbf{x}') 
% = \sigma(\mathbf{x}) \,\sigma(\mathbf{x}') \,\bar{k}(\mathbf{x}, \mathbf{x}')\quad \text{product form},
% \]
% \[
% k(\mathbf{x}, \mathbf{x}') 
% = \bar{k}(\mathbf{x}, \mathbf{x}')- r(\BFx,\BFx)\quad  \text{additive form},
% \]
% where  \(\bar{k}\) represents any kernel function suitable for the underlying problem, \(\sigma(\mathbf{x})\) and \(r(\BFx, \BFx')\) are functions that depend on the projections of points \(\BFx\) and \(\BFx'\) onto the boundary, and \(\{\mathbf{c}_i\}\) denote the boundary points. By appropriately selecting \(\sigma(\mathbf{x})\) and \(r(\BFx, \BFx')\), the kernel function can be constructed to satisfy the Dirichlet boundary conditions at each \(\mathbf{c}_i\), ensuring that \(k(\mathbf{c}_i, \BFx) = 0\) for any \(\BFx\) and boundary points \(\mathbf{c}_i\). However, since finding \(\sigma(\mathbf{x})\) and \(r(\BFx, \BFx')\) that enforce derivative-based boundary conditions on \(k\) is challenging, this approach is generally limited to handling Dirichlet constraints.

\textbf{Eigenfunction-based approaches:} Another class of boundary-integrated GPs (see \cite{solin2019know,solin2020hilbert,gulian2022gaussian}) makes use of an eigenfunction approach for kernel construction. The key idea is to employ basis functions that inherently satisfy the provided boundary conditions. Here, such works consider the setting of zero (i.e., homogeneous) boundary conditions, thus $\mu_{\mathcal{B}}(\mathbf{x})$ is set as 0. The covariance kernel $k$ is then constructed as:
\[
k_{\mathcal{B}}(\mathbf{x},\mathbf{x}') = \sum_{j=1} \lambda_j \phi_j(\mathbf{x}) \phi_j(\mathbf{x}'),
\]
where \(\{\phi_j\}\) are basis functions that satisfy known boundary conditions. For rectangular domains $\mathcal{X}$, these basis functions can be taken as the eigenfunctions of the Laplace operator under Dirichlet boundary conditions, which can be derived in closed form \citep{dalton2024boundary}. 

There are two limitations of such eigenfunction-based approaches for boundary-integrated surrogates. First, the derivation of closed-form eigenfunctions is possible only under restricted boundary conditions on simple domains $\mathcal{X}$, and thus may not be possible in broad applications. Second, there is little work investigating smoothness control (of GP sample paths) for such approaches. As mentioned previously, this smoothness control is desirable for surrogate modeling, as one may have prior knowledge on the smoothness of $f$ from the considered PDE system.

% This approach ensures that the resulting Gaussian process model naturally respects boundary constraints while maintaining computational efficiency through low-rank approximations.

\textbf{ADF-based approaches:} The last class of boundary-integrated GPs makes use of the so-called ``approximate distance functions''. Such a term was coined by \cite{dalton2024boundary}, but related concepts have been explored in \cite{tan2018gaussian,li2022improving}. The key idea is to construct the boundary-integrated kernel $k_{\mathcal{B}}$ by modifying a base kernel $k$ (e.g., the Mat\'ern kernel) by a distance function $\varsigma(\Bx)$:
\begin{equation}
k_{\mathcal{B}}(\mathbf{x},\mathbf{x}') = \varsigma(\Bx) \varsigma(\Bx') k(\mathbf{x},\mathbf{x}').
\label{eq:adf}
\end{equation}
Here, $\varsigma(\Bx)$ should be specified such that $\lim_{\Bx \rightarrow \mathbf{b}} \varsigma(\Bx) = 0$ for any boundary point $\mathbf{b} \in \mathcal{B}$. This condition ensures that the prior variance of the GP, i.e., $k_{\mathcal{B}}(\mathbf{x},\mathbf{x})$, converges to zero as $\Bx$ approaches a known boundary, thus reflecting a modeler's certainty of boundary information. Such distance functions can further be used to construct a mean function $\mu_{\mathcal{B}}(\mathbf{x})$ that interpolates known boundaries. \cite{tan2018gaussian} makes use of a product distance function that incorporates Dirichlet boundaries at discrete points. \cite{dalton2024boundary} employs the so-called $R$-function method \citep{rvachev1995r} to construct ADFs for incorporating Dirichlet boundaries on irregular domains $\mathcal{X}$. The same paper further extends this approach (with appropriate modifications on \eqref{eq:adf}) for incorporating Neumann and Robin boundaries.

While such approaches offer a flexible framework for boundary integration, they share a similar limitation with earlier methods: they do not provide a means for investigating smoothness control of the resulting GP sample paths, which is desirable for surrogate modeling. Furthermore, given the need for approximation with complex irregular boundaries, these approaches do not provide analysis of such approximation error. This error analysis is useful in surrogate modeling, where it is important to quantify and aggregate all sources of uncertainty for confident scientific inference. The proposed BdryMat\'ern GP presented next addresses these needs for smoothness control and approximation error analysis.

% utilized the closure property of Gaussian processes (GPs) under linear operators \citep{Adler81}. Consider a linear operator \(\mathcal{L}\) acting on \(u\) with boundary constraints \(B[u] = 0\) defined over a domain \(\mathcal{X}\). The Green’s function \(G(\mathbf{x}, \mathbf{x}')\) for \(\mathcal{L}\) is given by:

% \[
% \mathcal{L} \int_{\mathcal{X}} G(\cdot, \mathbf{x})\,f(\mathbf{x})\,d\mathbf{x} = f, 
% \quad G(\mathbf{x}', \mathbf{x}) = G(\mathbf{x}, \mathbf{x}'),
% \quad B\bigl[G(\cdot, \mathbf{x})\bigr] = 0 \quad \forall \,\mathbf{x}, \mathbf{x}' \in \mathcal{X}.
% \]
% From the closure property of GPs under linear transformations, one can construct a kernel function that satisfies the boundary constraints \(B\) on \(\mathcal{X}\) as:
% \[
% k(\mathbf{x}, \mathbf{x}')
% = \int_{\mathcal{X} \times \mathcal{X}}
% G(\mathbf{x}, \mathbf{s}) \,\overline{k}(\mathbf{s}, \mathbf{s}') \,G(\mathbf{s}', \mathbf{x}')
% \,d\mathbf{s}\,d\mathbf{s}',
% \]
% where \(\overline{k}\) is any kernel function appropriate for the underlying problem. However, because \(k\) generally has no closed-form expression, applying this approach typically requires direct observations of the GP with covariance kernel \(k\) as training data.  

%\cmtS{can we have the following as a new Section 3? should split up discussions of Dirichlet and Robin boundaries as separate subsections. some discussions providing context on our constructive approach and related literature, e.g., the SPDE stuff from spatial statistics (Lindgren et al.). }

\section{The BdryMat\'ern GP}
\label{sec:Bdrymat_SPDE}
% The methodologies discussed in Section \ref{sec:GP_review} for constructing GPs with boundary conditions each have inherent limitations. For instance, the Markov random field, Projection, and eigenfunction approaches cannot handle problems involving Robin boundary conditions, which involve constraints on both the function values and their derivatives on the boundary and are widely used in applications. Additionally, the Projection and eigenfunction approaches are not physical-informed, as they modify the kernel function directly without explicitly revealing the underlying physical law governing the kernel, such as those represented by PDEs in the Markov random field and Green’s function approaches. We summarize the advantages and disadvantages of each methodology in Table \ref{tab:methods}.


 % \begin{table}
 % \begin{center}
 %    \begin{tabular}{ | c | c | c |c |c |}
 %      \hline
 %       &\thead{ Robin \\ Boundary} & \thead{Computational \\ Efficient} & \thead{Physical \\ Informed} & \thead{Approximation\\ Error} \\
 %      \hline
 %     \thead{Markov\\ Random Field} &  \XSolidBrush & \CheckmarkBold & \CheckmarkBold &  \XSolidBrush  \\
 %     \hline
 %     \thead{Projection\\ Approach} &  \XSolidBrush & \XSolidBrush & \XSolidBrush & \text{Exact Expression}   \\
 %      \hline
 %    \thead{Eigenfunction\\ Approach}&   \XSolidBrush & \CheckmarkBold &  \XSolidBrush & \CheckmarkBold  \\
 %    \hline
 %       \thead{Green's Function \\ Approach}&   \CheckmarkBold &   \XSolidBrush &\CheckmarkBold & \XSolidBrush   \\
 %    \hline
 %    \thead{BdryGP\\ +  Regression Kernel }&   \CheckmarkBold &  \CheckmarkBold &\CheckmarkBold & \CheckmarkBold   \\
 %    \hline
 %    \end{tabular}
 %  \end{center}
 %      \caption{Comparison of methodologies for constructing GPs with boundary constraints.}
 %      \label{tab:methods}
 %    \end{table}

% We note that, apart from the projection approach which is largely application-oriented, the Markov random field, eigenfunction, and Green’s function methods all rely on kernel approximation. In this study, we propose a  framework that unifies these three approach. In specific, our approach extends the SPDE representation from the Markov random field approach, solves the associated Green’s function, and then approximates the convolution of the kernel with this Green’s function using basis functions. As shown in Table \ref{tab:methods}, our approach addresses the limitations present in all previously mentioned methods.

We now present the proposed BdryMat\'ern GP model. We first outline its kernel specification for the Dirichlet boundary setting via a boundary-modified formulation of the SPDE \eqref{eq:SPDE2}. Using an extension of the Feynman-Kac formula for elliptical PDEs, we prove that this BdryMat\'ern kernel admits a path integral form on a connected domain with a boundary set that is twice-differentiable almost everywhere. Such a form permits efficient kernel approximation via finite-element-based algorithms (see Section \ref{sec:kernel_approx}). We then derive analogous BdryMat\'ern kernels with path integral forms for the Neumann and Robin boundary settings. 

\subsection{The BdryMat\'ern Kernel for Dirichlet boundaries}
\label{sec:BdryGP_Green}

Consider first the Dirichlet boundary setting, where the black-box function $f$ is known on $\mathcal{B} = \partial \mathcal{X}$, the boundary set of a potentially irregular domain $\mathcal{X}$. As before, in the homogeneous setting with $f(\partial \mathcal{X}) = 0$, $\mu_{\mathcal{B}}(\mathbf{x})$ can be set as 0. In the inhomogeneous setting, $\mu_{\mathcal{B}}(\mathbf{x})$ can be specified via a careful interpolation of known boundaries, using, e.g., Remark 10 of \cite{dalton2024boundary}. Thus, without loss of generality, we presume in this subsection (unless otherwise stated) the homogeneous setting of $f(\partial \mathcal{X}) = 0$.

To incorporate this boundary information, we investigate the SPDE system in \eqref{eq:SPDE2} with an additional boundary condition on domain $\mathcal{X}$:
\begin{align}
\begin{split}
    (\kappa^2-\bigtriangleup)f(\BFx)&=\CalW_\nu(\BFx), \quad \BFx\in\CalX, 
    \\
    f(\Bx)&=0,\quad \quad \; \; \; \; \; \Bx\in\partial\CalX,\\
\end{split}
\label{eq:bdry_SPDE}
\end{align}
where $\CalW_\nu$ is the Mat\'ern GP with smoothness $\nu>0$ as defined earlier. Constrained to a general irregular domain $\mathcal{X}$ and not $\Real^d$, the SPDE \eqref{eq:bdry_SPDE} may not have a closed-form solution. However, since $(\kappa^2-\bigtriangleup)$ is a linear operator, a unique solution $f$ of \eqref{eq:bdry_SPDE} should also a GP (see \cite[Chapter 4]{lifshits2012lectures}). Our goal is then to represent the covariance function of $f$ in a form that can be efficiently approximated. To do so, we leverage the Green's function representation \citep{folland2009fourier} of the solution $f$; further details can be found in \cite{sakellaris2021scale}. Here, the Green's function associated with the operator $(\kappa^2-\bigtriangleup)$ on domain $\mathcal{X}$ is defined as a function $G(\cdot)$ satisfying:
% and our goal is then to write the covariance of $f$ in a form that can be easily calculated and efficiently approximated. To this end, we first need to represent the BdryGP $f$ using the Green's function representation \cite{folland2009fourier}. Specifically, the Green's function associated to the operator $(\kappa^2-\bigtriangleup)$ on domain $\CalX$ is defined as a function $G$ such that
\begin{align}
\begin{split}
    (\kappa^2-\bigtriangleup)\int_\CalX G(\BFx,\Bs)\phi(\Bs) d\Bs&=\int_\CalX G(\BFx,\Bs)(\kappa^2-\bigtriangleup)\phi(\Bs) d\Bs=\phi(\Bs),\\
    G(\BFx,\Bx')&=G(\BFx',\Bx),\\
    G(\Bx,\cdot)&=0, \quad \quad \Bx\in \partial\CalX.
\end{split}
\label{eq:greendef}
\end{align}
for any integrable function $\phi : \mathcal{X} \rightarrow \Real$.

While the Green's function $G$ has no closed-form solution here, it does provide a useful integral representation of $f$ as:
\begin{equation}
\label{eq:Green_SPDE}
    f(\BFx)=\int_{\CalX}G(\BFx,\Bs)\CalW_\nu(\Bs)d\Bs.
\end{equation}
Using such a representation, we can write the covariance function of $f$ as:
\begin{align}
    k_{\nu+2,\mathcal{B}}(\BFx,\BFx') 
    &= \E\left[\int_{\CalX} G(\BFx, \Bs) \CalW_\nu(\Bs) \, d\Bs \int_{\CalX} G(\BFx', \bold{s}') \CalW_\nu(\bold{s}') \, d\bold{s}' \right] \nonumber \\
    &= \int_{\CalX \times \CalX} G(\BFx, \Bs) k_\nu(\Bs, \bold{s}') G(\bold{s}', \BFx') \, d\Bs \, d\bold{u}.
    \label{eq:kernel_green}
\end{align}
Here, $k_\nu(\Bs, \bold{s}')$ is the Mat\'ern kernel from \eqref{eq:matern}, and the second line follows from applying Fubini's theorem to exchange the integral and expectation. The notation $k_{\nu+2,\mathcal{B}}(\BFx,\BFx')$ will be explained later. Combining \eqref{eq:greendef} and \eqref{eq:kernel_green}, we see that if we let the linear operator \((\kappa^2 - \Delta)\) act on one variable of \(k_{\nu+2,\mathcal{B}}(\BFx, \BFx')\) and then act on the other variable, the resulting function reduces to the Mat\'ern kernel \(k_\nu(\BFx, \BFx')\).

% We call the kernel function $k^{\rm BM}_{\nu+2}(\BFx,\BFx')$ BdryMat\'ern kernel of order $\nu+2$.

% From \eqref{eq:Green_SPDE}, we can observe that if we let the linear operator \((\kappa^2 - \Delta)\) act on one variable of \(k^{\rm BM}_{\nu+2}(\BFx, \BFx')\) and then act on the other variable, the resulting function becomes the Mat\'ern kernel \(k_\nu^{\rm mat}(\BFx, \BFx')\) according to the definition of Green's function.

Using \eqref{eq:kernel_green}, we can then leverage the Feynman-Kac formula for Dirichlet boundaries \citep{ito1957fundamental,stroock1971diffusion,papanicolaou1990probabilistic} to derive a path integral form for the new kernel $k_{\nu+2,\mathcal{B}}(\BFx, \BFx')$. This is given by the following theorem:
\begin{theorem}[BdryMat\'ern kernel for Dirichlet boundary]
\label{thm:bdryMatern_path_integral_dirichlet}
    Suppose the domain $\CalX$ is connected, with its boundary twice-differentiable almost everywhere. Then the kernel $k_{\nu+2,\mathcal{B}}(\BFx,\BFx')$ in \eqref{eq:kernel_green} takes the path integral form:
    % we have the following path integral representation of the BdryMat\'ern kernel when the GP follows Dirichlet boundary condition
    \begin{align}
    \begin{split}
        k_{\nu+2,\mathcal{B}}(\BFx,\BFx')
        =  &k_{\nu+2}(\BFx,\BFx')-\E\left[k_{\nu+2}(\BFx,\bold{B}'_{\tau})e^{-\kappa^2\tau}\right]\\
        &-\E\left[k_{\nu+2}(\bold{B}_{\gamma},\BFx')e^{-\kappa^2\gamma}\right]+\E\left[k_{\nu+2}(\bold{B}_{\gamma},\bold{B}'_{\tau})e^{-\kappa^2(\tau+\gamma)}\right], \quad \mathbf{x},\mathbf{x}' \in \mathcal{X}.
        \end{split}
        \label{eq:bdryMatern_path_integral_dirichlet}
    \end{align}
    Here, $\bold{B}_t$ and $\bold{B}_s'$ are mutually independent $d$-dimensional Brownian motions initialized at $\bold{B}_0=\Bx$ and $\bold{B}'_0=\Bx'$, respectively, and $\tau$ and $\gamma$ are the hitting times to the boundary of the Brownian motions $\bold{B}_t$ and $\bold{B}_s'$, i.e.:
    \begin{align*}
        &\tau=\inf\{t:\bold{B}_t\in\partial\CalX\},\quad \gamma=\inf\{s:\bold{B}'_s\in\partial\CalX\}.
    \end{align*}
\end{theorem}

% Proofs of Theorems \ref{thm:bdryMatern_path_integral_dirichlet} and \ref{thm:bdryMatern_path_integral_robin} can be found in Appendices \ref{appendix:proof_path_integral_dirichlet} and \ref{appendix:proof_path_integral_robin}, respectively. The proofs are extensions of Feynman-Kac Formula for elliptical equations \citep{ito1957fundamental,stroock1971diffusion,papanicolaou1990probabilistic}.

% \subsection{Proof of Theorem \ref{thm:bdryMatern_path_integral_dirichlet}}
% \label{appendix:proof_path_integral_dirichlet}

% \cmtS{can you check notation? I used $\gamma$ and $\gamma'$ to avoid notation conflict with earlier. but pls check.}({ \color{blue} I'd not to change as I have used the same notation in our proof. Changing it will make us hard to check the notation of our proof })

\begin{proof}
Our proof first decomposes \eqref{eq:kernel_green} into four terms corresponding to the terms in \eqref{eq:bdryMatern_path_integral_dirichlet}. For each term, we then rewrite its Green's function representation in path integral integral form. Recall that the Mat\'ern kernel $k_\nu$ admits the spectral representation \cite{Wendland10}:
\begin{equation}
    \label{eq:pf_path_integral_1}
    k_\nu(\BFx,\BFx')=\CalF^{-1}[{(\kappa^2+\|\Bomega\|^2)^{-\alpha}}](\BFx-\BFx')=C_{\CalF^{-1}}\int_{\Real^d}\frac{e^{ \im (\BFx-\BFx')^T\Bomega }}{(\kappa^2+\|\Bomega\|^2)^{\alpha}}d\Bomega,
\end{equation}
where $\mathcal{F}$ is the Fourier transform, and $C_{\CalF^{-1}}$ is the inverse Fourier transform constant that depends only on dimension $d$. Substituting \eqref{eq:pf_path_integral_1} into  \eqref{eq:kernel_green} and applying Fubini's theorem, we get:
\begin{align}
    k_{\nu+2,\mathcal{B}}(\BFx,\BFx')=&C_{\CalF^{-1}}\int_{\Real^d}\left[\int_{\CalX\times\CalX}\frac{G(\BFx,\Bs)e^{ \im (\Bs-\Bu)^T\Bomega}G(\BFx',\Bu)}{(\kappa^2+\|\Bomega\|^2)^{\alpha}}\right]d\Bs d\Bu d\Bomega\nonumber\\
    =& C_{\CalF^{-1}}\int_{\Real^d}(\kappa^2+\|\Bomega\|^2)^{-\alpha}\left[\int_{\CalX}G(\BFx,\Bs)e^{ \im\Bs^T\Bomega} d\Bs\right]\overline{\left[ \int_{\CalX}G(\BFx',\Bu)e^{ \im\Bu^T\Bomega}d\Bu\right]} d\Bomega\nonumber\\
    =& C_{\CalF^{-1}}\int_{\Real^d}\psi(\Bomega) g_{\Bomega}(\BFx)\overline{g_{\Bomega}(\BFx')}d\Bomega\label{eq:pf_path_integral_2},
\end{align}
where $\psi(\Bomega)=(\kappa^2+\|\Bomega\|^2)^{-\alpha}$. Here, following the definition of Green's function, $g_{\Bomega}$ is the solution to the following PDE with zero boundary condition on $\partial\CalX$:
\begin{align*}
    (\kappa^2-\bigtriangleup)g_{\Bomega}=e^{ \im \BFx^T\Bomega},\quad \BFx\in\CalX, \quad g_{\Bomega}(\partial\CalX)=0. 
\end{align*}

% $=h_{\Bomega}-v_{\Bomega}$ with

From this, $g_{\Bomega}$ can then be decomposed into a homogeneous solution $h_{\Bomega}$ and a particular solution $v_{\Bomega} = h_{\Bomega} - g_{\Bomega}$ such that:
\begin{align}
    &(\kappa^2-\bigtriangleup) h_{\Bomega}(\BFx)=e^{ \im \BFx^T\Bomega},\quad  \BFx\in\Real^d, \label{eq:pf_path_integral_3}\\
    &(\kappa^2-\bigtriangleup) v_{\Bomega}(\BFx)=0,\quad \quad \quad \BFx\in\CalX, \quad v_{\Bomega}(\partial\CalX)=h_\omega(\partial\CalX).\label{eq:pf_path_integral_4}
\end{align}
Substituting the decomposition $g_{\Bomega}=h_{\Bomega}-v_{\Bomega}$ into \eqref{eq:pf_path_integral_2}, the kernel $k_{\nu+2,\mathcal{B}}(\BFx,\BFx')$ can then be decomposed into four different terms:
\begin{equation}
    k_{\nu+2,\mathcal{B}}(\BFx,\BFx')=K_{hh}(\BFx,\BFx')-K_{hv}(\BFx,\BFx')-K_{hv}(\BFx',\BFx)+K_{vv}(\BFx,\BFx'),\label{eq:pf_path_integral_5}
\end{equation}
where:
\begin{align}
\begin{split}
    &K_{hh}(\BFx,\BFx')=C_{\CalF^{-1}}\int_{\Real^d}\psi(\Bomega)h_{\Bomega}(\BFx)\overline{h_{\Bomega}(\BFx')}d\Bomega,\\
    &K_{hv}(\BFx,\BFx')=C_{\CalF^{-1}}\int_{\Real^d}\psi(\Bomega)h_{\Bomega}(\BFx)\overline{v_{\Bomega}(\BFx')}d\Bomega,\\
    &K_{vv}(\BFx,\BFx')=C_{\CalF^{-1}}\int_{\Real^d}\psi(\Bomega)v_{\Bomega}(\BFx)\overline{v_{\Bomega}(\BFx')}d\Bomega.
    \label{eq:khh_int}
\end{split}
\end{align}

Consider the first term $K_{hh}(\BFx,\BFx')$. From \eqref{eq:pf_path_integral_3}, we see that $h_{\Bomega}$ is the inverse Fourier transform of the Green's function of the operator $(\kappa^2-\bigtriangleup)$ over the whole domain $\Real^d$. From our direct calculations, we have:
\begin{equation}
    h_{\Bomega}(\BFx)=\frac{e^{\im \BFx^T\Bomega}}{\kappa^2+\|\Bomega\|^2}\label{eq:pf_path_integral_6}.
\end{equation}
Substituting \eqref{eq:pf_path_integral_6} into the integral representation of $K_{hh}$, we can show that $K_{hh}$ corresponds to the Mat\'ern-$(\nu+2)$ kernel, since:
\begin{equation}
    K_{hh}(\BFx,\BFx')=C_{\CalF^{-1}}\int_{\Real^d}\frac{e^{\im(\Bx-\Bx')^T\Bomega}}{(1+\|\Bomega\|^2)^{\alpha+2}}d\Bomega=k_{\nu+2}(\BFx,\BFx'). \label{eq:pf_path_integral_7}
\end{equation}

Consider the next three terms in \eqref{eq:pf_path_integral_5}. Note that $v_{\Bomega}$ has no closed-form solution in general. However, since it is the solution of the elliptical PDE \eqref{eq:pf_path_integral_4} and given the twice-differentiability of the boundary set $\partial \mathcal{X}$ almost everywhere, we can apply the Feynman-Kac formula for Dirichlet boundaries \citep{Chung1981} to obtain the representation:
\begin{equation}
    \label{eq:pf_path_integral_8}
    v_{\Bomega}(\BFx)=\E\left[h_{\Bomega}(\bold{B}_\gamma)e^{-\kappa^2\gamma}|\bold{B}_0=\BFx\right]=\frac{\E\left[e^{\im \bold{B}_\gamma^T\Bomega}e^{-\kappa^2\gamma}|\bold{B}_0=\BFx\right]}{\kappa^2+\|\Bomega\|^2}.
\end{equation} 
% Here, $\bold{B}_t$ is a Brownian motion initialized at $\bold{B}_0=\BFx$, and $\gamma$ is the hitting time of $\bold{B}_t$ to the boundary $\partial\CalX$, i.e., $\gamma=\text{inf}\{t:\bold{B}_t\in\partial\CalX\}$.
% \cmtS{should below be $\bold{B}_\gamma$ instead of $\bold{B}_\tau$? also for next eqn as well. pls check.}({\color{blue} The notation is correct because $K_{hv}$ and $K_{vh}$ are symmetric})
Substituting \eqref{eq:pf_path_integral_8} into the integral representation of $K_{hv}$ and $K_{vv}$ in \eqref{eq:khh_int}, we can apply Fubini's theorem for Gaussian measures to obtain: 
\small
\begin{equation}
\label{eq:pf_path_integral_9}
    K_{hv}(\BFx,\BFx')=\E\left[C_{\CalF^{-1}}\int_{\Real^d}\frac{e^{\im(\Bx-\bold{B}_\tau)^T\Bomega}}{(1+\|\Bomega\|^2)^{\alpha+2}}d\Bomega e^{-\kappa^2\tau}\bigg|\bold{B}_0=\Bx'\right]=\E\left[k_{\nu+2}(\BFx,\bold{B}_\tau)e^{-\kappa^2\tau}|\bold{B}_0=\Bx'\right],
\end{equation}
\normalsize
and:
\begin{align}
\label{eq:pf_path_integral_10}
\begin{split}
    K_{vv}(\BFx,\BFx')=&\E\left[C_{\CalF^{-1}}\int_{\Real^d}\frac{e^{\im(\bold{B}_\tau-\bold{B}'_\gamma)^T\Bomega}}{(1+\|\Bomega\|^2)^{\alpha+2}}d\Bomega e^{-\kappa^2(\tau+\gamma)}\bigg|\bold{B}_0=\Bx, \bold{B}'_0=\Bx'\right]\\
    =&\E\left[k_{\nu+2}(\bold{B}_{\tau},\bold{B}'_\gamma)e^{-\kappa^2(\tau+\gamma)}|\bold{B}_0=\Bx,\bold{B}'_0=\Bx'\right].
\end{split}
\end{align}
The proof is then complete by substituting \eqref{eq:pf_path_integral_7}-\eqref{eq:pf_path_integral_10} into \eqref{eq:pf_path_integral_5}.

\end{proof}

In what follows, we define the kernel $k_{\nu,\mathcal{B}}(\BFx, \BFx')$ in \eqref{eq:bdryMatern_path_integral_dirichlet} as the \textit{BdryMat\'ern kernel} with smoothness parameter $\nu > 2$ for \textit{Dirichlet} boundaries, with process variance $\sigma^2$ and inverse length-scale $\kappa > 0$ inherited from the base Mat\'ern kernel $k_\nu(\BFx, \BFx')$. Here, the restriction of $\nu > 2$ is needed for the Feynman-Kac representation \eqref{eq:pf_path_integral_8}; we thus presume $\nu > 2$ for the BdryMat\'ern kernel for the remainder of the paper. With the mean function $\mu_{\mathcal{B}}(\mathbf{x})$ specified as described at the start of the subsection, the BdryMat\'ern GP for Dirichlet boundaries is then defined as $f \sim \text{GP}\{\mu_{\mathcal{B}}(\mathbf{x}),k_{\nu,\mathcal{B}}(\BFx, \BFx')\}$; such a model embeds the desired boundary information. We will discuss later in Section \ref{sec:kernel_approx} how the posterior predictive distribution of this model can be efficiently approximated.

% Provided $\mu_{\mathcal{B}}$ and $k_{\nu,\mathcal{B}}$ are computed, one can then directly use the GP equations in \eqref{eq:kriging} with $\mu(\Bx) = \mu_{\mathcal{B}}(\Bx)$ and $k(\Bx,\Bx') = k_{\nu,\mathcal{B}}(\BFx, \BFx')$ for boundary-integrated predictive modeling.

A key appeal of the BdryMat\'ern kernel is that it not only ensures the resulting GP sample paths satisfy known Dirichlet boundary conditions, but also provides smoothness control of such paths via the smoothness parameter $\nu$. The latter is analogous to the smoothness control from the Mat\'ern kernel \eqref{eq:matern}. The following theorem proves this for the zero-mean BdryMat\'ern GP given the homogeneous Dirichlet boundary condition $f(\partial \mathcal{X}) = 0$:
\begin{theorem}[Smoothness control for Dirichlet BdryMat\'ern GP]
Suppose the domain $\CalX$ is connected, with its boundary twice-differentiable almost everywhere. Consider the zero-mean BdryMat\'ern GP $f(\cdot) \sim \textup{GP}\{0,k_{\nu,\mathcal{B}}(\BFx, \BFx')\}$, where $k_{\nu,\mathcal{B}}(\BFx, \BFx')$ is the BdryMat\'ern kernel defined in \eqref{eq:bdryMatern_path_integral_dirichlet}. Then:
\begin{enumerate}[label=(\alph*)]
    \item Its sample paths satisfy the zero Dirichlet boundary condition $f(\partial \mathcal{X}) = 0$ almost surely,
    \item Its sample paths are \textup{($\lceil\nu\rceil$-1)}-differentiable almost surely, where $\lceil\cdot\rceil$ denotes the ceiling function.
\end{enumerate}
\label{thm:smoothdir}
\end{theorem}
\noindent The proof of this theorem is provided in Appendix \ref{pf:smoothdir}. This theorem holds analogously for inhomogeneous Dirichlet boundary information on $\partial \mathcal{X}$ (i.e., $f(\partial \mathcal{X}) \neq 0$), as long as the specified mean function $\mu_{\mathcal{B}}(\Bx)$ (see Remark 10 of \cite{dalton2024boundary}) satisfies the same smoothness properties from part (b).

\subsection{The BdryMat\'ern Kernel for Neumann and Robin Boundaries}
Consider next the Neumann and Robin boundary setting. Homogeneous Robin boundaries typically presume (see \cite{zwillinger2021handbook}) known boundaries of the form:
\begin{equation}
\frac{\partial f(\Bx)}{\partial \bold{n}} + c(\Bx)f(\Bx) = 0, \quad \quad \Bx \in \partial \mathcal{X},
\label{eq:robinbdry}
\end{equation}
where $c(\mathbf{x}) \geq 0$ is given and $\bold{n}$ is the inward unit
normal vector on $\partial\CalX$. Again, without loss of generality, we presume in this subsection the homogeneous setting; the inhomogeneous setting can be accommodated via a careful specification of the mean function $\mu_{\mathcal{B}}(\mathbf{x})$ via interpolation; see Remark 10 of \cite{dalton2024boundary}. Note that Neumann boundaries correspond to a specific case of Robin boundaries with $c(\mathbf{x})=0$; the following development on Robin boundaries thus holds for Neumann boundaries as well. 

% \cmtS{discussion on mean function specification? this might be more involved than in Dirichlet case so worth clarification}

To incorporate the above Robin boundary information, we investigate the SPDE system in \eqref{eq:SPDE2} with an additional boundary condition on domain $\mathcal{X}$:
\begin{align}
\label{eq:bdry_SPDE3}
\begin{split}
    (\kappa^2-\bigtriangleup)f(\BFx) & =\CalW_\nu(\BFx), \quad \BFx\in\CalX, 
    \\
    \frac{\partial f(\Bx)}{\partial \bold{n}}+c(\mathbf{x})f(\Bx) & =0,\quad \quad \quad  \Bx\in\partial \CalX,
\end{split}
\end{align}
where $\mathcal{W}_\nu$ is the Mat\'ern GP with smoothness parameter $\nu > 0$. Again, the SPDE \eqref{eq:bdry_SPDE3} may not have a closed-form solution on an irregular domain $\mathcal{X}$. We thus apply a similar Green's function approach as for the earlier Dirichlet boundary setting. Here, the Green's function $G(\cdot$) for operator $(\kappa^2 - \bigtriangleup)$ on $\mathcal{X}$ satisfies:
% $c\geq 0$ is some non-negative function, and $\bold{n}$ is the inward unit
% normal vector on $\partial\CalX$. We can notice that the Neumann boundary condition corresponds to the Robin case with $c=0$ so it can be seen as a special case of \eqref{eq:bdry_SPDE}.
\begin{align}
\begin{split}
    (\kappa^2-\bigtriangleup)\int_\CalX G(\BFx,\Bs)\phi(\Bs) d\Bs & =\int_\CalX G(\BFx,\Bs)(\kappa^2-\bigtriangleup)\phi(\Bs) d\Bs=\phi(\Bs),\\
    G(\BFx,\Bx')&=G(\BFx',\Bx),\\
    \partial_{\bold{n}}G(\Bx,\cdot)+c(\mathbf{x})G(\Bx,\cdot)&=0, \quad \Bx\in \partial\CalX,
\end{split}
\label{eq:greendef2}
\end{align}
for any integrable function $\phi: \mathcal{X} \rightarrow \Real$.

Again, while $G$ has no closed-form solution here, it does provide a useful representation of a solution $f$ of \eqref{eq:bdry_SPDE3}, which should be a GP since $(\kappa^2 - \bigtriangleup)$ is a linear operator. In particular, the covariance function of such a solution $f$, again denoted $k_{\nu+2,\mathcal{B}}(\BFx, \BFx')$, can be written in the same form as \eqref{eq:kernel_green}. Using this representation with the Feynman-Kac formula for Robin boundaries \citep{ito1957fundamental,stroock1971diffusion,papanicolaou1990probabilistic}, we can then derive a path integral form for the kernel $k_{\nu+2,\mathcal{B}}(\BFx, \BFx')$. This is given by the following theorem:
\begin{theorem}[BdryMat\'ern kernel for Robin boundary]
\label{thm:bdryMatern_path_integral_robin}
% have the following path integral representation of the BdryMat\'ern kernel when the GP follows Robin boundary condition
    Suppose the domain $\CalX$ is connected, with its boundary twice-differentiable almost everywhere. Then the kernel $k_{\nu+2,\mathcal{B}}(\BFx, \BFx')$ defined above takes the path integral form:
    \begin{align}
    \begin{split}
        k_{\nu+2,\mathcal{B}}(\BFx,\BFx')
        =  &k_{\nu+2}(\BFx,\BFx')-\E\left[\int_0^{\infty} k_{\nu+2}(\BFx,\bold{B}'_{t})e^{-\kappa^2t-\int_0^tc(\bold{B}'_{s})dL_s}dL_t\right]\\
        &-\E\left[\int_0^{\infty} k_{\nu+2}(\bold{B}_{\tau},\BFx')e^{-\kappa^2\tau-\int_0^\tau c(\bold{B}_{s})dL_s}dL_\tau\right]\\
        &+\E\left[\int_0^{\infty}\int_0^{\infty} k_{\nu+2}(\bold{B}_{\tau},\bold{B}'_{t})e^{-\kappa^2(t+\tau)-\int_0^\tau c(\bold{B}_{s})dL_s-\int_0^tc(\bold{B}'_{s})dL_s}dL_tdL_\tau\right].
        \end{split}
        \label{eq:bdryMatern_path_integral_robin}
    \end{align}
    Here, $\bold{B}_t$ and $\bold{B}_s'$ are mutually independent $d$-dimensional Brownian motions initialized at $\bold{B}_0=\Bx$ and $\bold{B}'_0=\Bx'$, respectively, and $L_t$ and $L_\tau$ are the so-called boundary local times defined as:
    \begin{equation}
    L_t=\lim_{\delta\to 0}\frac{\int_0^t\ind_{\{\bold{B}_s\in\CalX_\delta\}}ds}{\delta},\quad L_\tau =\lim_{\delta\to 0}\frac{\int_0^\tau \ind_{\{\bold{B}'_s\in\CalX_\delta\}}ds}{\delta},
        \label{eq:local}
    \end{equation}
    where $\CalX_\delta = \{\BFx\in\CalX: \|\BFx-\partial\CalX\|_2 \leq\delta\}$ is a region of width $\delta$ around the boundary $\partial\CalX$.
\end{theorem}

\begin{proof}

This proof is similar to that for Theorem \ref{thm:bdryMatern_path_integral_dirichlet}. The kernel solution of the form \eqref{eq:kernel_green} can again be decomposed into the four terms:
\begin{equation}
    k_{\nu+2,\mathcal{B}}(\BFx,\BFx')=K_{hh}(\BFx,\BFx')-K_{hv}(\BFx,\BFx')-K_{hv}(\BFx',\BFx)+K_{vv}(\BFx,\BFx'),
    \label{eq:pf_path_integral_robin_2}
\end{equation}
where:
\begin{align}
\begin{split}
    &K_{hh}(\BFx,\BFx')=C_{\CalF^{-1}}\int_{\Real^d}\psi(\Bomega)h_{\Bomega}(\BFx)\overline{h_{\Bomega}(\BFx')}d\Bomega,\\
    &K_{hv}(\BFx,\BFx')=C_{\CalF^{-1}}\int_{\Real^d}\psi(\Bomega)h_{\Bomega}(\BFx)\overline{v_{\Bomega}(\BFx')}d\Bomega,\\
    &K_{vv}(\BFx,\BFx')=C_{\CalF^{-1}}\int_{\Real^d}\psi(\Bomega)v_{\Bomega}(\BFx)\overline{v_{\Bomega}(\BFx')}d\Bomega.
    \label{eq:krep}
\end{split}
\end{align}
Here, $h_{\Bomega}(\BFx)$ is the homogeneous solution of \eqref{eq:pf_path_integral_3}, and $v_{\Bomega}(\BFx)$ is the solution to:
\begin{align}
\begin{split}
    (\kappa^2-\bigtriangleup) v_{\Bomega}(\BFx)&=0,\quad \BFx\in\CalX,\\
    \frac{\partial v_{\Bomega}(\BFx)}{\partial\bold{n}}+c(\BFx)v_{\Bomega}(\BFx)&=0,\quad \BFx\in\partial\CalX. \label{eq:pf_path_integral_robin_1}
\end{split}
\end{align}

As before, the first term in \eqref{eq:pf_path_integral_robin_2} reduces to the Mat\'ern kernel $k_{\nu+2,\mathcal{B}}(\mathbf{x},\mathbf{x}')$. For the remaining terms, note that $v_{\Bomega}$ again has no closed-form solution in general. However, as it is the solution to the elliptical PDE \eqref{eq:pf_path_integral_robin_1} and given the twice-differentiability of $\partial \mathcal{X}$ almost everywhere, one can apply the Feynman-Kac formula for Robin boundaries \cite[Corollary 4.4]{papanicolaou1990probabilistic} to obtain the representation:
\begin{equation}
    v_{\Bomega}(\BFx)=\frac{\E\left[\int_0^\infty \exp\{-\kappa^2t-\int_0^tc(\bold{B}_s)dL_s +i\bold{B}_t^T{\Bomega}\}dL_t|\bold{B}_0=\BFx\right]}{\kappa^2+\|\Bomega\|^2}. \label{eq:pf_path_integral_robin_3}
\end{equation}
Plugging this expression into \eqref{eq:krep}, we can then apply Fubini's theorem (in a similar fashion to \eqref{eq:pf_path_integral_9} and \eqref{eq:pf_path_integral_10}) to exchange the order of integration and expectation, under our assumption that $c(\mathbf{x}) \geq 0$. This yields:
\begin{equation}
    \label{eq:pf_path_integral_robin_4}
    K_{hv}(\BFx,\BFx')=\E\left[\int_0^{\infty} k_{\nu+2}(\bold{B}_{\tau},\BFx)e^{-\kappa^2\tau-\int_0^\tau c(\bold{B}_{s})dL_s}dL_\tau\big|\bold{B}_0=\BFx'\right]
\end{equation}
and:
\begin{align}
\label{eq:pf_path_integral_robin_5}
    K_{vv}(\BFx,\BFx')=\E\left[\int_0^{\infty}\int_0^{\infty} k_{\nu+2}(\bold{B}_{\tau},\bold{B}'_{t})e^{-\kappa^2(t+\tau)-\int_0^\tau c(\bold{B}_{s})dL_s-\int_0^tc(\bold{B}'_{s})dL_s}dL_tdL_\tau\right].
\end{align}
The result is then proven by plugging in \eqref{eq:pf_path_integral_robin_4} and \eqref{eq:pf_path_integral_robin_5} into \eqref{eq:pf_path_integral_robin_2}.
\end{proof}

We define the kernel $ k_{\nu,\mathcal{B}}(\BFx,\BFx')$ in \eqref{eq:bdryMatern_path_integral_robin} as the \textit{BdryMat\'ern kernel} with smoothness $\nu > 2$ for \textit{Robin} boundaries, with process variance $\sigma^2$ and inverse length-scale $\kappa > 0$ inherited from the base Mat\'ern kernel $k_\nu(\BFx, \BFx')$. (In what follows, the notation $k_{\nu,\mathcal{B}}(\BFx,\BFx')$ does not distinguish between Dirichlet and Robin boundaries for brevity; this should be clear from context.) Again, the restriction of $\nu > 2$ is needed for the Feynman-Kac representation \eqref{eq:pf_path_integral_8}. The BdryMat\'ern GP for Robin boundaries then takes the form $f(\cdot) \sim \text{GP}\{\mu_{\mathcal{B}}(\Bx),k_{\nu,\mathcal{B}}(\BFx,\BFx')\}$, where $\mu_{\mathcal{B}}(\Bx)$ is the mean function specified at the start of the subsection. We will discuss in Section \ref{sec:kernel_approx} how to efficiently approximate the posterior predictive distribution of this model.

As before, the above BdryMat\'ern GP not only guarantees its sample paths satisfy known Robin boundary conditions, but also ensures the smoothness of such paths can be controlled by the parameter $\nu$. The following theorem shows this for the zero-mean BdryMat\'ern GP given a zero Robin boundary condition:
\begin{theorem}[Smoothness control for Robin BdryMat\'ern GP]
Suppose the domain $\CalX$ is connected, with its boundary twice-differentiable almost everywhere. Consider the zero-mean BdryMat\'ern GP for Robin boundaries $f(\cdot) \sim \textup{GP}\{0,k_{\nu,\mathcal{B}}(\BFx, \BFx')\}$, where $k_{\nu,\mathcal{B}}(\BFx, \BFx')$ is the BdryMat\'ern kernel defined in \eqref{eq:bdryMatern_path_integral_robin}. Then:
\begin{enumerate}[label=(\alph*)]
    \item Its sample paths satisfy the zero Robin boundary condition $\partial f(\Bx)/\partial \bold{n} + c(\Bx) f(\Bx) = 0$ on $\Bx \in \partial \mathcal{X}$ almost surely,
    \item Its sample paths are \textup{($\lceil\nu\rceil$-1)}-differentiable almost surely.
\end{enumerate}
\label{thm:smooth}
\end{theorem}
\noindent The proof of this theorem is provided in Appendix \ref{pf:smoothrobin}. This theorem again holds analogously for inhomogeneous Robin boundaries on $\partial \mathcal{X}$, as long as the specified mean function $\mu_{\mathcal{B}}(\Bx)$ (see Remark 10 of \cite{dalton2024boundary}) satisfies the required smoothness properties.

For general irregular domains $\mathcal{X}$, the BdryMat\'ern kernel (both for the Dirichlet setting \eqref{eq:bdryMatern_path_integral_dirichlet} and the Robin setting \eqref{eq:bdryMatern_path_integral_robin}) cannot be obtained in closed form and thus requires careful approximation. We present next a finite-element-method approach that performs this kernel approximation to facilitate efficient posterior predictions with reliable error analysis.

% Another key appeal of the BdryMat\'ern GP, both for the Dirichlet setting \eqref{eq:bdryMatern_path_integral_dirichlet} and the Robin setting \eqref{eq:bdryMatern_path_integral_robin}, is that its posterior predictive distribution can be readily approximated via a careful Monte Carlo simulation of the Brownian motion boundary hitting times or local times. Such a Monte Carlo approximation further provides error bounds that can guide approximation error analysis. This approximation procedure, however, needs to be carefully designed for numerical stability and computational efficiency. We will tackle this next in Section \ref{sec:kernel_approx}.

\section{Approximation of the Posterior Predictive Distribution}
\label{sec:kernel_approx}

% For general irregular domains $\mathcal{X}$, recall the BdryMat\'ern kernels in Equations \eqref{eq:bdryMatern_path_integral_dirichlet} and \eqref{eq:bdryMatern_path_integral_robin} cannot be evaluated in closed form and thus needs to be approximated. Such an approximation is then used for estimating the desired posterior predictive distribution $[f(\mathbf{x}_{\rm new})|f(\mathbf{X})]$ (see \eqref{eq:kriging}) at many untested points $\mathbf{x}_{\rm new} \in \mathcal{X}$.

We first consider the approximation of the BdryMat\'ern kernel, which is needed within the GP equations \eqref{eq:kriging} for posterior predictions. This kernel approximation approach requires two steps. First, we construct a coupled Monte Carlo estimator for estimating the Gram kernel matrix for the BdryMat\'ern kernel. Next, we employ a finite-element-method approach that leverages such an estimator for efficient approximation of the posterior predictive distribution with error analysis. In what follows, we again presume the domain $\mathcal{X}$ is connected with its boundary twice-differentiable almost everywhere, as required for defining the BdryMat\'ern GP.

% For surrogate modeling, recall that it is important to have error control over kernel approximations, as such uncertainties need to be propagated downstream \citep{mak2018support} for confident scientific inference. We present next a coupled Monte Carlo approach for estimating the BdryMat\'ern kernel with approximation error bounds, then outline an extended method for scaling up this kernel approximation for larger sample sizes $n$.

\subsection{Kernel Matrix Estimator}\label{sec:kermat}

Consider first the BdryMat\'ern kernel $k_{\nu,\mathcal{B}}$ in \eqref{eq:bdryMatern_path_integral_dirichlet} for the Dirichlet setting. Let $\mathbf{U} = \{\mathbf{u}_1, \cdots, \mathbf{u}_m\}$ denote a set of $m$ points on $\mathcal{X}$, and consider the approximation of its kernel matrix $\mathbf{K}:=\mathbf{K}(\mathbf{U},\mathbf{U}) = [k_{\nu,\mathcal{B}}(\mathbf{u}_i,\mathbf{u}_j)]_{i,j=1}^m$. A naive approach might be to \textit{separately} approximate each matrix entry $K_{ij} = k_{\nu,\mathcal{B}}(\mathbf{u}_i, \mathbf{u}_j)$ via the Monte Carlo (MC) estimator:
\small
   \begin{align}
    \begin{split}
        \tilde{K}_{ij}
        =  &k_{\nu}(\Bu_i,\Bu_j)-k_{\nu}(\Bu_i,\mathbf{W}'_{ij})e^{-\kappa^2\tau_{ij}} - k_{\nu}(\mathbf{W}_{ij},\Bu_j)e^{-\kappa^2\gamma_{ij}} + k_{\nu}(\mathbf{W}_{ij},\mathbf{W}'_{ij})e^{-\kappa^2(\tau_{ij}+\gamma_{ij})},
        \end{split}
        \label{eq:naive_dir}
    \end{align}
    \normalsize
where $(\mathbf{W}_{ij},\tau_{ij})$ and $(\mathbf{W}'_{ij},\gamma_{ij})$ are the simulated hitting locations and hitting times on $\partial \mathcal{X}$ of mutually independent Brownian motions initialized at $\mathbf{u}_i$ and $\mathbf{u}_j$, respectively. The Brownian motions over different indices $(i,j)$ are also mutually independent. Here, only one Monte Carlo draw is used, for reasons we explain later. However, there is a key drawback to such an approach: the estimated kernel matrix $\tilde{\mathbf{K}} = [\tilde{K}_{ij}]_{i,j=1}^m$ may no longer be positive definite, as the eigenvalues of the Monte Carlo error matrix $\tilde{\mathbf{K}}-\mathbf{K}$ can be larger than the minimum eigenvalue of the desired matrix $\mathbf{K}$. A similar issue arises if an analogous naive MC estimator is used for the BdryMat\'ern kernel \eqref{eq:bdryMatern_path_integral_robin} in the Robin setting.
% Second, its computational cost may be high, as it requires the Monte Carlo simulation of $\mathcal{O}(Mn^2)$ hitting times. 

To address this in the Dirichlet setting, we adopt the following \textit{coupled} Monte Carlo approximation of $\mathbf{K}$. For each point $\mathbf{u}_i$, let $(\mathbf{W}_{i},\tau_{i})$ denote the simulated hitting locations and times on $\partial \mathcal{X}$ of a Brownian motion initialized at $\mathbf{u}_i$; the Brownian motions over different indices $i$ are also mutually independent. We can then estimate each matrix entry $K_{ij}$ via the coupled MC estimator:
   \begin{align}
        \hat{K}_{ij}^{\rm C}
        =  &k_{\nu}(\mathbf{u}_i,\mathbf{u}_j)-k_{\nu}(\Bu_i,\mathbf{W}_{j})e^{-\kappa^2\tau_{j}} - k_{\nu}(\mathbf{W}_{i},\Bu_j)e^{-\kappa^2\tau_{i}}+k_{\nu}(\mathbf{W}_{i},\mathbf{W}_{j})e^{-\kappa^2(\tau_{i}+\tau_{j})}.
       \label{eq:coupled_dir}
    \end{align}
For two entries on the same row $i$ (or same column $i$), its estimators from \eqref{eq:coupled_dir} depend on the same hitting time $\tau_{i}$; such estimators are thus ``coupled'' and dependent.

The following theorem shows that this coupled MC estimator $\hat{\mathbf{K}}^{\rm C}(\mathbf{U},\mathbf{U}) = [\hat{K}_{ij}^{\rm C}]_{i,j=1}^m$ for Dirichlet boundary can preserve positive definiteness almost surely:
\begin{theorem}[Positive Definiteness of Dirichlet Coupled MC Estimator]
\label{thm:empirical_BdryMat_Dirichlet}
   Suppose the points \(\BU = \{\Bu_1, \cdots, \Bu_m\} \subset \mathcal{X}\) are distinct and within the interior of $\mathcal{X}$. Further let $\{(\mathbf{W}_{i},\tau_{i})\}_{i=1}$ denote the simulated hitting locations and times on $\partial \mathcal{X}$ of the mutually independent Brownian motions initialized at the points $\{\mathbf{u}_i\}_{i=1}^m$. Then the coupled matrix estimator $\hat{\mathbf{K}}^{\rm C}(\mathbf{U},\mathbf{U})$ with entries $\hat{K}_{ij}^{\rm C}$ defined in \eqref{eq:coupled_dir} is positive definite almost surely. 
\end{theorem}

\begin{proof}
% Note that $\E(\hat{\BK}) =k_{\nu+2,\mathcal{B}}(\BX,\BX)$ as it is an unbiased estimator. We only show that $\BK(\BU,\BU)$ is positive-definite almost surely for any $m$ distinct points $\BU=\{\Bu_i\}_{i=1}$.
The positive definiteness of $\hat{\mathbf{K}}^{\rm C}(\mathbf{U},\mathbf{U})$ can be shown as follows. Define $\bold{W}_i := \Bu_i+\boldsymbol{\delta}_i$, and let $\mathbf{v}=(v_1,\cdots,v_m)^T \in \mathbb{R}^m$ be an arbitrary vector. Using the path integral representation of $K_{hh}$, $K_{hv}$, and $K_{vv}$ in \eqref{eq:pf_path_integral_7}-\eqref{eq:pf_path_integral_10}, the quadratic form $\mathbf{v}^T\hat{\BK}^{\rm C}(\mathbf{U},\mathbf{U})\mathbf{v}$ can be written as:
\small
\begin{align}
    \begin{split}
    \mathbf{v}^T\hat{\BK}^{\rm C}(\mathbf{U},\mathbf{U})\mathbf{v}=&C_{\CalF^{-1}}\int_{\Real^d}\frac{\sum_{i,j}v_iv_je^{\im(\mathbf{u}_i-\mathbf{u}_j)^T\Bomega}\left[1-e^{\im \boldsymbol{\delta}_i^T\Bomega-\kappa^2 \tau_i}-e^{-\im \boldsymbol{\delta}_j^T\Bomega-\kappa^2 \tau_j}+e^{\im(\boldsymbol{\delta}_i-\boldsymbol{\delta}_j)^T\Bomega-\kappa^2 (\tau_i+\tau_j)}\right]}{(\kappa^2+\|\Bomega\|^2)^{\alpha+2}}d\Bomega\\
    =& C_{\CalF^{-1}}\int_{\Real^d}\frac{{\mathbf{v}}^T\left[\bold{I}-\text{diag}[e^{\im \boldsymbol{\delta}_i^T\Bomega-\kappa^2 \tau_i}]\right][e^{\im (\mathbf{u}_i-\mathbf{u}_j)^T\Bomega}]_{i,j}\overline{\left[\bold{I}-\text{diag}[e^{\im \boldsymbol{\delta}_i^T\Bomega-\kappa^2 \tau_i}]\right]{\mathbf{v}}}}{(\kappa^2+\|\Bomega\|^2)^{\alpha+2}}d\Bomega\\
    \geq & C_{\CalF^{-1}}\int_{\Real^d}\frac{{\mathbf{v}}^T\left[\bold{I}-\text{diag}[e^{-\kappa^2\tau_*}]\right][e^{\im (\mathbf{u}_i-\mathbf{u}_j)^T\Bomega}]_{i,j}\overline{\left[\bold{I}-\text{diag}[e^{-\kappa^2\tau_*}]\right]{\mathbf{v}}}}{(\kappa^2+\|\Bomega\|^2)^{\alpha+2}}d\Bomega,\label{eq:pf_empirical_BdryMat_1}
\end{split}
\end{align}
\normalsize
where $\tau_*=\min_{i}\tau_i>0$   and $\lambda_*=1-e^{-\kappa^2\tau_*}>0 $ almost surely for any $\mathbf{U}$ in the interior of $\CalX$. We thus have:
\begin{align}
    \left[\bold{I}-\text{diag}[e^{\im \boldsymbol{\delta}_i-\kappa^2\tau_i}]\right]^2> \left[\bold{I}-\text{diag}[e^{-\kappa^2\tau_*}]\right]^2=\lambda_*^2\bold{I}.
\end{align}
Using this with the last line of \eqref{eq:pf_empirical_BdryMat_1}, it follows that:
\begin{equation}
\mathbf{v}^T\hat{\BK}^{\rm C}(\mathbf{U},\mathbf{U})\mathbf{v} \geq \lambda^2_* C_{\CalF^{-1}}\int_{\Real^d}\frac{\|\sum_{i=1}v_ie^{\im \Bu_i^T\Bomega}\|^2}{(\kappa^2+\|\Bomega\|^2)^{\alpha+2}}d\Bomega>0\quad\quad \text{a.s.,} 
\end{equation}
which completes the proof.

\end{proof}

A similar coupled estimator can be used in the Robin setting for the kernel matrix $\mathbf{K} = [K_{ij}]_{i,j=1}^m = k_{\nu ,\mathcal{B}}(\mathbf{U},\mathbf{U})$. For each point $i$, let $\{(\mathbf{B}_{t}^{(i)}, L_{t}^{(i)})\}_{i=1}^m$ be the mutually independent Brownian motions initialized at $\{\mathbf{u}_i\}_{i=1}^m$ along with its boundary local times as defined in \eqref{eq:local}. We then estimate each matrix entry $K_{ij}^{\rm C}$ via the coupled MC estimator:
\begin{align}
\begin{split}
     \hat{K}_{ij}^{\rm C} &= k_{\nu}(\Bx_i,\Bx_j)-  \int_0^{\infty} k_{\nu}(\Bx_i,\bold{B}^{(j)}_{t})e^{-\kappa^2t-\int_0^tc(\bold{B}^{(j)}_{s})dL_s^{(j)}}dL_t^{(j)}  \\
        & \quad - \int_0^{\infty} k_{\nu}(\bold{B}^{(i)}_{\tau},\Bx_j)e^{-\kappa^2\tau-\int_0^\tau c(\bold{B}^{(i)}_{s})dL_s}dL^{(i)}_\tau  \\
        &\quad +  \int_0^{\infty}\int_0^{\infty} k_{\nu}(\bold{B}^{(i)}_{\tau},\bold{B}^{(j)}_{t})e^{-\kappa^2(t+\tau)-\int_0^\tau c(\bold{B}^{(i)}_{s})dL_s^{(i)}-\int_0^tc(\bold{B}^{(j)}_{s})dL_s^{(j)}}dL_t^{(i)}dL_\tau^{(j)}.
        \label{eq:coupled_robin}
\end{split}
\end{align}
Again, note that for matrix entries on the same row $i$ (or same column $i$), its estimators from \eqref{eq:coupled_robin} depend on the same simulated Brownian motions, and thus are dependent.

The following theorem shows that this coupled MC estimator $\hat{\mathbf{K}}^{\rm C}(\mathbf{U},\mathbf{U}) = [\hat{K}^{\rm C}_{ij}]_{i,j=1}^m$ for Robin boundary again preserves positive definiteness almost surely:
\begin{theorem}[Positive Definiteness of Robin Coupled MC Estimator]
\label{thm:empirical_BdryMat_Robin}
% a standard reflecting Brownian motion (SRBM; see \cmtS{ref})
% , i.e., $\mathbf{B}_{t}^{(i)}=\mathbf{B}_t(\mathbf{u}_i)$ 
   Suppose the points \(\BU = \{\Bu_1, \cdots, \Bu_m\} \subset \mathcal{X}\) are distinct and within the interior of $\mathcal{X}$. Further, let $\{\mathbf{B}_{t}^{(i)}\}_{i=1}^m$ denote the mutually independent standard reflecting Brownian motions (SRBMs; see \cite{harrison1981reflected}) initialized at the points $\{\mathbf{u}_i\}_{i=1}^m$, with $\{L_{t}^{(i)}\}_{i=1}^m$ their corresponding boundary local times as defined in \eqref{eq:local}. If the following inequality is satisfied: 
   \begin{equation}
    \label{eq:empirical_BdryMat_Robin}
    \sup_{\BFx\in\CalX}\E[L_t|\bold{B}_0=\BFx] < (1-e^{-1})\kappa \sqrt{t},
\end{equation}
then the coupled matrix estimator $\hat{\mathbf{K}}^{\rm C}(\mathbf{U},\mathbf{U})$ with entries $\hat{K}_{ij}^{\rm C}$ defined in \eqref{eq:coupled_robin} is positive definite almost surely. 
\end{theorem}

\begin{proof}
Let $\Bu_i+\bold{W}^{(i)}_t=\bold{B}^{(i)}_t$. Similar to the proof for Theorem \ref{thm:empirical_BdryMat_Dirichlet}, note that for any $\bold{v}\in\Real^m$, the quadratic form $\bold{v}^T\BK^{\rm C}(\BU,\BU)\bold{v}$ can be written as:
\begin{align}
\begin{split}
     \bold{v}^T\BK(\BU,\BU)\bold{v}    =&C_{\CalF^{-1}}\int_{\Real^d}\frac{\sum_{i,j}v_iv_je^{\im(\Bu_i-\Bu_j)^T\Bomega}\left[1-S_i-S_j+S_{i}S_j\right]}{(\kappa^2+\|\Bomega\|^2)^{\alpha+2}}d\Bomega\\
    =& C_{\CalF^{-1}}\int_{\Real^d}\frac{{\bold{V}}^T\left[\bold{I}-\text{diag}[S_i]\right][e^{\im (\Bu_i-\Bu_j)^T\Bomega}]_{i,j}\overline{\left[\bold{I}-\text{diag}[S_i]\right]{\bold{V}}}}{(\kappa^2+\|\Bomega\|^2)^{\alpha+2}}d\Bomega,
    \end{split}
\end{align}
where:
\begin{equation*}
    S_i=\int_0^\infty\exp\left\{-\kappa^2t-\int_0^tc(\bold{B}^{(i)}_s)dL_s+\im{\Bomega}^T{\bold{W}^{(i)}_t}\right\}dL_t.
\end{equation*}
Further note that:
\begin{align}
\begin{split}
    \left|S_i\right|=&\left|\int_0^\infty\exp\left\{-\kappa^2t-\int_0^tc(\bold{B}^{(i)}_s)dL_s+\im{\Bomega}^T{\bold{W}^{(i)}_t}\right\}dL_t\right|\\
    \leq & \int_0^\infty\exp\left\{-\kappa^2t-\int_0^tc(\bold{B}^{(i)}_s)dL_s\right\}dL_t\\
    =&\sum_{l=0}^\infty e^{-\kappa^2\vartriangle_t l-\int_0^{\vartriangle_t l}c(\bold{B}^{(i)}_s)dL_s}G_{\vartriangle_t}(\bold{B}^{(i)}_{\vartriangle_t l}),\label{eq:pf_empirical_BdryMat_robin_1}
\end{split}
\end{align}
where the last line of \eqref{eq:pf_empirical_BdryMat_robin_1} follows from the Markov property of the SRBM, which holds for any $\vartriangle_t>0$. Moreover, we have:
\begin{align}
    G_{\vartriangle_t}(\BFx)=&\E\left[\int_0^{\vartriangle_t}\exp\left\{-\kappa^2t-\int_0^tc(\bold{B}_s)dL_s\right\}dL_t|\bold{B}_0=\BFx \right]\leq \sup_{\BFx\in\CalX}\E[L_t|\bold{B}_0=\BFx],\label{eq:pf_empirical_BdryMat_robin_2}
\end{align}
where $\bold{B}_s$ is an SRBM and $L_s$ its boundary local time.

Substituting \eqref{eq:pf_empirical_BdryMat_robin_2} into \eqref{eq:pf_empirical_BdryMat_robin_1}, and using the fact that $c(\BFx)\geq 0$, we have:
\begin{align*}
    |S_i|\leq &\sup_{\BFx\in\CalX}\E[L_t|\bold{B}_0=\BFx]\sum_{l\geq 0}e^{-\kappa^2\vartriangle_t l}
    \leq \frac{\sup_{\BFx\in\CalX}\E[L_t|\bold{B}_0=\BFx]}{1-e^{-\kappa^2 \vartriangle_t}}
    < \frac{(1-e^{-1})\kappa \sqrt{\vartriangle_t} }{1-e^{-\kappa^2 \vartriangle_t}}
    \leq 1,
\end{align*}
where the second-last inequality follows from \eqref{eq:empirical_BdryMat_Robin} and the last inequality is obtained by setting $\vartriangle_t=\kappa^2$. It follows that $\left[\bold{I}-\text{diag}[S_i]\right]^2\geq \lambda_*^2\bold{I}> 0$ and thus $\bold{v}^T\BK^{\rm C}(\BU,\BU)\bold{v} \succ 0$.
\end{proof}

% but they 
 % allows for MC estimation of the corresponding BdryMatérn kernels because simulation provides an efficient way to sample the hitting times \(\tau\) and \(\gamma\), the hitting positions \(\bold{B}_\tau\) and \(\bold{B}_\gamma\) and the boundary local times $L_t$ and $L_\tau$ . Given any pair of input points \(\BFx_i\) and \(\BFx_j\) from the data \(\BFX\), we can generate  samples using the path integral representation \eqref{eq:bdryMatern_path_integral_dirichlet} or \eqref{eq:bdryMatern_path_integral_robin} to obtain an MC estimates \(\hat{\bold{K}}_{i,j}\) for the value \(k^{\rm BM}_{\nu+2}(\BFx_i, \BFx_j)\). However, this approach leads to two major issues. Firstly, suppose we independently construct \(n^2\) MC estimates and use the matrix \(\hat{\bold{K}} = [\hat{\bold{K}}_{i,j}]_{i,j=1}^n\) to estimate the kernel matrix \(k^{\rm BM}_{\nu+2}(\BFX, \BFX)\). The kernel matrix \(\hat{\bold{K}}\) may not be positive definite because the sampling noise \(\bold{\varepsilon} = \hat{\bold{K}} - k^{\rm BM}_{\nu+2}(\BFX, \BFX)\) can be larger than the minimum eigenvalue of \(k^{\rm BM}_{\nu+2}(\BFX, \BFX)\), which is extremely small when dataset is large. Secondly, the computational cost for constructing the estimation is high because we need to run MC simulations for all input and prediction points.



% In this section, we propose a regression method based on efficient sampling of the BdryMatérn kernel \eqref{eq:bdryMatern_path_integral_dirichlet}, which addresses these two major issues. Specifically, the resulting estimator can be computed efficiently, and the positive definite property is preserved. We first introduce  efficient sampling methods for Dirichlet and Robin boundary problems, and then, based on these methods, we explain how to construct a regression estimator for the BdryMatérn kernel.

% \subsection{Sampling of BdryMat\'ern Kernel}

% Instead of sampling the values of the kernel matrix \(k^{\rm BM}_{\nu+2}(\BFX, \BFX)\) entry-wise, we propose a sampling method that generates positive definite matrices directly. Using these generated positive definite matrices, we can construct a finite-rank estimation of the infinite-rank BdryMatérn kernel function. Our sampling method is based on the following theorem:

% \begin{theorem}[Dirichlet]
% \label{thm:empirical_BdryMat_Dirichlet}
%    Let \(\BU = \{\Bu_i \in \CalX\}_{i=1}\) be a set of \(m\) distinct points. Let \((\bold{W}_i, \tau_i)\) be the hitting position and hitting time of a Brownian motion hitting the boundary \(\partial \CalX\) with initial position \(\Bu_i\). Let \({{\BK}}(\BU, \BU)\) be a random matrix with its \((i,j)\) entries defined as:

% \begin{align*}
%     {{\BK}}(\Bu_i, \Bu_j) &= k_{\nu+2}(\Bu_i, \Bu_j) - k_{\nu+2}(\Bu_i, \bold{W}_j) e^{-\kappa^2 \tau_j} \\
%     &\quad - k_{\nu+2}(\bold{W}_i, \Bu_j) e^{-\kappa^2 \tau_i} + k_{\nu+2}(\bold{W}_i, \bold{W}_j) e^{-\kappa^2 (\tau_i + \tau_j)}.
% \end{align*}
% Then the matrix \({{\BK}}(\BU, \BU)\) is positive definite almost surely. 
% \end{theorem}

% \begin{theorem}[Robin]
% \label{thm:empirical_BdryMat_Robin}
%    Let \(\BU = \{\Bu_i \in \CalX\}_{i=1}\) be a set of \(m\) distinct points. Let \((\bold{B}_t^{(i)},L^{(i)}_t)\) be the SRBMs with initial position \(\Bu_i\) and the corresponding boundary local times. Let \({{\BK}}(\BU, \BU)\) be a random matrix with its \((i,j)\) entries defined as:

% \begin{align*}
%      {{\BK}}(\Bu_i, \Bu_j) &= k_{\nu+2}(\Bu_i,\Bu_j)-  \int_0^{\infty} k_{\nu+2}(\Bu_i,\bold{B}^{(j)}_{t})e^{-\kappa^2t-\int_0^tc(\bold{B}^{(j)}_{s})dL_s^{(j)}}dL_t^{(j)}  \\
%         &- \int_0^{\infty} k_{\nu+2}(\bold{B}^{(i)}_{\tau},\Bu_j)e^{-\kappa^2\tau-\int_0^\tau c(\bold{B}^{(i)}_{s})dL_s}dL^{(i)}_\tau  \\
%         &+  \int_0^{\infty}\int_0^{\infty} k_{\nu+2}(\bold{B}^{(i)}_{\tau},\bold{B}^{(j)}_{t})e^{-\kappa^2(t+\tau)-\int_0^\tau c(\bold{B}^{(i)}_{s})dL_s^{(i)}-\int_0^tc(\bold{B}^{(j)}_{s})dL_s^{(j)}}dL_t^{(i)}dL_\tau^{(j)}.
% \end{align*}
% If the following inequality is satisfied
% \begin{equation}
%     \label{eq:empirical_BdryMat_Robin}
%     \sup_{\BFx\in\CalX}\E[L_t|\bold{B}_0=\BFx] < (1-e^{-1})\kappa \sqrt{t},
% \end{equation}
% then the matrix \({{\BK}}(\BU, \BU)\) is positive definite almost surely. 
% \end{theorem}
% \begin{remark}

It should be noted that the condition \eqref{eq:empirical_BdryMat_Robin} in Theorem \ref{thm:empirical_BdryMat_Robin} is quite mild. As shown in \cite{hsu1984reflecting,hsu1985probabilistic}, the boundary local time \( L_t \) satisfies:
\[
\sup_{\BFx \in \CalX} \E[L_t | \bold{B}_0 = \BFx] \leq A\sqrt{t},
\]
where \( A \) is a constant dependent only on the domain \( \CalX \). Thus, our condition \eqref{eq:empirical_BdryMat_Robin} easily satisfied for any \( \kappa > {A}(1 - 1/e) \); see \citep{grebenkov2019probability}. 
% \end{remark}

%Theorem \ref{thm:empirical_BdryMat_Dirichlet} comes from the intuition that the BdryMatérn kernel can be treated as the path-integral expectation over a space of positive-definite, infinite-rank operators with a probability measure induced by the following set of Brownian motions and hitting times:
%\[\mathcal{M}(\CalX) = \left\{(\bold{B}_t(\Bx), \tau(\Bx)) : \bold{B}_0(\Bx) = \Bx, \tau = \inf_t \{\bold{B}_t(\Bx) \in \partial \CalX\}, \Bx \in \CalX\right\}.\]

Finally, it should be noted that the simulation of hitting times and locations (for the Dirichlet setting) and the simulation of boundary local times (for the Robin setting) can be efficiently performed with high accuracy and with theoretical guarantees. Efficient Monte Carlo simulation algorithms for such tasks have been proposed and investigated in the literature; see, e.g., \cite{roger_RSDE,milstein2004stochastic,leimkuhler2023simplest}. In our later experiments, we make use of a MATLAB implementation of the approach in \cite{milstein2004stochastic} for simulating hitting times and local times.

% \cmtS{need here? it's a bit confusing without the context below} It is important to note that our samples retain the positive-definiteness of the kernel matrices—a fundamental property that ensures their validity as covariances of GPs. Because positive-definiteness is maintained, each sampled matrix can be regarded as a noisy realization of the full kernel— not merely as isolated point evaluations at $k_{\nu}(\Bu_i,\Bu_j)$. This perspective lets us employ the sampled matrices directly in a regression framework to recover the latent kernel function, which can be treated as a infinite dimensional matrix.

 % \cmtS{brief sentence on what we used later in implementation?}({\color{blue} done}) 

% generations of  random components $k_{\nu+2}(\Bu_i, \bold{W}_j) e^{-\kappa^2 \tau_j}$ as shown in Theorem  \ref{thm:empirical_BdryMat_Dirichlet} and the one  $\int_0^{\infty} k_{\nu+2}(\Bu_i,\bold{B}^{(j)}_{t})e^{-\kappa^2t-\int_0^tc(\bold{B}^{(j)}_{s})dL_s^{(j)}}dL_t^{(j)}$ as shown in Theorem \ref{thm:empirical_BdryMat_Robin} can be easily implemented with high accuracy.

% \cmtS{can we state a kernel approximation theorem for the above set-up without inducing points? i have concerns on including inducing point approximations: the paper has too much content, and our arguments for using boundary information is to supplement limited sample sizes (so not the setting of large $n$). paper also too long (we still need to include the anisotropic part).}

\subsection{Finite-Element-Method Kernel Approximation} \label{sec:fea}

With the kernel matrix estimation method in Section \ref{sec:kermat}, one may be tempted to directly apply this for estimating the kernel matrix $\mathbf{K}(\mathbf{X},\mathbf{X})$ then compute the posterior predictive distribution $[f(\mathbf{x}_{\rm new})|f(\mathbf{X})]$ from Equation \eqref{eq:kriging}. There are, however, several limitations with this approach. First (i), even with $\mathbf{K}(\mathbf{X},\mathbf{X})$ estimated, one still needs to approximate the kernel $\mathbf{k}(\mathbf{x}_{\rm new},\mathbf{X})$ within the predictive equations \eqref{eq:kriging}. Note that each prediction point $\mathbf{x}_{\rm new}$ requires a separate MC approximation, thus the prediction of multiple new points can be computationally expensive. Second (ii), there is little guarantee that the MC-approximated kernel satisfies the desired boundary conditions in any sense. Using the earlier coupled MC estimator, we present next an FEM-based approach that addresses these limitations for effective posterior approximation with error analysis. In what follows, we combine the discussion for both Dirichlet and Robin boundaries; the distinction should be clear from context.

Finite element methods are broadly used for numerically solving differential equation systems \citep{huebner2001finite}. Let $V^Q=\{\phi_q(\mathbf{x})\}_{q=1}^Q$ be the set of basis functions, with each basis having support set $S_q = {\rm supp}\{\phi_q\}$. FEM leverages the so-called finite elements $\{(\phi_q, S_q)\}_{q=1}^Q$ to approximate the system over a desired domain $\mathcal{X}$. The domain for such support sets, denoted $\hat{\CalX}_Q := \bigcup_{q=1}^Q S_q$, should cover the desired domain $\mathcal{X}$, i.e., $\CalX \subseteq \hat{\CalX}_Q$. There are established strategies for selecting the finite elements $\{(\phi_q, S_q)\}_{q=1}^Q$, e.g., via simplicial methods \citep{dang2012triangulations} and B-spline approaches \citep{hollig2003finite}. Error analysis for the recommended B-spline approach is provided in the next subsection. For this subsection, we adopt the following general form for the finite elements on the Sobolev space  \citep[Chapter 2]{Brenner08}:
\begin{equation}
\ \|\phi_q\|_{\mathscr{H}_{\nu}}<\infty \;\;\; \text{for all} \;\;\; \phi_q\in V^Q, Q\in\NatInt, \quad \overline{\lim_{Q\to\infty}{\rm span}\{V^Q\}} = \mathscr{H}_{\nu}(\CalX).
\label{eq:femreg}
\end{equation}
Here, $\|\cdot\|_{\mathscr{H}_{\nu}}$ is the norm of the reproducing kernel Hilbert space $\mathscr{H}_{\nu}(\CalX)$ (see \cite{Wendland10}) induced by the Mat\'ern kernel with smoothness $\nu$, i.e., $k_{\nu}$. Such a function space can be shown to be equivalent to the $(\nu+d/2)$-order Sobolev space; see \cite[Corollary 1]{Tuo16}. 

% We first select a set of basis functions $\{\phi_q(\mathbf{x})\}_{q=1}^Q$, with each basis having support set $S_q = {\rm supp}\{\phi_q\}$ \cmtS{add characterization of basis functions for FEM}. Here, . The pairs $\{(\phi_q, S_q)\}_{q=1}^Q$ are known as \textit{finite elements}, and are widely used for numerically solving differential equation systems. Figure \ref{fig:triangulation} shows an example of a partition of a disk domain $\mathcal{X}$ by basis functions supported on triangles; this is known as a ``triangulation'' of the domain. The choice of finite elements varies depending on the domain $\CalX$ and the restricted boundary conditions \cite{Brenner08, Iserles09}. Here, we assume that the finite elements $\{(\phi_q, S_q)\}_{q=1}^Q$ span a $Q$-dimensional subspace of the ambient function space in which the underlying function resides.

%\cmtS{this might need a bit of clarification. what do we mean by well-chosen? is this referring to the b-spline approach later?}


% \begin{figure}
% \begin{center}
% \includegraphics[width=0.26\textwidth]
% {figures/triangulation/triangulation.pdf}
% \includegraphics[width=0.3\textwidth]{figures/triangulation/triangulation_Dirichlet.pdf}
% \includegraphics[width=0.3\textwidth]{figures/triangulation/triangulation_Robin.pdf}
% % \includegraphics[width=0.26\textwidth]
% % {figures/triangulation/rectangulation.pdf}
% % \includegraphics[width=0.3\textwidth]{figures/triangulation/rectangulation_Dirichlet.pdf}
% % \includegraphics[width=0.3\textwidth]{figures/triangulation/rectangulation_Robin.pdf}
% \end{center}
% \caption{Left column: discretization of a disk by triangles (upper row) and rectangles(lower row); Middle column: values of a single basis function on elements subject to Dirichlet boundary conditions.
% Right: values of a single basis function on elements with Neumann boundary conditions.
% }
% \label{fig:triangulation}
% \end{figure}

With the choice of finite elements $\{(\phi_q, S_q)\}_{q=1}^Q$ satisfying \eqref{eq:femreg} in hand, consider the following FEM-based approximation of the BdryMat\'ern kernel:
\begin{equation}
\label{eq:kernel_regression}
    \hat{k}_{m,Q}(\BFx, \BFx') = \BPhi(\BFx)^T [\BPhi(\BU)]^\dagger \hat{\BK}^{\rm C}(\BU, \BU) [\BPhi(\BU)^T]^\dagger \BPhi(\BFx'), \quad \mathbf{u}_i \overset{\text{i.i.d.}}{\sim}\text{Unif}(\mathcal{X}), \quad i = 1, \cdots, m,
\end{equation}
where $\hat{\BK}^{\rm C}(\BU, \BU)$ is the earlier coupled MC estimator: \eqref{eq:coupled_dir} for the Dirichlet setting, \eqref{eq:coupled_robin} for the Robin setting. We refer to the point set $\mathbf{U}$ as \textit{inducing points}, following \cite{snelson2005sparse,li2025prospar}. Here, $\BPhi(\BFx) = [\phi_1(\BFx), \cdots, \phi_Q(\BFx)]^T$ is the vector of basis functions at $\mathbf{x}$, $\BPhi(\BU) =[\BPhi(\mathbf{u}_1), \cdots, \BPhi(\mathbf{u}_m)]^T$ consists of basis functions for inducing points $\mathbf{U}$, and $\mathbf{M}^\dagger$ denotes the pseudo-inverse of matrix $\mathbf{M}$. This kernel approximation offers a solution to the earlier two limitations. For (i), note that $\mathbf{M} = [\BPhi(\BU)]^\dagger \hat{\BK}^{\rm C}(\BU, \BU) [\BPhi(\BU)^T]^\dagger$ requires just a one-shot evaluation, as it does not depend on kernel inputs $\mathbf{x}$ and $\mathbf{x}'$. With $\mathbf{M}$ evaluated, $\hat{k}_{m,Q}(\mathbf{x},\mathbf{x}')$ can be quickly computed for any choice of $\mathbf{x}$ and $\mathbf{x}'$ via matrix multiplication, without need for MC simulation or pseudo-inverses. We discuss later an implementation using B-splines that further reduces this computation. For (ii), the following proposition shows that $\hat{k}_{m,Q}(\mathbf{x},\mathbf{x}')$ satisfies the desired boundary conditions on the approximated domain $\hat{\CalX}_Q$:
\begin{corollary}[Boundary Conditions on FEM-Based Kernel]
\label{coro:posotive_hat_K}
    Let $\hat{\CalX}_Q = \bigcup_{q=1}^Q S_q$ be the approximated domain of \(\CalX\) via the finite elements \(\{(\phi_q, S_q)\}_{q=1}^Q\). Then, for the Dirichlet setting, the approximated kernel $\hat{k}_{m,Q}(\mathbf{x},\mathbf{x}')$ in \eqref{eq:kernel_regression} almost surely satisfies: 
    \begin{equation}
        \hat{k}_{m,Q}(\BFx, \cdot) = 0, \quad \BFx \in \partial \hat{\CalX}.
    \end{equation}
    Similarly, for the Robin setting, the approximated kernel $\hat{k}_{m,Q}(\mathbf{x},\mathbf{x}')$ almost surely satisfies:
    \begin{equation}
    \frac{\partial\hat{k}_{m,Q}(\BFx, \cdot)}{\partial\bold{n}}+c(\BFx)\hat{k}_{m,Q}(\BFx, \cdot) = 0,\quad \BFx \in \partial \hat{\CalX}.
    \end{equation}    
\end{corollary}
\noindent This corollary follows directly from Theorems \ref{thm:empirical_BdryMat_Dirichlet} and \ref{thm:empirical_BdryMat_Robin}.


% Theorems \ref{thm:empirical_BdryMat_Dirichlet} and \ref{thm:empirical_BdryMat_Robin} also represent finite-rank approximations of the infinite-rank BdryMatérn kernel. While we can now work with finite-rank samples, the challenge remains that we cannot generate samples from the infinite-rank operator—only its finite-rank approximation. Nevertheless, inducing-point approximation methods \cite{seeger2003fast, snelson2005sparse} provide finite-rank approximations of infinite-rank kernel functions. Generally speaking, inducing-point methods approximate the infinite-rank kernel matrix \([k(\BFx,\BFx')]_{\BFx,\BFx' \in \CalX}\) by a rank-\(m\) kernel matrix \([k(\Bu_i, \Bu_j)]_{\Bu_i, \Bu_j \in \bold{U}}\), where the \(m\) points \(\bold{U} = \{\Bu_i\}_{i=1}\) are called the inducing points. If we reduce the path-integral representation of the kernel function to the finite-rank matrix space with samples induced by \(m\) Brownian motions, then the expectation of random kernel operators can be strictly defined. Essentially, the expectation is now over a finite-dimensional space with a probability measure induced by  \(m\) Brownian motions. As a result, the kernel matrix ${\BK} (\bold{U},\bold{U})$ in both Theorems \ref{thm:empirical_BdryMat_Dirichlet} and \ref{thm:empirical_BdryMat_Robin} are the realization of a positive-definite kernel matrix induced by a probability measure defined on a space of positive-definite matrices.





%Secondly, the sample is unbiased  only for \(i \neq j\): \(\E{{\BK}}(\Bu_i, \Bu_j)= k^{\rm BM}_{\nu+2}(\Bu_i,\Bu_j)\). When $i=j$, \(\E[k_{\nu+2}(\bold{W}_i, \bold{W}_i) e^{-2\kappa^2 \tau_i}]-\E[k_{\nu+2}(\bold{W}_i, \bold{W}_i') e^{-\kappa^2 (\tau_i + \tau'_i)}] > 0\), meaning the sample \(\BK(\BU, \BU)\) is biased on the diagonal. Nevertheless, this bias on the diagonal, along with the variance of independently distributed noise, can be corrected through bias-variance tradeoff, as we will show later.


%Secondly, the intuition cannot be directly applied. We can neither generate samples according to the probability measure induced by \(\mathcal{M}(\CalX)\) nor compute any element in the space consisting of positive-definite operators, because both the dimension of the set \(\mathcal{M}(\CalX)\) and the rank of the operators are infinite.


% However, the inducing point method is still computationally expensive. Classical inducing point methods require \(\CalO(n m^2)\) time complexity, where \(n\) and \(m\) are the numbers of data points and inducing points, respectively. Note that the kernel matrix \({\BK}(\bold{U}, \bold{U})\) is only a realization of the true inducing-point approximation. We must run MC estimation to produce the inducing approximation of posterior distribution at each predictive point \(\BFx\). Therefore, direct inducing-point approximation combined with the MC method still lacks efficiency.


% \subsection{Regression Kernel Approximation}

%\cmtS{can we provably show that our approximation approaches below is computationally preferable to simply discretizing the boundary with sample points then conditioning on them? this would help motivate the following developments on approximation} 

% We propose an estimate of the BdryMatérn kernel based on regression for noisy data.

% Given a set of inducing points \(\BU\) and a single sample \(\BK(\BU, \BU)\) as in Theorems \ref{thm:empirical_BdryMat_Dirichlet} and \ref{thm:empirical_BdryMat_Robin}, the BdryMatérn kernel is estimated as follows:

% \begin{equation}
%     \label{eq:kernel_regression}
%     \hat{K}(\BFx, \BFx') = \BPhi(\BFx)^T [\BPhi(\BU)]^\dagger \BK(\BU, \BU) [\BPhi(\BU)^T]^\dagger \BPhi(\BFx'),
% \end{equation}
% where the vector \(\BPhi(\BFx) = [\phi_1(\BFx), \cdots, \phi_p(\BFx)]^T\) consists of all \(p\) basis functions, the \((i,j)\)-entry of the matrix \(\BPhi(\BU)\) equals the value of the \(i^{\rm th}\) basis \(\phi_i\) at the \(j^{\rm th}\) inducing point \(\Bu_j\), and \(\dagger\) denotes the pseudo-inverse.  We call this kernel $\hat{K}$ \emph{regression kernel}. In practice, the number of basis functions \(p\) is much less than the number of observations \(m\) in \eqref{eq:kernel_regression} to balance the variance caused by noise. The finite elements are chosen such that \(S_i \cap S_j = \emptyset\) for most \(i, j = 1, \cdots, m\). Therefore, computing \(\hat{K}(\BFx, \BFx')\) is much more efficient than using MC methods. Additionally, \(\hat{K}(\BFx, \BFx')\) can safely approximate the BdryMatérn kernel because it is semi-positive definite and satisfies the boundary conditions of the approximated domain by the finite elements. Specifically, we have the following corollary of Theorems \ref{thm:empirical_BdryMat_Dirichlet} and \ref{thm:empirical_BdryMat_Robin}:

With this, the posterior predictive distribution $[f(\mathbf{x}_{\rm new})|f(\mathbf{X})]$ can be approximated as follows. First, sample a large number of i.i.d. samples $\mathbf{U} = \{\mathbf{u}_1, \cdots, \mathbf{u}_m\}$ uniformly-at-random on $\mathcal{X}$. Next, given the finite elements \(\{(\phi_q, S_q)\}_{q=1}^Q\), evaluate the intermediate matrix $\mathbf{M} = [\BPhi(\BU)]^\dagger \hat{\BK}^{\rm C}(\BU, \BU) [\BPhi(\BU)^T]^\dagger$. Finally, using the form $\hat{k}_{m,Q}(\mathbf{x},\mathbf{x}') = \BPhi(\BFx)^T \mathbf{M} \BPhi(\BFx')$ in \eqref{eq:kernel_regression}, evaluate the GP predictive mean and variance in \eqref{eq:kriging} using the approximated kernel $k = \hat{k}_{m,Q}$. The following corollary guarantees that, with a sufficiently large number of finite elements $Q$, the kernel matrix $\hat{\mathbf{K}}_{m,Q}(\mathbf{X},\mathbf{X}) = [\hat{k}_{m,Q}(\mathbf{x}_i,\mathbf{x}_j)]_{i,j=1}^n$ required in the GP equations \eqref{eq:kriging} is almost surely positive definite:
\begin{corollary}[Positive Definiteness of FEM-Based Kernel Matrix]
Let $\mathbf{X}$ be a set of $n$ distinct design points on the interior of $\mathcal{X}$, i.e., $\CalX \setminus \partial \CalX$. Choose the number of finite elements $Q$ to be greater than $n$. Then, using the approximated kernel in \eqref{eq:kernel_regression}, its kernel matrix $\hat{\mathbf{K}}_{m,Q}(\mathbf{X},\mathbf{X}) = [\hat{k}_{m,Q}(\mathbf{x}_i,\mathbf{x}_j)]_{i,j=1}^n$ is positive definite almost surely. 
% Also, \(\hat{K}(\BFx, \BFx')\) is semi-positive definite on \(\CalX \setminus \partial \CalX\) almost surely for any set of inducing points \(\BU\subset\CalX\), i.e., given any point set \(\BFX \subset \CalX\) of size \(n\),
%     \[\hat{K}(\BFX, \BFX)\succeq 0,\ \text{with}\  {\rm rank}\left[\hat{K}(\BFX, \BFX)\right] = \min\{p, n\}, \quad \text{a.s.}\]
\label{coro:pd}
\end{corollary}
\noindent This corollary again follows from Theorems \ref{thm:empirical_BdryMat_Dirichlet} and \ref{thm:empirical_BdryMat_Robin}. Algorithm \ref{alg:fem} summarizes the above steps for this FEM-based approximation of the posterior predictive distribution.

\begin{algorithm}[!t]
    \caption{FEM-Based Approximation of the Posterior Predictive Distribution}

\justifying

        \noindent \textbf{Inputs}: Number of inducing points $m$, number of finite elements $Q$, simulated data points $\{(\mathbf{x}_i,f(\mathbf{x}_i)\}_{i=1}^n$, prediction point $\mathbf{x}_{\rm new}$. Ensure $Q > n$.\\
        \vspace{-0.35cm}
        
        \noindent $\bullet$ Generate the finite elements $\{(\phi_q, S_q)\}_{q=1}^Q$ satisfying regularity conditions \eqref{eq:femreg}.\\
        \vspace{-0.2cm}

        \noindent $\bullet$ Sample inducing points $\mathbf{U} = \{\mathbf{u}_1, \cdots, \mathbf{u}_m\} \overset{\text{i.i.d.}}{\sim}\text{Unif}(\mathcal{X})$.\\
        \vspace{-0.35cm}

        \noindent $\bullet$ Compute $\mathbf{M} = [\BPhi(\BU)]^\dagger \hat{\BK}^{\rm C}(\BU, \BU) [\BPhi(\BU)^T]^\dagger$, where $ \hat{\BK}^{\rm C}(\BU, \BU)$ is the coupled MC kernel matrix estimator (Equation \eqref{eq:coupled_dir} for Dirichlet boundary; Equation \eqref{eq:coupled_robin} for Robin boundary).\\
        \vspace{-0.35cm}

        \noindent $\bullet$ Compute the FEM-based approximated kernel $\hat{k}_{m,Q}(\mathbf{x},\mathbf{x}') = \BPhi(\BFx)^T \mathbf{M} \BPhi(\BFx')$, and its approximated kernel matrix $\hat{\mathbf{K}}_{m,Q}(\mathbf{X},\mathbf{X}) = [\hat{{k}}_{m,Q}(\mathbf{x}_i,\mathbf{x}_j)]_{i,j=1}^n$.\\
        \vspace{-0.35cm}

        \noindent $\bullet$ Evaluate the approximated posterior predictive distribution ${\Pi}(\mathbf{x}_{\rm new})$ via Equation \eqref{eq:kriging} with $k = \hat{k}_{m,Q}$.\\
        \vspace{-0.35cm}
        
   \noindent \textbf{Output}: Return ${\Pi}(\mathbf{x}_{\rm new})$.
    % \end{algorithmic}
    \label{alg:fem}
\end{algorithm}

% Such analysis relies on the following decomposition. \cmtS{pls update notation} To construct a precise regression kernel \(\hat{K}\) for the BdryMatérn kernel \(k^{\rm BM}_\nu\), we also need to estimate the error of \(\hat{K}\) constructed from a specific set of inducing points \(\BU\) and finite elements \(\{\phi_i\}\). To this end, we argue that \(\hat{K}\) is a regression estimate of \(k^{\rm BM}_\nu\) with noisy samples on the set \(\BU\). Given the inducing point set \(\BU\), a sample matrix \({\BK}(\BU, \BU)\)  can be treated as noisy observations:
% \begin{equation}
%     {\BK}(\BU, \BU) = k^{\rm BM}_{\nu+2}(\BU, \BU) + \left({\BK}(\BU, \BU) - \E[{\BK}(\BU, \BU)]\right) \coloneqq k^{\rm BM}_{\nu+2}(\BU, \BU) + \boldsymbol{\varepsilon}_{i,j}.\label{eq:noise_decomposition}
% \end{equation}
% The noise decomposition \eqref{eq:noise_decomposition} shows that, conditioned on a specific row \(i\), or more specifically, conditioned on a pair $(\bold{W}_i,\tau_i)$, the corresponding column entries \(j\) of \(\boldsymbol{\varepsilon}_{i,j}\) are independently distributed sampling noise.  The following theorem for the estimation error of the regression kernel \(\hat{K}\) shows that if we choose appropriate inducing points and finite elements, the error always shrinks to \(0\) at a decent convergence rate.

Another appeal of the FEM-based kernel $\hat{k}_{m,Q}$ in \eqref{eq:kernel_regression} is that it facilitates approximation error analysis. The following theorem provides an error bound on how well such a kernel approximates the desired BdryMat\'ern kernel $k_{\nu,\mathcal{B}}$:
\begin{theorem}[FEM-Based Kernel Approximation Error]
    Consider the FEM-based approximated kernel in \eqref{eq:kernel_regression}. Define the \(L^2\)-error between two kernels \(k_1\) and \(k_2\) as:
\begin{equation}
     \|k_1 - k_2\|^2_{L^2(\CalX \times \CalX)} = {\int_{\CalX \times \CalX} \left|k_1(\BFx, \BFx') - k_2(\BFx, \BFx')\right|^2 d\BFx d\BFx'}. \label{eq:HS_norm}
\end{equation}
Then, with probability at least \(1 - v_* - Q(\sqrt{2}e^{-1/2})^{m\underline{\lambda}/S_\phi} - Q\left(\frac{e^{1/2}}{(3/2)^{3/2}}\right)^{m\overline{\lambda}/S_\phi}\), we have:
\begin{equation}
    \label{eq:convergence_general}
    \|\hat{k}_{m,Q} - k_{\nu,\mathcal{B}}\|_{L^2(\CalX \times \CalX)}^2 \leq C \left(\frac{\overline{\lambda}}{\underline{\lambda}}\right)^2 \left[\int_\CalX (\bold{I} - \CalP)[k_\nu + \bar{k}](\bold{I} - \CalP)^T(\Bx, \Bx) d\Bx + \frac{\overline{\lambda}S_\phi}{\underline{\lambda}^2m}\right],
\end{equation}
where $C>0$ is a constant independent of basis choice. Here, $\overline{\lambda}$ and $\underline{\lambda}$ are the maximum and minimum eigenvalues of the Gram matrix $\boldsymbol{\Lambda} = [\lambda_{ij}]_{i,j=1}^Q$, $\lambda_{ij} = \int_\CalX \phi_i(\BFx) \phi_j(\BFx) d\BFx$. Further, \(S_\phi = \max_{\Bx \in \CalX} \sum_{j=1}^p \phi_j^2(\Bx)\), $\bar{k}$ is defined as the kernel:
\small
\begin{equation}
    \bar{k}(\mathbf{x},\mathbf{x}')=
    \begin{cases}
        &\E\left[k_\nu(\BB_\tau,\BB'_\gamma)\bigg|\BB_0=\mathbf{x},\BB'_0=\mathbf{x}'\right]\quad\textup{(Dirichlet)},\\
        &\E\left[\int_{\Real_+^2}e^{-\kappa^2(t+\tau)-\int_0^tc(\BB_s)dL_s-\int_0^\tau c(\BB'_s)dL'_s}k_\nu(\BB_t,\BB'_\tau)dL_tdL'_\tau\bigg|\BB_0=\mathbf{x},\BB'_0=\mathbf{x}'\right]\  \textup{(Robin)},
    \end{cases}
    \label{eq:kbar}
\end{equation}
\normalsize
$\mathcal{P}$ is the $L^2$-projection operator $\CalP f = \arg\min_{h \in \textup{span}\{V^Q\}} \|h - f\|_{L^2(\CalX)}^2$, $\bold{I}$ is the identity operator, and $(\bold{I} - \CalP)k(\cdot, \cdot)(\bold{I} - \CalP)^T(\BFx, \BFx')$ is an operator that applies $\bold{I} - \CalP$ on the first variable of $k(\mathbf{x},\mathbf{x}')$, then $\bold{I} - \CalP$ on its second variable. Finally, $v_* = \int_\CalX \left|(\bold{I} - \CalP)[k_\nu + \bar{k}](\bold{I} - \CalP)^T(\Bx, \Bx)\right|^2 d\Bx$.
    \label{thm:convergence_general}

\end{theorem}
 % Let $\Bu_i\sim\text{Unif}(\CalX)$ be identical distributed for $i=1,\cdots,m$. 
%  For a basis \(\{\phi_j\}_{j=1}^p\), define \(S_\phi = \max_{\Bx \in \CalX} \sum_{j=1}^p \phi_j(\Bx)^2\). Let \(\CalP\) denote the \(L^2\) projection onto \(\text{span}\{\phi_j\}_{j=1}^p\), i.e., \(\CalP f = \arg\min_{h \in \text{span}\{\phi_j\}} \|h - f\|_{L^2(\CalX)}^2\). For the identity operator \(\bold{I}\) and projection \(\CalP\), define
% \[ k(\BFx, \cdot)(\bold{I} - \CalP)^T[(\BFx')] = (\bold{I} - \CalP)[k(\BFx, \cdot)](\BFx') \]
% and 
% \[ \tilde{k}(\BFx, \BFx') = (\bold{I} - \CalP)k(\cdot, \cdot)(\bold{I} - \CalP)^T(\BFx, \BFx'). \]
% This means \(\bold{I} - \CalP\) first acts on the first variable \(\BFx\) and then on the second variable \(\BFx'\) of \(k(\BFx, \BFx')\) to get the projection \(\tilde{k}(\BFx, \BFx')\).
% Define the \(L^2\) error between two kernels \(k_1\) and \(k_2\) as
% \begin{equation}
%      \|k_1 - k_2\|^2_{L^2(\CalX \times \CalX)} = \int_{\CalX \times \CalX} \left|k_1(\BFx, \BFx') - k_2(\BFx, \BFx')\right|^2 d\BFx d\BFx'. \label{eq:HS_norm}
% \end{equation}
% Define the Gram matrix of basis functions \(M_{ij} = \int_\CalX \phi_i(\BFx) \phi_j(\BFx) d\BFx\) and let \(\overline{\lambda} = \lambda_{\max}(M)\) and \(\underline{\lambda} = \lambda_{\min}(M)\) be the maximum and minimum eigenvalues of $M$, respectively. Then for Dirichlet boundary  and Robin boundary satisfying condition \eqref{eq:empirical_BdryMat_Robin}, we have, with probability at least \(1 - v_* - p(\sqrt{2}e^{-1/2})^{m\underline{\lambda}/S_\phi} - p\left(\frac{e^{1/2}}{(3/2)^{3/2}}\right)^{m\overline{\lambda}/S_\phi}\):
% \begin{equation}
%     \label{eq:convergence_general}
%     \|\hat{K} - k^{\rm BM}_\nu\|_{L^2(\CalX \times \CalX)}^2 \leq C \left(\frac{\overline{\lambda}}{\underline{\lambda}}\right)^2 \left[\int_\CalX (\bold{I} - \CalP)[k_\nu + \bar{k}](\bold{I} - \CalP)^T(\Bx, \Bx) d\Bx + \frac{\overline{\lambda}S_\phi}{\underline{\lambda}^2m}\right]
% \end{equation}
% where \(v_* = \int_\CalX \left|(\bold{I} - \CalP)[k + \bar{k}](\bold{I} - \CalP)^T(\Bx, \Bx)\right|^2 d\Bx\), $C$ is some  constant independent of  $\{\phi_j\}$ and $\BU$, and
% \begin{align*}
%     &\bar{k}(\cdot,\cdot)=\\
%     &\begin{cases}
%         &\E\left[k(\BB_\tau,\BB'\gamma)\bigg|\BB_0=\cdot,\BB'_0=\cdot\right]\quad\text{(Dirichlet)}\\
%         &\E\left[\int_{\Real_+^2}e^{-\kappa^2(t+\tau)-\int_0^tc(\BB_s)dL_s-\int_0^\tau c(\BB'_s)dL'_s}k(\BB_t,\BB'_\tau)dL_tdL'_\tau\bigg|\BB_0=\cdot,\BB'_0=\cdot\right]\  \text{(Robin)}.
%     \end{cases}
% \end{align*}
\noindent The proof of this theorem is provided in Appendix \ref{pf:convergence_general}. Note that the norm defined in \eqref{eq:HS_norm} is equivalent to the Hilbert–Schmidt norm if we treat $k_1-k_2$ as an integral operator. Equation \eqref{eq:convergence_general} shows that the approximation error bound depends on the base Mat\'ern kernel $k_{\nu}$ and its projection onto the employed finite element basis. This error bound is derived from the decomposition in Theorem \ref{thm:empirical_BdryMat_Dirichlet}, and we make use of such a bound next to derive a kernel approximation rate using the specific choice of B-spline bases.

\subsection{B-Spline Implementation}
\label{sec:Bspline}
Next, we consider the choice of B-splines \cite{hollig2003finite} for the finite elements $\{(\phi_q,S_q)\}_{q=1}^Q$, which we recommend in implementation. Without loss of generality, we presume that in the following that $\mathcal{X} \subseteq [0,1]^d$; this can easily be done by rescaling the domain. The use of B-splines here has three key advantages. First (i), such bases have been shown to provide effective approximation of both Dirichlet and Robin boundaries in finite-element modeling (see \cite{shu2012differential}). Second (ii), by construction, the support sets of most B-spline basis functions are disjoint, which facilitates an efficient evaluation of the approximated kernel $\hat{k}_{m,Q}$ \eqref{eq:kernel_regression} via sparse matrix operations. Finally (iii), using the famous de Boor's conjecture proven in \cite{shadrin2001norm}, we can nicely characterize the kernel approximation quality of $\hat{k}_{m,Q}$ to the BdryMat\'ern kernel ${k}_{\nu,\mathcal{B}}$ using B-splines.

% . This can be achieved via the Differential Quadrature Method (DQM; \citep{shu2012differential}), which is commonly used for handling complex boundary conditions. DQM approximates such conditions using weighted sums of function values at discrete grid points. In specific, given a irregular domain $\CalX$, we can use discretize $\CalX$ by rectangles as illustrated by the lower row of Figure \ref{fig:triangulation}. Suppose we need to construct a basis function that satisfies the general boundary condition on $A$ boudnary points
% \[c_1(\BFx_i)\frac{\partial f(\BFx_i)}{\partial\bold{n}}+c_2(\BFx_i)f(\BFx_i)=0, \quad i=1,\cdots,A,\ \BFx_i\in\partial\CalX.\]
% We can select $A$ B-splines $\phi_i$ whose supports cover $\{\BFx_i\}_{i=1}^A$ to construct a basis function $\phi_{\bold{a}}(\BFx)=\sum_{i=1}a_i\phi_i$ and solve for the coefficient $\bold{a}=(a_1,\cdots,a_A)$
% \[c_1(\BFx_i)\frac{\partial \phi_{\bold{a}}(\BFx_i)}{\partial\bold{n}}+c_2(\BFx_i)\phi_{\bold{a}}(\BFx_i)=0,\quad i=1,\cdots,A.\]
% There are $A$ unknown $(a_1,\cdots,a_A)$ and $A$ equations in total so the solution is unique. The new B-spline then can exactly satisfy the boundary condition on these discrete points. Implementation and details of this approach for different types of numerical problems can be found in \cite{shu2012differential} \cmtS{did we use this? how so?}({\color{blue} done}).

% If we choose the eigenfunctions of the Green’s function \eqref{eq:Green_SPDE} as basis functions for the regression kernel \eqref{eq:kernel_regression}, the approximation reduces to the eigenfunction approach \cite{solin2019know,solin2020hilbert,gulian2022gaussian} reviewed in Section \ref{sec:GP_review}. However, the eigenfunctions generally have no closed-form on irregular domains. Therefore, we adopt a more flexible class of basis functions, namely B-splines, to construct the regression kernels.

We first briefly review B-splines, following \cite{hollig2003finite}. Let $ N(x) = \boldsymbol{1}_{x \in [0,1]}$, and define:
\begin{equation}
    \label{eq:B_spline}
    N_s(x) = \underbrace{N \ast N \ast \cdots \ast N}_{ \text{$s+1$ times}}(x),
\end{equation}
where \( f \ast g(x) := \int f(x - t)g(t) \, dt \) is the convolution operator. One can show that the support of $N_s$ is $[0,s+1]$. For given multi-indices \( \boldsymbol{\zeta} = (\zeta_1, \ldots, \zeta_d) \in \mathbb{N}^d \) and \( \mathbf{j} = (j_1, \ldots, j_d) \in \mathbb{N}^d \), the multi-dimensional tensor B-spline function indexed by $(\boldsymbol{\zeta}, \mathbf{j})$ is defined as 
$M_{\boldsymbol{\zeta}, \mathbf{j}}^d(x) = \prod_{l=1}^d N_s(2^{\zeta_l}x_l - j_l)$.
Here, \( \boldsymbol{\zeta} \) is a degree vector that controls the resolution of the basis function in each dimension, and \( \mathbf{j} \) specifies the location on which the basis is put. We can similarly define this spline function for a given scalar degree $\zeta \in \mathbb{N}$ as $M_{\zeta, \mathbf{j}}^d(x) = \prod_{l=1}^d N_s(2^\zeta x_l - j_l)$. As the support of $N_s$ is $[0,s+1]$, it is straightforward to show that the support of $M_{\zeta, \mathbf{j}}^d$ is:
\begin{equation}
    \label{eq:B_spline_support}
    {\rm supp}\{M_{\zeta, \mathbf{j}}^d\}=[0,2^{-\zeta}(s+1)]^d+\bold{j}2^{-\zeta}.
\end{equation}
 For a given order $s$, the set of cardinal B-spline basis functions is given by $F_\zeta=\{M_{\zeta, \mathbf{j}}^d, \mathbf{j}\in \mathcal{J}(\zeta)\}$, where $\mathcal{J}(\zeta)=\{-s,-s+1,...,2^\zeta-1,2^\zeta\}^d$. This serves as our set of basis functions $\{\phi_q\}_{q=1}^Q$, with the number of basis functions given as $Q= | F_{\zeta}|$. 


% For a given degree vector \( \mathbf{l} = (l_1, \ldots, l_d) \in \mathbb{N}^d \) and \( \mathbf{j} = (j_1, \ldots, j_d) \in \mathbb{N}^d \), multi-dimensional tensor B-spline labeled by $(\mathbf{l}, \mathbf{j})$ is defined as 
% $M_{\mathbf{l}, \mathbf{j}}^d(x) = \prod_{i=1}^d N_s(2^{l_i}x_i - l_i)$.
% Even for \( l \in \mathbb{N} \), we also use the same notation to express $M_{\mathbf{l}, \mathbf{j}}^d(x) = \prod_{i=1}^d N_s(2^l x_i - j_i)$.
% Here, \( \mathbf{l} \) controls the spatial “resolution” and \( \mathbf{j} \) specifies the location on which the basis is put. From the fact that the support of $N_s$ is $[0,s+1]$, it is straightforward to further show that the support of $M_{\mathbf{l}, \mathbf{j}}^d$ is
% \begin{equation}
%     \label{eq:B_spline_support}
%     {\rm supp}\{M_{\mathbf{l}, \mathbf{j}}^d\}=[0,2^{-l}(s+1)]^d+\bold{j}2^{-l}.
% \end{equation}
%  For  order $s$ cardinal B-spline bases, let $\mathcal{J}(l)=\{-s,-s+1,...,2^l-1,2^l\}^d$ and $F_l=\{M_{\mathbf{l}, \mathbf{j}}^d, \mathbf{j}\in \mathcal{J}(l)\}$. For notation simplicity, we define $\phi_q=M_{\mathbf{l}, \mathbf{j}}^d$ as our basis function and $Q=|F_l|$. 

% The supports of almost all \( M_{\mathbf{l}, \mathbf{j}}^d \in \mathcal{J}(l) \) are disjoint, giving B-splines a sparse structure \cmtS{can we state this }. This sparsity enables the discretization of complex domain and enhances their computational efficiency,  making them ideal basis functions for BdryGP's approximation.

 % Firstly, the compact supports of B-splines enables the  discretization for irregular boundaries and their derivatives. From the compact support property, a finite number of B-splines can be used to construct basis functions with local supports. These basis functions allow for the efficient and accurate local approximation of both Dirichlet and Robin boundary conditions. In this study, the approximation is achieved through the Differential Quadrature Method (DQM), which is commonly used  for handling complex boundary conditions. The implementation of DQM for approximating boundary conditions involves evaluating derivatives of the boundary conditions using weighted sums of function values at discrete grid points. For detailed implementation steps and theoretical background on DQM for boundary condition approximation, please refer to \cite{shu2012differential}.


Consider now point (ii) above. Recall that our kernel approximator takes the form:
\begin{equation}
\label{eq:kernel_reg2}
 \hat{k}_{m,Q}(\BFx, \BFx') = \sum_{q,q'=1}^Q m_{q,q'}{\phi}_q(\BFx) {\phi}_{q'}(\BFx'), \; \mathbf{M} = [m_{q,q'}]_{q,q'=1}^Q = [\BPhi(\BU)]^\dagger \hat{\BK}^{\rm C}(\BU, \BU) [\BPhi(\BU)^T]^\dagger.
\end{equation}
The minimal support overlap between B-spline basis functions (see \cite{hollig2003finite}) can greatly speed up this computation. We see from \eqref{eq:B_spline_support} that (a) many basis functions in $\{\phi_q(\mathbf{x})\}_{q=1}^Q$ are zero for a given $\mathbf{x}$, and (b) the product terms $\{\phi_q(\mathbf{x})\phi_{q'}(\mathbf{x})\}_{q,q'=1}^Q$ are also largely zero for any pair of distinct points $\mathbf{x}$ and $\mathbf{x}'$. Property (a) permits the use of sparse algorithms \citep{ding2020generalization,ding2024sample} for efficient evaluation of the pseudoinverses and matrix products in $\mathbf{M}$ within \eqref{eq:kernel_reg2}. Property (b) permits the efficient evaluation of $\hat{k}_{m,Q}$, as most terms in the sum within \eqref{eq:kernel_reg2} equal zero. With this setup (see Appendix \ref{sec:B_spline_time_complexity} for details), after a one-time pre-processing step with computational cost $\mathcal{O}(Q^3)$, the approximated kernel $\hat{k}_{m,Q}(\BFx,\BFx')$ can be evaluated in $\mathcal{O}\{(2s)^{2d}\}$ work for any pair $(\BFx, \BFx')$. The latter thus provides a constant-time complexity independent of both $Q$ and $n$.








% the supports of most B-splines are disjoint, this enables efficient computation of the regression kernel and subsequent matrix inversions for the posterior. The regression kernel \eqref{eq:kernel_regression} can be expressed as:
% \begin{equation}
%     \hat{K}(\BFx, \BFx') = \sum_{i,j=1}^p \bold{M}_{ij} \phi_i(\BFx) \phi_j(\BFx'),\label{eq:regression:ker_sparse}
% \end{equation}
% where \(\bold{M} = [\BPhi(\BU)]^\dagger \BK(\BU, \BU) [\BPhi(\BU)^T]^\dagger\). 
% Since most B-splines in \eqref{eq:regression:ker_sparse} have disjoint supports, only \(\CalO(1)\) B-splines are non-zero for any input pair \(\BFx\) and \(\BFx'\). This allows the regression kernel computation to be performed in \(\CalO(1)\) time. Consequently, kernel matrix multiplications and inversions involving the regression kernel, \(m\) inducing points, and \(n\) data points can be completed in \(\CalO(\min(m,p)n)\) time. For further details on implementing fast kernel matrix operations for sparse basis kernels, see \cite{ding2020generalization, ding2024sample}.

Consider next point (iii) above. Using B-spline bases, one can use Theorem \ref{thm:convergence_general} to characterize the approximation quality of the kernel $\hat{k}_{m,Q}$ to the desired BdryMat\'ern kernel $k_{\nu,\mathcal{B}}$. This relies on showing that the ratio $\overline{\lambda}/\underline{\lambda} = \mathcal{O}(1)$, which follows from the famous de Boor's conjecture proven in \cite{shadrin2001norm}. This approximation rate using B-splines is presented below:


% Finally, the statistical quantities required in Theorem \ref{thm:convergence_general} to analyze the estimation error of the regression kernel approximation can be efficiently evaluated due to the well-established properties of B-splines. Specifically, it is straightforward to show that \( S_\phi = \CalO(1) \) for B-splines, as only a finite number of non-zero B-splines exist at any given point \(\BFx\). Moreover, the ratio \(\overline{\lambda}/\underline{\lambda} = \CalO(1)\), determined by the Gram matrix induced by B-splines, is independent of the basis functions. This independence follows from the famous de Boor's conjecture, which was proven in \cite{shadrin2001norm}. Based on these results, we can establish the following theorem for B-spline regression kernels:
\begin{theorem}[B-Spline Kernel Approximation Error Rate]
\label{thm:b_spline_regression}
Suppose $\CalX\subseteq[0,1]^d$ and $\nu > d/2$. Let $\{\phi_q\}_{q=1}^Q$ be the set of B-spline basis functions with order $s>\nu+d/2+1/2$ and degree $\zeta = \lfloor (\log_2m)/(2\nu+2d) \rfloor$. Using this basis, the approximated kernel $\hat{k}_{m,Q}$ in \eqref{eq:kernel_regression} satisfies:
\begin{equation}
    \|\hat{k}_{m,Q}-k_{\nu,\mathcal{B}}\|_{L^2(\CalX\times\CalX)}= \CalO(m^{-\frac{\nu}{2\nu+d}}),
    \label{eq:approx_rate}
\end{equation}
     with probability at least $1-C_1m^{-\frac{2\nu-d}{2\nu+d}}-e^{-C_2(m^{\frac{2\nu}{2\nu+d}}-\log m)}$, where $C_1$ and $C_2$ are constants in $m$.
     % \[\|\hat{K}-k^{\rm BM}_{\nu}\|^2_{L^2(\CalX\times\CalX)}\leq \CalO(m^{-\frac{2\nu}{2\nu+d}})\]
     % with probability at least $1-C_1m^{-\frac{2\nu-d}{2\nu+d}}-e^{-C_2(m^{\frac{2\nu}{2\nu+d}}-\log m)}$.
     % Suppose the Mat\'ern kernel has smoothness $\nu>d/2$. For $\CalX\subseteq[0,1]^d$, let $\{\phi_j\}_{j=1}^p=F_l$ be the set of B-spline functions with  smoothness $s-1/2>\nu+d/2$, let $l=(\log_2m)/(2\nu+2d)$.  Then we have the following estimation error for $\hat{K}$:
     % \[\|\hat{K}-k^{\rm BM}_{\nu}\|^2_{L^2(\CalX\times\CalX)}\leq \CalO(m^{-\frac{2\nu}{2\nu+d}})\]
     % with probability at least $1-C_1m^{-\frac{2\nu-d}{2\nu+d}}-e^{-C_2(m^{\frac{2\nu}{2\nu+d}}-\log m)}$.
 \end{theorem}
\noindent The proof of this theorem is provided in Appendix \ref{pf:b_spline_regression}. This approximation rate holds for both the Dirichlet and Robin boundary settings. Note that the rate in \eqref{eq:approx_rate} coincides with the $L^2$-minimax lower bound for GP regression \citep[Theorem 10]{wang2022gaussian}. This is not too surprising, as the task of kernel approximation is related to GP regression in the sense that a (zero-mean) GP is fully characterized by its covariance kernel specification.

An important practical use of such error analysis is that it can guide the specification of $m$, i.e., the number of inducing points, for reliable kernel approximation. Note that, given a desired smoothness parameter $\nu$, only $m$ needs to be specified using the approximation set-up in Theorem \ref{thm:b_spline_regression}; the order and degree of the B-spline bases are provided given this choice of $\nu$ and $m$. One way to specify $m$ is to match the approximation rate in \eqref{eq:approx_rate} with known minimax contraction rates (similarly in $L^2$) for GPs. Using the Mat\'ern-$\nu$ kernel, the latter is known to be of order $\mathcal{O}(n^{-\nu/d})$, where $n$ is the number of evaluation runs on $f$. Setting both rates together and solving for $m$, we get (ignoring constants) that $m = n^{(2\nu+d)/d}$. This specification of $m$ can be used as a rule-of-thumb to ensure the kernel approximation error does not affect the overall GP contraction rate in an asymptotic sense.

% Given $n$ uniformly distributed observations from a Matérn-$\nu$ process, the optimal posterior convergence rate is $\mathcal{O}(n^{-\nu/d})$ in the noiseless case \citep{Tuo17}, and $\mathcal{O}(n^{-\nu/(2\nu + d)})$ in the noisy case \citep{wang2022gaussian}. Therefore, it suffices to choose $m$ large enough such that the approximation error in \eqref{eq:approx_rate} matches the order of the optimal posterior rate. Specifically, we require $m = n^{\frac{2\nu+d}{d}}$ in the noiseless case, and $m = n$ in the noisy case.  This ensures that the approximation error does not affect the overall convergence rate of the posterior.


% For high-dimensional problems, regression kernel approximation may suffer from the  curse of dimensionality, as the convergence rate in Theorem \eqref{thm:b_spline_regression} decays exponentially with increasing \( d \). This rate is likely unimprovable, as it matches the \( L^2 \) minimax convergence rate for classical regression problems.  For high-dimensional boundary problems, a common approach is to impose tensor-type boundary conditions (Section \ref{sec:tensor_mat}). As shown in the next numerical section, this approach significantly enhances modeling accuracy.


%Also, the approximation error  \(\int_\CalX (\bold{I} - \CalP)[k^{\rm BM}_\nu + \bar{k}](\bold{I} - \CalP)^T(\Bx, \Bx) d\Bx\) associated to B-spli



%In particular, the quantity \( S_\phi \), introduced in \cite{cohen2013stability}, can be considered as the 'rank' of the projection \(\mathcal{P}\), while \(\overline{\lambda}\) and \(\underline{\lambda}\) are its maximum and minimum eigenvalues, respectively. These three quantities remain invariant under any rotation applied to the basis functions \(\{\phi_j\}_{j=1}^p\) and are therefore independent of the choice of basis.  %The former set has a Gram matrix with constant eigenvalues, while the latter has a constant \( S_\phi \). Despite these differences, we will demonstrate that both sets yield similar estimation errors.


 %From Lemma \ref{lem:eigen_mat} , it is straightforward to check that, if $\CalP$  is the projection  to the space $\text{span}\{e_i\}_{i=1}^p$, then the calculations of error rate \eqref{eq:convergence_general} becomes calculations of the truncation error. Moreover, the eigenvalues of $M$ are constant because $M=[\langle e_i,e_j\rangle_{L^2}]=\delta_{ij}$ is the identity matrix. So we will have the follow Theorem
% \begin{theorem}
% \label{thm:mat_eigen_regression}
%     Suppose the Mat\'ern kernel has smoothness $\nu>d/2$. Let $\{\phi_i=e_i\}_{i=1}^p$ be the eigenfunctions of Mat\'ern kernel $k_\nu$ on $\CalX$. Let $p=\CalO(m^{\frac{1}{2\nu/d+1}})$. Then we have the following estimation error for $\hat{K}$:
%     \[\|\hat{K}-k^{\rm BM}_{\nu}\|^2_{L^2(\CalX\times\CalX)}\leq \CalO(m^{-\frac{2\nu}{2\nu+d}})\]
     %with probability at least $1-C_1m^{-\frac{4\nu-d}{2\nu+d}}-e^{-C_2(m^{\frac{2\nu}{2\nu+d}}-\log m)}$.
 %\end{theorem}






%\begin{figure}
%\hspace{2cm}
%\begin{center}
%\includegraphics[width=0.32\textwidth]{figures/boundary/bdrymat_ellipse.pdf}
%\includegraphics[width=0.32\textwidth]{figures/boundary/boundary_ellipse.pdf}
%\includegraphics[width=0.32\textwidth]{figures/bdryEffect/B_spline_cubic.pdf}
%\end{center}
%\caption{Left: Plot of \(\hat{k}(\cdot, \bold{0})\), the regression kernel approximation of the BdryMatérn kernel \(k^{\rm BM}_{5/2}\) with zero boundary condition on an ellipse. Riddle: Centers (blue) of B-splines basis functions for regression, the ellipse is approximated by rectangular finite elements. Right: cubic B-splines on the grid points. }
%\label{fig:bdryMat_ellpise}
%\end{figure}
 

 
%The key difference between eigenvalues and B-splines for regression kernel lies in the probability of achieving the convergence rate \( m^{-\frac{2\nu}{2\nu+d}} \). B-spline regression kernel has a lower probability of reaching this convergence rate compared to eigenfunctions. However, the sparse support of B-splines makes them computationally efficient. %Figure \ref{fig:bdryMat_ellpise} illustrates regression kernel with cubic B-splines (smoothness $s=3$).


\section{The Tensor BdryMat\'ern GP Model}
\label{sec:tensor_mat}

%\cmtS{does the earlier boundary GP reduce to this tensor construction in the case of hypercube boundaries? if not, this section might seem a bit out of place; we might not need this with sufficient discussion on the above developments.}

%When the computer simulator has many input variables with different degrees of importance, the BdryMat\'ern GP presented earlier may require modification to account for such varying importances within the surrogate model. For the Mat\'ern GP, a popular way for modeling such anisotropy (see \cite{santner2003}) is via a product form of the one-dimensional Mat\'ern kernels on each variable, where each one-dimensional kernel has distinct length-scale parameters. We present below a similar tensor form modification of the BdryMat\'ern GP that can account for varying variable importance.

Finally, we present a tensor form of the BdryMat\;ern GP. Such a form may be useful when the domain is a high-dimensional unit hypercube, i.e., $\mathcal{X} = [0,1]^d$ for large $d$, and boundary information is known on $\mathcal{X}$. For this high-dimensional setting, the earlier approximation approach may suffer from a ``curse-of-dimensionality'', in that approximation becomes more difficult as $d$ increases. Here, the tensor BdryMat\'ern GP offers a key advantage of a closed-form kernel, which can directly integrate within the GP equations \eqref{eq:kriging} for efficient posterior predictions.

% In high-dimensional settings, approximating the boundary becomes exponentially more challenging. However, accurate specification of boundary conditions is crucial for estimating the underlying functions, as a large proportion of the volume concentrates near the boundary in high-dimensional domains. For high-dimensional problems, we consider domains that can be embedded within a hyper-rectangle, which allows us to model the boundary conditions along its sides using Tensor BdryMatérn GPs.

As before, without loss of generality, we consider the homogeneous boundary setting below; the inhomogeneous setting can be accounted for via a careful specification of its mean function $\mu_{\mathcal{B}}$ (see Section \ref{sec:Bdrymat_SPDE}). Suppose we know the following Robin boundary conditions on $f$:
\begin{align}
\label{eq:boundar_mixed_full}
f(\BFx)+c_l\partial_{x_l}f(\BFx)=0,\quad c_l\geq 0, \quad x_l=0,1, \quad \text{for all } l=1,\cdots,d.
\end{align}
% \begin{align}
% \label{eq:boundar_mixed_1D}
% f(x)+cf'(x)=0,\quad c\geq 0,\quad  x=0,1.
% \end{align}
%  Note that this covers the case of Dirichlet boundaries when $c(\mathbf{x})$ is set as 0.
To incorporate boundary information of this form, we employ the so-called tensor BdryMat\'ern kernel:
\begin{equation}
k_{\nu,\mathcal{B}}^{\rm T}(\mathbf{x},\mathbf{x}') = \prod_{l=1}^d k_{\nu,\mathcal{B},c_l}({x_l},{x_l}').
\label{eq:tensorbdry}
\end{equation}
Here, $k_{\nu,\mathcal{B},c}$ is the one-dimensional BdryMat\'ern kernel from \eqref{eq:bdryMatern_path_integral_robin} with smoothness $\nu$, incorporating homogeneous Robin boundaries of the form:
\begin{align}
\label{eq:boundar_mixed_1D}
   f(x)+cf'(x)=0, \quad c\geq 0, \quad x=0,1,
\end{align}
on the one-dimensional domain $[0,1]$.

% the homogeneous Robin boundaries \eqref{eq:robinbdry} with $c(\mathbf{x}) = c_l$ on the one-dimensional domain $[0,1]$.

% It is also worth noting that the two-point boundary condition on the interval in \eqref{eq:tensorbdry} is generalized to hypersurface boundary conditions on  the boundaries of the hypercube as follows:

% The different inverse length-scales $\kappa_1, \cdots, \kappa_d$ for each input variable can be trained from data (e.g., via maximum likelihood estimation \cite{casella2002}), and provides modeling flexibility for variable importance.

With this, we can derive a closed-form expression for the one-dimensional BdryMat\'ern kernel $k_{\nu,\mathcal{B},c}$ in \eqref{eq:tensorbdry}. This is given in the following theorem:
\begin{theorem}
\label{thm:tensor}
Consider the one-dimensional BdryMat\'ern kernel $k_{\nu,\mathcal{B},c}$ with smoothness $\nu$ for the homogeneous Robin boundaries \eqref{eq:boundar_mixed_1D} on domain $[0,1]$. This kernel has the closed-form representation:
\begin{align}
\begin{split}
k_{\nu,\mathcal{B},c}(x,x') = k_{\nu}(x,x') - k_{hp}(x,x') - k_{hp}(x',x) + k_{hh}(x,x'),
\end{split}
\label{eq:closedform}
\end{align}
\label{thm:closedform}
where $k_{\nu}(\cdot,\cdot)$ is the base Mat\'ern kernel with smoothness $\nu > 2$ and inverse length-scale $\kappa > 0$, and:
\small
\begin{align}
\begin{split}
k_{hp}(x,x') &= \frac{1/2}{\sinh(\kappa)}\left[K_1(x-1)\left(\frac{e^{\kappa x'}}{1+c\kappa}-\frac{e^{-\kappa x'}}{1-c\kappa}\right)+K_1(x)\left(\frac{e^{\kappa(1- x')}}{1-c\kappa}-\frac{e^{-\kappa(1- x')}}{1+c\kappa}\right)\right],\\
k_{hh}(x,x')&=\frac{1}{\sinh({\kappa})^{2}} \times\\
& \quad \left(\left[\cosh(\kappa)K_2(0)-K_2(1)\right]\left(\frac{e^{\kappa(1-x-x')}}{2(1-c\kappa)^2}+\frac{e^{-\kappa(1-x-x')}}{2(1+c\kappa)^2}\right) + \left[\cosh( {\kappa})K_2(1)-K_2(0)\right]\frac{\cosh( {\kappa}(x-x'))}{1-c^2\kappa^2}\right).
\end{split}
\label{eq:closedform2}
\end{align}
\normalsize
Here, $K_1(x)= k_{\nu}(x,x)-c\partial_x k_{\nu}(x,x)$ and $K_2(x)=k_{\nu}(x,x)-c^2\partial_{x x}k_{\nu}(x,x)$.
\end{theorem}

\begin{proof}
% In \eqref{eq:boundar_mixed_1D}, Dirichlet boundary condition corresponds to $c=0$ and general Robin boundary correspons to $c>0$.
% Given the one-dimensional domain $[0,1]$, the considered one-dimensional BdryMat\'ern kernel $k_{\nu,\mathcal{B},c}$ incorporates the Robin boundary condition of the form:
% % we impose the following boundary condition that combine both the Dirichlet and Robin boundary conditions:
% \begin{align}
% \label{eq:boundar_mixed_1D}
%    f(x)+cf'(x)=0,\quad x=0,1,\quad c\geq 0.
% \end{align}
Since the BdryMat\'ern kernel (with $\nu > 2$) can be represented as the convolution between Green's functions and the Mat\'ern kernel, we have:
\begin{align}
    k_{\nu,\mathcal{B}}(x,x')=&\int_{\CalX\times\CalX}G(x,s)k_{\nu-2}(s,u)G(u,x')dsdu
    =C_{\CalF^{-1}}\int_\Real\frac{f(x,\omega)\overline{f(x',\omega)}}{(\kappa^2+\omega^2)^{\nu-3/2}}d\omega.
     \label{eq:kernel_green_1d}
\end{align}
Here, the second equality follows from the spectral density representation of $k_\nu$ and switch of integral, and $f(x,\omega)=\int_{\CalX}G(x,s)e^{i\omega s}ds$ is the solution to the one-dimensional PDE:
\begin{align}
\begin{split}
   & \left(\kappa^2-\frac{\partial^2}{\partial x^2}\right)f(x)=e^{i\omega x},\\
   &f(x)+cf'(x)=0,\quad x=0,1,\quad c\geq 0.
   \end{split}
\label{eq:green_1_character}
\end{align}
The solution to \eqref{eq:green_1_character} can be represented as a particular solution $f_p$ minus a homogeneous $f_h$, i.e., $f(x,\omega)=f_p(x,\omega)-f_h(x,\omega)$, with each part satisfying:
\begin{align*}
    &\left(\kappa^2-\frac{\partial^2}{\partial x^2}
    \right)u_p(x)=e^{i\omega x},\quad \left(\kappa^2-\frac{\partial^2}{\partial x^2}\right)u_h(x)=0,\\
   & u_h(x)+cu_h'(x)=u_p(x),\quad x=0,1.
\end{align*}
Here, both $f_p$ and $f_h$ can be solved analytically as:
\begin{align}
\begin{split}
    & f_p(x)=\frac{e^{i\omega x}}{\kappa^2+\omega^2},\\
    & f_h(x)=\frac{(e^\kappa-e^{i\omega})(1+ic\omega)e^{-\kappa x}}{2(1-c\kappa)\sinh(\kappa)(\kappa^2+\omega^2)}+\frac{(e^{i\omega}-e^{-\kappa})(1+ic\omega)e^{\kappa x}}{2(1+c\kappa)\sinh(\kappa)(\kappa^2+\omega^2)}.
\end{split}
\label{eq:u_p_u_h_solution}
\end{align}
Substitute \eqref{eq:u_p_u_h_solution} into \eqref{eq:kernel_green_1d}, we have the closed form of $k_{\nu,\mathcal{B},c}$ in Equations \eqref{eq:closedform} and \eqref{eq:closedform2}.
% \begin{align}
%     & k^{\rm BM}_{\nu+2}(x,x')=k_{\nu+2}(x,x')-k_{hp}(x,x')-k_{hp}(x',x)+k_{hh}(x,x')  \label{eq:kernel_green_1d_closed} \\
%       & k_{hp}(x,x')=\frac{1/2}{\sinh(\kappa)}\left[K_1(x-1)(\frac{e^{\kappa x'}}{1+c\kappa}-\frac{e^{-\kappa x'}}{1-c\kappa})+K_1(x)(\frac{e^{\kappa(1- x')}}{1-c\kappa}-\frac{e^{-\kappa(1- x')}}{1+c\kappa})\right],\nonumber\\
%     & k_{hh}(x,x')=\bigg(\left[\cosh(\kappa)K_2(0)-K_2(1)\right]\left(\frac{e^{\kappa(1-x-x')}}{2(1-c\kappa)^2}+\frac{e^{-\kappa(1-x-x')}}{2(1+c\kappa)^2}\right)\nonumber\\
%     &\quad+ \left[\cosh( {\kappa})K_2(1)-K_2(0)\right]\frac{\cosh( {\kappa}(x-x'))}{1-c^2\kappa^2}\bigg)\frac{1}{(\sinh({\kappa}))^2},\nonumber\\
%     &K_1(x)=\int_\Real\frac{e^{i\omega x}(1-ic\omega)}{(\kappa^2+\omega^2)^{\nu+5/2}}d\omega= k_{\nu+2}(x)-c\partial_x k_{\nu+2}(x),\nonumber\\
%     &K_2(x)=\int_\Real\frac{e^{i\omega x}(1+c^2\omega^2)}{(\kappa^2+\omega^2)^{\nu+5/2}}d\omega=k_{\nu+2}(x)-c^2\partial_{x x}k_{\nu+2}(x).\nonumber
% \end{align}
\end{proof}

  % The SPDE representation in Section \ref{sec:BdryGP_Green} generally has no closed-form solution. However, \cite{solin2020hilbert} demonstrated that, in one dimension, eigenfunctions of the Green’s function \eqref{eq:Green_SPDE} with Dirichlet boundary conditions can be derived analytically. They approximated the underlying kernel function by a truncated series of eigenfunctions. In this section, we show that BdryGP \eqref{eq:Green_SPDE} and its kernel \eqref{eq:kernel_green} indeed admits closed-form solutions in one dimension for both Dirichlet and Robin boundary conditions. Consequently, as in \cite{solin2020hilbert}, we can construct a tensor BdryGP that handles Dirichlet or Robin boundary conditions on the rectangular domain \([0,1]^d\).

\noindent An analogous closed-form kernel for the Dirichlet boundary setting can be obtained by setting $c=0$. Note that the closed-form kernel from \eqref{eq:closedform} and \eqref{eq:closedform2} aligns with the BdryMat\'ern kernel form in \eqref{eq:bdryMatern_path_integral_dirichlet} and \eqref{eq:bdryMatern_path_integral_robin}; $k_{hp}$ and $k_{hh}$ provide the exact Brownian motion expectations in \eqref{eq:bdryMatern_path_integral_dirichlet} for Dirichlet boundaries and \eqref{eq:bdryMatern_path_integral_robin} for Robin boundaries.

With this, the kernel from \eqref{eq:closedform} and \eqref{eq:closedform2} can then be integrated within \eqref{eq:tensorbdry} to obtain a closed-form form tensor BdryMat\'ern kernel $k_{\nu,\mathcal{B}}^{\rm T}$. One can show that the GP with such a tensor kernel has corresponding sample paths that satisfy the desired boundary information \eqref{eq:boundar_mixed_full}. This closed-form tensor kernel further permits closed-form GP predictive equations from \eqref{eq:kriging} with $k = k_{\nu,\mathcal{B}}^{\rm T}$. Thus, in this simpler setting of a unit hypercube domain $\mathcal{X} = [0,1]^d$ with boundary information of the form \eqref{eq:boundar_mixed_full}, one can bypass the need for kernel approximation via a carefully-specified tensor kernel.

% For the general setting where $\mathcal{X}$ is not a unit hypercube, the kernel approximation approach from Section \ref{sec:kernel_approx} is needed, as the Green's function representation (see \eqref{eq:greendef} and \eqref{eq:greendef2}) may not yield a closed-form solution.

% for larger choices of $\kappa$, the resulting sample paths become increasingly more ``wiggly'', which is expected.

Figure \ref{fig:bdrymat_plot} visualizes different sample paths for the BdryMat\'ern GP using the closed-form kernel from Theorem \ref{thm:tensor}. Here, two choices of smoothness parameters ($\nu = 5/2$ and $7/2$) are tested with inverse length-scale parameter fixed at $\kappa = 5$, under homogeneous Dirichlet, Robin ($c=10$) and Neumann boundary conditions. There are two observations of interest. First, for the Dirichlet (or Neumann) settings, all sample paths appear to go through zero (or have a derivative of zero) at left and right end-points, which satisfies provided boundaries; Robin boundaries are more difficult to see visually. Second, for larger $\nu$, the sample paths appear to have a higher degree of smoothness (in terms of differentiability), which is as desired.

% at the left and right end-points, thus satisfying the provided boundaries. For the Neumann setting, all sample paths appear to have a derivative of zero at end-points, which again satisfies provided boundaries. 

% We observe that under the Dirichlet boundary condition, the sample paths are zero at the endpoints, while under the Neumann condition, their derivatives vanish at the endpoints. The sample paths under the Robin boundary condition with $c = 10$ closely resemble those of the Neumann case.

 
 % illustrates the curves for one dimensional normalized $k^{\rm BM}_{\nu+2}(x,1/2)/k^{\rm BM}_{\nu+2}(1/2,1/2)$ with $\nu=5/2$, and $\nu=7/2$ and general boundary condition \eqref{eq:boundar_mixed_1D} with $c=0$, $c=1$, and $c=10$.


\begin{figure}
\begin{center}
\includegraphics[width=0.3\textwidth]{figures/1dsample_path/BdryMatern-5_2_Dirichlet.pdf}
\includegraphics[width=0.3\textwidth]{figures/1dsample_path/BdryMatern-5_2_Robin.pdf}
\includegraphics[width=0.3\textwidth]
{figures/1dsample_path/BdryMatern-5_2_neumann.pdf}
\includegraphics[width=0.3\textwidth]
{figures/1dsample_path/BdryMatern-7_2_Dirichlet.pdf}
\includegraphics[width=0.3\textwidth]{figures/1dsample_path/BdryMatern-7_2_Robin.pdf}
\includegraphics[width=0.3\textwidth]{figures/1dsample_path/BdryMatern-7_2_neumann.pdf}
\end{center}
\caption{[Top] Sample paths from the BdryMat\'ern GP with smoothness parameter $\nu = 5/2$ given Dirichlet, Robin and Neumann boundary conditions on $[0,1]$. Different colors correspond to different sample paths. [Bottom] Sample paths from the BdryMat\'ern GP with smoothness parameter $\nu = 7/2$ given Dirichlet, Robin and Neumann boundary conditions on $[0,1]$.}
\label{fig:bdrymat_plot}
\end{figure}

% Left Column: BdryGP with Dirichlet boundary condition; Middle Column: BdryGP with Robin boundary condition with $c=10$; Right Column: BdryGP with Neumann boundary condition; Upper Row: BdryMat\'ern-$5/2$ covariance; Lower Row: BdryMat\'ern-$7/2$ covariance. 

% Compared to the BdryMatérn kernel defined through SPDE, the tensor BdryMatérn kernel is easier to implement if the target function has a mixed boundary  condition \eqref{eq:boundar_mixed_1D} on  \([0,1]^d\), which is a sufficient assumption for many problems. On the other hand, due to its SPDE representation which can provide interpretation of uncertainty, the general BdryMat\'ern kernels are more suitable for applications involving physical information, such as geostatistics \cite{LindgrenRueLindstrom11}, physics \cite{ji2024conglomerate}, or complex PDEs emulations \cite{sung2024mesh}.


\section{Numerical Experiments}
\label{sec:numeric}

Finally, we investigate the performance of the proposed boundary-integrated GP model in a suite of numerical experiments. We first explore the BdryMat\'ern GP (Sections \ref{sec:Bdrymat_SPDE} and \ref{sec:kernel_approx}) for a variety of irregular domains in two dimensions, then examine the tensor BdryMat\'ern GP (Section \ref{sec:tensor_mat}) on the higher-dimensional unit hypercube domain $[0,1]^{30}$. 

The general experimental set-up is as follows. In both cases, we compare the proposed model with its corresponding Mat\'ern GP counterpart without boundary conditions (details provided later), to gauge the effect of incorporating boundary information. All models share the same training design points, and model parameters are estimated via maximum likelihood \citep{casella2002}. Using the posterior mean $\hat{f}_n(\cdot)$ as our predictor, predictive performance is evaluated via its log-mean-squared-error (log-MSE) on an independent test set of $10^4$ randomly-selected points on $\mathcal{X}$. All experiments are performed in Matlab R2023a on a MacBook with Apple M2 Max Chip and 32GB of RAM.

% In this section, we will compare GPs with BdryMatérn kernels (BdryGPs) to those with the classical Matérn kernels. 
% In the fisrt part we compare regression kernels for both Dirichlet and Robin boundary conditions on irregular domains. In the second part, we compare tensor BdryGPs on a 30-dimensional domain \([0,1]^{30}\). Our experiments  illustrate that incorporating boundary information can significantly improve the accuracy of high-dimensional predictors.

% Comparisons are based on the logarithmic Mean Squared Error (log-MSE). Let $\hat{f}$ denote a predictor. We first choose a test set $\BX_{\rm test}$ from the input space. For each input $\Bx\in\BX_{\rm test}$, we sample $10^4$ independent samples from the true underlying system and from the predictor to get data sets  $\{{f}_i(\Bx)\}_{i=1}^{10^4}$ and $\{\hat{f}_i(\Bx)\}_{i=1}^{10^4}$, respectively.  The log-MSE is then defined as :
% \[\log-{\rm MSE}= \log\left( \frac{1}{10^4}\sum_{i=1}^{10^4}\left|f(\Bx_i)-\hat{f}(\Bx_i)\right|^2\right)\]
% to quantify the similarity between $f$ and $\hat{f}$. All of the experiments are implemented in Matlab R2023a on a MacBook with Apple M2 Max Chip and 32GB of RAM.

\subsection{BdryMat\'ern GP}
We first explore the BdryMat\'ern GP on a variety of irregular (i.e., non-hypercube) domains. The following four domains are considered:
\begin{itemize}
\item T-shaped: $\mathcal{X}_{\rm T} = ([-1, 1] \times [0, 1]) \cup ([-0.5, 0.5] \times [0, 1])$,
\item Ring-shaped: $\mathcal{X}_{\rm ring} = \left\{ (x_1, x_2) : 1 \leq \sqrt{x_1^2+x_2^2} \leq 2 \right\}$,
\item Disk-shaped: $\mathcal{X}_{\rm disk} = \left\{ (x_1, x_2): \sqrt{x_1^2+x_2^2} \leq 1 \right\}$,
\item Holed-rectangle: $\mathcal{X}_{\rm hrec} = \left\{ (x_1, x_2): |x_1|\leq 1,-1\leq x_2\leq \frac{1}{2}, \sqrt{(x_1-0.5)^2+(x_2+0.5)^2}\geq 0.25\right\}$. 
\end{itemize}
% T-shaped domain \(\mathcal{X}_{\rm T}\), ring-shaped domain \(\mathcal{X}_{\rm ring}\),   disk  \(\mathcal{X}_{\rm disk}\), and  half-space \(\mathcal{X}_{\rm h-rec}\):
% \begin{align*}
%     &\mathcal{X}_{\rm T} = \{ (x_1, x_2) \in ([-1, 1] \times [0, 1]) \cup ([-0.5, 0.5] \times [0, 1]) \}, \\
%     &\mathcal{X}_{\rm ring} = \{ (x_1, x_2) \mid 1 \leq \|(x_1, x_2)\| \leq 2 \}, \\
%     &\mathcal{X}_{\rm disk} = \{ (x_1, x_2): x_1^2+x_2^2<1 \},\quad \mathcal{X}_{\rm h-rec} = \{ (x_1, x_2): x_1<1 \}.
% \end{align*}
For each domain choice, we generate the true function $f$ from the BdryMat\'ern GP with kernel $k_{\nu,\mathcal{B}}$ with smoothness $\nu = 5/2$, inverse length-scale $\kappa = 5 $ and normalized variance parameter $\sigma^2 = \max_{\BFx\in\CalX} \hat{k}_{m,Q}(\BFx,\BFx)$.  For the T-shaped and ring-shaped domains, we adopt the homogeneous (i.e., zero) Dirichlet boundary conditions $f(\partial \mathcal{X}) = 0$, which we presume to be known for surrogate modeling. For the disk and holed-rectangle domains, we adopt Robin boundary conditions of the form \eqref{eq:robinbdry} with $c(\mathbf{x})=20$; these are again presumed to be known for surrogate modeling. 

% In particular, we adopt second-order splines on triangulations with mesh size $0.4$, resulting in $Q = 64, 52, 40, 51$ basis functions for the T-shaped, ring-shaped, disk, and holed-rectangle domains, respectively.

 % At each inducing point, we evaluate the kernel function 10 times and take the average as the observed value for regression.

Next, we sampled $n=50$ training design points $\mathbf{x}_1, \cdots, \mathbf{x}_n$ uniformly-at-random on $\mathcal{X}$, and collected training data $f(\mathbf{x}_1), \cdots, f(\mathbf{x}_n)$. Using such data, we fit the proposed BdryMat\'ern GP (from Section \ref{sec:Bdrymat_SPDE}) with smoothness parameter $\nu = 5/2$, using the FEM-based approximation procedure from Algorithm \ref{alg:fem} for posterior predictions. Here, B-splines are used, with $m = 153, 107, 97, 127$ inducing points for the T-shaped, ring-shaped, disk and holed-rectangle domains, following the nodes of their respective triangulations. We then compare with an (isotropic) Mat\'ern GP with smoothness $\nu = 5/2$, which is the counterpart for the BdryMat\'ern GP from Section \ref{sec:Bdrymat_SPDE} without boundary information. This simulation is replicated 100 times to ensure reproducibility.

% For the T-shape domain $\CalX_{\rm T}$ and ring-shape domain $\CalX_{\rm ring}$ with zero Dirichlet condition, we first use Monte Carlo method to generate a realization of BdryGP with  BdryMat\'ern-$5$ kernel $k^{\rm BM}_5$ in path-integral form \eqref{eq:bdryMatern_path_integral_dirichlet} on the test points. We then randomly collect $n=5,10,\cdots,50$ data $\{(\Bx_i,f(\Bx_i))\}_{i=1}^n$ as our training set and then, we compare the posterior of BdryGP with $k^{\rm BM}_{5/2}$ regression kernels with posteriors of classical Mat\'ern-$5/2$ kernels. For the disk and half-space with Robin boundary condition $\partial_{\bold{n}}f+ 10f=0$, we use the same setup. Each experiment is repeated 100 times with different randomly collected training set. We then plot the log-MSE of our BdryGPs  and the classical Mat\'ern GPs.

Figure \ref{fig:2d_log_err} shows the corresponding mean log-MSEs of the compared models for each domain choice. For all domains, we see that the integration of known boundary information using the BdryMat\'ern GP (with the kernel approximation approach in Section \ref{sec:kernel_approx}) indeed improves predictive performance. Furthermore, as sample size $n$ increases, the BdryMat\'ern GP appears to offer improved prediction error rates in $n$ over its Mat\'ern GP counterpart. This suggests that the integration of boundary information can perhaps accelerate GP prediction rates in the irregular domain setting, much like the improved rates for boundary-integrated surrogates shown in \cite{ding2019bdrygp} for unit hypercube domains. We aim to investigate this rate improvement theoretically as future work. Figure \ref{fig:2d_log_err} visualizes the true response surfaces $f$ for a single experiment replication, and the posterior mean predictor $\hat{f}_n$ from the two compared models. We see that, without boundary information, the Mat\'ern GP may yield erratic predictions particularly when the domain is highly irregular. By incorporating known boundary information on $\mathcal{X}$, the BdryMat\'ern GP provides a considerably improved visual fit of $f$ with limited samples, which corroborates the results in Figure \ref{fig:2d_log_err}.

% It is also useful to explore the varying degrees of improvement provided by boundary integration for different domain choices. From Figure \ref{fig:2d_log_err}, we see that BdryMa\'tern GPs yield noticeably lower prediction errors from boundary integration on all irregular domains over the Mat\'ern GP. A plausible reason is that the covariance functions on irregular domains with boundary conditions are markedly different from those defined for the whole Euclidean space domain. Figure \ref{fig:2d_plot} shows the true function $f$ and the fitted posterior mean predictor $\hat{f}_n$ over different domains in a simulation run. We see that the Mat\'ern GP appears to empirically struggle in balancing the posterior within the domain and outside of the domain.

% The comparisons results are shown in Figure \ref{fig:2d_log_err}. For the  domain $\CalX_{\rm T}$ and $\CalX_{\rm ring}$ with zero Dirichlet conditions and the domain $\CalX_{\rm disk}$ with Robin boundary condition, our experiments reveal that correctly specifying the boundary information allows the predictors to converge to the true underlying function at a fast rate. In contrast, if boundary information is neglected, the convergence can be very slow.   For the  domain $\CalX_{\rm h-rec}$ with Robin boundary condition, BdryGP still slightly outperforms Mat\'ern GP.


% The intuition for the different convergence rates between  BdryGPs and Matérn GPs is illustrated Figure \ref{fig:2d_plot}. In these plots, the random fields with and without boundary information exhibit totally different patterns.  The shapes of domain become are important factor in approximation. For domains highly different than the euclidean space $\Real^2$, such as $\CalX_{\rm T}$, $\CalX_{\rm ring}$, and $\CalX_{\rm disk}$, the Mat\'ern GPs cannot capture the true underlying GPs because it tries to balance the posterior inside and outside of the irregular domain, leading to a wrong posterior. For domain $\CalX_{\rm h-rec}$, because it is similar to $\Real^2$, Mat\'ern GPs still can capture the true underlying except on the boundary.
\begin{figure}[!t]
\hspace{2cm}
\begin{center}
\includegraphics[width=0.23\textwidth]
{figures/logerror/logErr_T.pdf}
\includegraphics[width=0.23\textwidth]{figures/logerror/logErr_ring.pdf}
\includegraphics[width=0.23\textwidth]{figures/logerror/logErr_disk.pdf}
\includegraphics[width=0.23\textwidth]{figures/logerror/logErr_hole_rec.pdf}
\end{center}
\caption{Log-MSEs for the T-shaped, ring-shaped, disk-shaped, and holed-rectangle domain experiments, using the BdryMat\'ern GP and the Mat\'ern GP.}
\label{fig:2d_log_err}
\end{figure}

\begin{figure}
\hspace{2cm}
\begin{center}
\includegraphics[width=0.3\textwidth]
{figures/2Dplot/underlying_T_new.pdf}
\includegraphics[width=0.3\textwidth]{figures/2Dplot/bdryGP_T_new.pdf}
\includegraphics[width=0.3\textwidth]{figures/2Dplot/mat_T_new.pdf}
\includegraphics[width=0.3\textwidth]
{figures/2Dplot/underlying_ring_new.pdf}
\includegraphics[width=0.3\textwidth]{figures/2Dplot/BdryGP_ring_new.pdf}
\includegraphics[width=0.3\textwidth]{figures/2Dplot/mat_ring_new.pdf}
\includegraphics[width=0.3\textwidth]
{figures/2Dplot/underlying_disk}
\includegraphics[width=0.3\textwidth]{figures/2Dplot/bdryGP_disk.pdf}
\includegraphics[width=0.3\textwidth]{figures/2Dplot/mat_disk_new.pdf}
\includegraphics[width=0.3\textwidth]
{figures/2Dplot/underlying_reccir.pdf}
\includegraphics[width=0.3\textwidth]{figures/2Dplot/BdryGP_reccir.pdf}
\includegraphics[width=0.3\textwidth]{figures/2Dplot/mat_reccir.pdf}
\end{center}
\caption{
Visualizing the true underlying function $f$ (left), the BdryMat\'ern GP posterior mean (middle), and the Mat\'ern GP posterior mean functions (right) for the T-shaped, ring-shaped, disk-shaped, and holed-rectangle domain experiments (separated by columns).
% The true underlying (left), the posterior of BdryGP with 50 data points (middle), and the posterior of classical Mat\'ern GP (right) for GPs with Dirichlet boundary on T-shape  (first row) and ring-shape domain (second row), and Robin boundary on disk (third row) and half-space (last row).
}
\label{fig:2d_plot}
\end{figure}

\subsection{Tensor BdryMat\'ern GP}

We next investigate the performance of the tensor BdryMat\'ern GP (Section \ref{sec:tensor_mat}) for incorporating boundary information of the form \eqref{eq:boundar_mixed_full} on a higher-dimensional unit hypercube domain. We consider the two test functions from \cite{Surjano16}:
% case of boundary integration using the tensor BdryMat\'ern GP (Section \ref{sec:tensor_mat}) in a high-dimensional setting, where even small improvements near the boundary can lead to substantial overall improvement, due to the fact that a large proportion of the volume is concentrated near the boundary.
\begin{align}
\begin{split}
    \textit{Griewank:} \quad f(\Bx) &= \sum_{j=1}^{30}\frac{x_j^2}{4000}-\prod_{j=1}^{30}\cos\left(\frac{x_j}{\sqrt{j}}\right)+1 , \quad \Bx\in [-5,5]^{30},\\
   % \textit{Schwefel:} \quad f(\Bx) &= 12000-\sum_{j=1}^{30}x_j\sin(\sqrt{|x_j|}),\ \quad \quad \quad \quad \  \quad \quad \Bx\in[-1,1 ]^{30}\\
    \textit{Product Sine:} \quad f(\Bx) &= \prod_{j=1}^{30}\left[1+\sin\left(\frac{\pi x_j}{2}\right)\right], \quad \Bx\in[-1,1 ]^{30}.
    %\textit{Levy:} \quad f(\Bx) &=\sum_{j=1}^{30}(\frac{x_j-1}{4})^2[1+10\sin^2(\pi\frac{x_j-1}{4}+1)], \quad \Bx\in [-2,2]^{30}
    \end{split}
\end{align}
Both domains are rescaled to the unit hypercube domain  prior to surrogate modeling. Dirichlet boundary information is provided for the Griewank test function; since this function has non-zero Dirichlet boundaries, the mean function $\mu_{\mathcal{B}}$ is specified via the approach in \cite{dalton2024boundary}. Neumann boundary information is provided for the product sine test function; since this function has a Neumann boundary of zero, the mean function is specified as $\mu_{\mathcal{B}} \equiv 0$.

% Neumann boundaries are incorporated for the product sine test function.


% where the Griewank test function is used for Tensor BdryMatérn GPs with Dirichlet boundary conditions on a rectangular domain, and the product sine test function is used for Tensor BdryMatérn GPs with Neumann boundary conditions on a rectangular domain.

% The Griewank function has non-zero boundary conditions; therefore, we ensure that all competing GPs share the same mean function 
% $\mu$, which satisfies the same boundary conditions as the test function. In contrast, the product sine test function satisfies zero Neumann boundary conditions, so we set the mean functions $\mu=0$ for all competing GPs in this case.

With this, we then sample $n$ training design points from a sparse grid design \cite{Bungartz04}, where $n$ is varied depending on the order of the sparse grid. Sparse grid designs are broadly used for softening the curse-of-dimensionality in high-dimensional approximation problems \cite{Plumlee14}, thus our use of it here for high-dimensional boundary incorporation. With such data, we fit the tensor BdryMat\'ern GP (from Section \ref{sec:tensor_mat}) with smoothness parameters \(\nu=5/2\) and \(\nu=7/2\). As a benchmark, we then compare with a GP with a product Mat\'ern kernel with smoothness parameters $\nu = 5/2$ and $\nu=7/2$, which is the counterpart for the tensor BdryMat\'ern GP without boundary information. 

% We then compare tensor BdryGPs with smoothness parameters \(\nu=7/2\) and \(\nu=5/2\) to standard tensor Matérn GPs with the same smoothness parameters. The training data are collected using the Sparse Grid (SG) design \cite{Bungartz04,Plumlee14}, a commonly used experimental design for high-dimensional problems, and the data size for SG is represented by its level \(l\). We then train all competing GPs on the training set and compare all trained models based on log-MSE with \(10^4\) randomly collected test points.


Figure \ref{fig:tensor} (left) shows the corresponding log-MSE of the compared models, and Figure \ref{fig:tensor} (middle and right) visualizes the fitted predictive models (and its uncertainties) over a slice of the prediction domain. From Figure \ref{fig:tensor} (left), we again see that for both test functions, the integration of boundary information via the tensor BdryMat\'ern GP can considerably improve predictive performance over its product Mat\'ern GP counterpart, which doesn't incorporate such information. The different error decay slopes over sparse grid level also suggest that this integration of boundary information may offer improved prediction rates; investigating this will be a topic of future work. Figure \ref{fig:tensor} (middle and right) shows that, in addition to improved point predictions, the incorporation of known boundaries  via the tensor BdryMat\'ern GP can further reduce predictive uncertainties, thus facilitating confident surrogate modeling with limited data.

% all competitors and one-dimensional slices for the posteriors. For all test functions, the log-MSE results appear to decrease linearly with the SG level \(l\). Furthermore, the two BdryGP models yield much lower errors than the standard GP without boundary information, with the error decay slopes for BdryGP being roughly double those of the standard GP model. 


\begin{figure}[!t]
\centering
\includegraphics[width=0.32\textwidth]{figures/logerror/logErr_tensor_grie.pdf}
\includegraphics[width=0.32\textwidth]{figures/1Dslice/griewank_full.pdf}
\includegraphics[width=0.32\textwidth]{figures/1Dslice/griewank_no.pdf}

%\hspace{0.05cm}
%\includegraphics[width=0.32\textwidth]{figures/logerror/logErr_tensor_schwefel.pdf}
%\includegraphics[width=0.32\textwidth]{figures/1Dslice/Schwefel_full.pdf}
%\includegraphics[width=0.32\textwidth]{figures/1Dslice/Schwefel_no.pdf}
\hspace{0.05cm}
\includegraphics[width=0.32\textwidth]{figures/logerror/logErr_tensor_prodsin.pdf}
\includegraphics[width=0.32\textwidth]{figures/1Dslice/prod_sin.pdf}
\includegraphics[width=0.32\textwidth]{figures/1Dslice/prod_sin_mat.pdf}
\caption{(Left) Log-MSEs for the Griewank and product-sine test functions using the tensor BdryMat\'ern GP and the product Mat\'ern GP. (Middle) Posterior means and 68\% predictive intervals of the tensor BdryMat\'ern GP for the Griewank and product-sine functions. This is shown along a slice of the domain for input $x_1$ with other inputs set at their left end-points. (Right) Posterior means and 68\% predictive intervals of the product Mat\'ern GP for the Griewank and product-sine functions along the same domain slice.
}
% (bottom panels). Slice of tensor BdryMat\'ern GP predictors (middle column)  and tensor Mat\'ern predictors (last column)
\label{fig:tensor}
\end{figure}

% Intuition of the inferior performances of Mat\'ern GPs are illustrated in the middle and last columns of Figure \ref{fig:tensor}. Matérn GP exhibits mean reversion, which causes the posterior mean to revert to the prior mean outside the observed data domain. Consequently, this results in relatively large errors near the boundary points. Additionally, this mean reversion impacts the posterior variance, leading to incorrect variance estimates near the boundaries.

% In contrast, the BdryMatérn GP enforces fixed values or fixed derivatives at the boundary points. In the Dirichlet case, the posterior variances are zero at the boundaries. In the Neumann case, the posterior variances increase more slowly due to the influence of the Neumann boundary condition. These property helps to prevent the large errors observed with the standard Matérn GP and provides more accurate variance estimates near the boundaries. By ensuring that the boundary values are respected, the BdryMatérn GP can deliver more reliable predictions in the vicinity of the boundaries, thus mitigating the issues associated with mean reversion and boundary value misspecification. In high-dimensional problems, the vicinity of the boundaries occupies nearly the total volume of the rectangle \([0,1]^d\). As a result, accurate predictions near the boundaries can lead to significant improvements in averaged performance.

\section{Conclusion}
\label{sec:conclusion}
This paper introduces a new Gaussian process model called the BdryMat\'ern GP, which targets the reliable incorporation of known Dirchlet, Neumann or Robin boundary information on a irregular (i.e., non-hypercube) domain $\mathcal{X}$. The BdryMat\'ern GP makes use of a new BdryMat\'ern covariance kernel, which we derive in path integral form using the Feynman-Kac formula \citep{ito1957fundamental,stroock1971diffusion,papanicolaou1990probabilistic}. Provided $\mathcal{X}$ is connected with a boundary set that is twice-differentiable almost surely, we show that sample paths from the BdryMat\'ern GP provide smoothness control (in terms of its derivatives) while satisfying the desired boundaries. We then present an efficient approach for approximating the BdryMat\'ern kernel on general domains $\mathcal{X}$ via finite element modeling with rigorous error analysis. When the domain takes the form of a high-dimensional unit hypercube, we present a tensor form of the BdryMat\'ern GP that facilitates efficient boundary integration via a closed-form covariance kernel. Finally, we demonstrate the effectiveness of the proposed model for reliable boundary integration in a suite of numerical experiments that feature broad boundary information on different domain choices.

% The key innovation in BdryGP is the BdryMatérn covariance function, which retains the smoothness properties of a Matérn kernel while ensuring sample paths almost surely satisfy boundary conditions. We provide a closed-form solution for tensor BdryMatérn kernels and propose a regression framework for accurately approximating general BdryMatérn kernels. Numerical simulations confirm these results, demonstrating that incorporating boundary information significantly improves model performance.
 
% While this paper provides an appealing theoretical framework for the BdryGP model, there are further developments which would be useful for practical implementation. For computational efficiency, one can leverage the sparsity of B-splines or other potential basis functions to eliminate matrix computation steps for prediction and likelihood evaluations, which improves the scalability of BdryGP for big datasets. It would also be useful to investigate the behavior of BdryGP (e.g., consistency and convergence rates) under maximum likelihood estimation of model parameters.

Given the presented theoretical framework for boundary integration, there are several promising future directions for future work. One direction is the exploration of posterior contraction rates of the BdryMat\'ern GP on irregular domains $\mathcal{X}$, which can shed light on how much the integration of boundary information can improve probabilistic predictive performance. Another direction is the investigation of consistency and contraction rates for the BdryMat\'ern GP under inference (and potential misspecifications) of model parameters. Finally, it would be fruitful to investigate new space-filling experimental designs \citep{johnson1990minimax,mak2018minimax} that incorporate boundary information on irregular domains.





%%%%%%%%%%%%%%%%%%%%%%%%%%%%%%%%%%%%%%%%%%%%%%
%% Example with single Appendix:            %%
%%%%%%%%%%%%%%%%%%%%%%%%%%%%%%%%%%%%%%%%%%%%%%
% \begin{appendix}
% \section*{Title}\label{appn} %% if no title is needed, leave empty \section*{}.
% Appendices should be provided in \verb|{appendix}| environment,
% before Acknowledgements.

% If there is only one appendix,
% then please refer to it in text as \ldots\ in the \hyperref[appn]{Appendix}.
% \end{appendix}
%%%%%%%%%%%%%%%%%%%%%%%%%%%%%%%%%%%%%%%%%%%%%%
%% Example with multiple Appendixes:        %%
%%%%%%%%%%%%%%%%%%%%%%%%%%%%%%%%%%%%%%%%%%%%%%
% \begin{appendix}
% \section{Construction of functions satisfying boundary conditions}\label{appdix:bdry_func}
% In the main text, we have assume without loss of generality that the mean function is 0 and the boundary conditions are homogeneous, i.e.,
% \begin{align*}
%         &f(\BFx) = 0,\quad \forall\BFx \in \partial {\CalX}\quad\text{(Dirichlet)},\\
%         &\frac{\partial f(\BFx)}{\partial\bold{n}}+c(\BFx)f(\BFx) = 0,\quad \forall\BFx \in \partial {\CalX}\quad\text{(Robin)}
%     \end{align*}

% To accommodate an underlying system with the following general inhomogeneous boundary conditions:
% \begin{align}
% &f(\BFx) = h(\BFx),\quad \forall \BFx \in \partial {\CalX} \quad \text{(Dirichlet)},\label{eq:inhomo_diri}\\
% &\frac{\partial f(\BFx)}{\partial \bold{n}} + c(\BFx)f(\BFx) = h(\BFx),\quad \forall \BFx \in \partial {\CalX} \quad \text{(Robin)},\label{eq:inhomo_robin}
% \end{align}
% the standard approach is to construct a mean function $f$ that satisfies the prescribed boundary conditions. Then, by the superposition principle of homogeneous and inhomogeneous boundary conditions, a GP with this mean function automatically inherits the inhomogeneous boundary condition.


% We suggest constructing the mean function as the solution to the deterministic counterpart of the SPDE \eqref{eq:SPDE}:
% $$
% (\kappa^2 - \bigtriangleup) f = 0,
% $$
% subject to either the inhomogeneous Dirichlet boundary condition \eqref{eq:inhomo_diri} or the inhomogeneous Robin boundary condition \eqref{eq:inhomo_robin}. Consequently, determining the mean function $f$ reduces to solving a linear PDE with standard boundary conditions—an extensively studied problem in the field of numerical PDEs and can be easily solved by current numerical solvers.

% In this article, we adopt one of the most widely used numerical techniques: the differential quadrature (DQ) method \cite{liu2011kernel, shu2003local, shu2012differential}, which is also the method for constructing B-spline basis functions satisfying the target boundary conditions in Section \ref{sec:Bspline}. DQ is a discretization method for approximating derivatives, where the partial derivative of an unknown function with respect to an independent variable is represented as a weighted linear combination of the function values at a set of discrete points within its domain of support.

%  We take the Robin boundary as an example. Suppose we have a set of basis functions $\{\phi_i\}_{i=1}$  whose supports are compact and cover the boundary $\partial {\CalX}$. Also, suppose that we have a set of boundary points $\{\bold{b}_i\}_{i=1}$ that can discretize the boundary $\partial\CalX$.  The local DQ method construct a function as a linear combination of $\{\phi_i\}_{i=1}$: $f=\sum_{i=1} f_i \phi_i$, so 
% \begin{equation}
%     \label{eq:DQ_1}
%     \frac{\partial f(\bold{b}_i)}{\partial \bold{n}}=\sum_{j\in {\rm Ne}[\bold{b}_i]}\bold{w}_j^T\bold{n} f(\bold{b}_j)
% \end{equation}
% and, at each boundary point $\bold{b}_i$,  the normal derivative of $f$ at $\bold{b}_i$ is approximated by its neighboring values:
% \[\frac{\partial f(\bold{b}_i)}{\partial \bold{n}}=\sum_{j\in {\rm Ne}[\bold{b}_i]}\bold{w}_j^T\bold{n} f(\bold{b}_j)\]
% so $f$ satisfies the inhomogeneous Robin boundary conditions at each  $\bold{b}_i$:
% \begin{equation}
% \label{eq:DQ_2}
%     \sum_{j\in {\rm Ne}[\bold{b}_i]}\bold{w}_j^T\bold{n} f(\bold{b}_j)+c(\bold{b}_j)f(\bold{b}_j)=h(\bold{b}_j),\quad \forall j=1,\cdots,m.
% \end{equation}
% Solving the linear system \eqref{eq:DQ_1} and \eqref{eq:DQ_2} yields a mean function $f$ that satisfies the inhomogeneous Robin boundary condition. Notice that the problem with Dirichlet boundaries is a special case of \eqref{eq:DQ_2} with all coefficients $\bold{w}_j$ equal 0. 



%\section{Title of the second appendix}\label{appB}
%\subsection{First subsection of Appendix \protect\ref{appB}}

%Use the standard \LaTeX commands for headings in \verb|{appendix}|.
%Headings and other objects will be numbered automatically.
%\begin{equation}
%\mathcal{P}=(j_{k,1},j_{k,2},\dots,j_{k,m(k)}). \label{path}
%\end{equation}

%Sample of cross-reference to the formula (\ref{path}) in Appendix \ref{appB}.
% \end{appendix}

%%%%%%%%%%%%%%%%%%%%%%%%%%%%%%%%%%%%%%%%%%%%%%
%% Support information, if any,             %%
%% should be provided in the                %%
%% Acknowledgements section.                %%
%%%%%%%%%%%%%%%%%%%%%%%%%%%%%%%%%%%%%%%%%%%%%%
% \begin{acks}[Acknowledgments]
% The authors would like to thank the anonymous referees, an Associate
% Editor and the Editor for their constructive comments that improved the
% quality of this paper.
% \end{acks}

%%%%%%%%%%%%%%%%%%%%%%%%%%%%%%%%%%%%%%%%%%%%%%
%% Funding information, if any,             %%
%% should be provided in the                %%
%% funding section.                         %%
%%%%%%%%%%%%%%%%%%%%%%%%%%%%%%%%%%%%%%%%%%%%%%
\begin{funding}
The second author gratefully recognizes funding from NSF CSSI Frameworks 2004571, NSF DMS 2210729, 2316012 and U.S. Department of Energy Grant DE-SC0024477.
\end{funding}

%%%%%%%%%%%%%%%%%%%%%%%%%%%%%%%%%%%%%%%%%%%%%%%%%%%%%%%%%%%%%
%%                  The Bibliography                       %%
%%                                                         %%
%%  imsart-???.bst  will be used to                        %%
%%  create a .BBL file for submission.                     %%
%%                                                         %%
%%  Note that the displayed Bibliography will not          %%
%%  necessarily be rendered by Latex exactly as specified  %%
%%  in the online Instructions for Authors.                %%
%%                                                         %%
%%  MR numbers will be added by VTeX.                      %%
%%                                                         %%
%%  Use \cite{...} to cite references in text.             %%
%%                                                         %%
%%%%%%%%%%%%%%%%%%%%%%%%%%%%%%%%%%%%%%%%%%%%%%%%%%%%%%%%%%%%%

%% if your bibliography is in bibtex format, uncomment commands:
\bibliographystyle{imsart-number} % Style BST file (imsart-number.bst or imsart-nameyear.bst)
\bibliography{bibliography}       % Bibliography file (usually '*.bib')

%% or include bibliography directly:

\pagebreak


%%%%%%%%%%%%%%%%%%%%%%%%%%%%%%%%%%%%%%%%%%%%%%
%% Supplementary Material, including data   %%
%% sets and code, should be provided in     %%
%% {supplement} environment with title      %%
%% and short description. It cannot be      %%
%% available exclusively as external link.  %%
%% All Supplementary Material must be       %%
%% available to the reader on Project       %%
%% Euclid with the published article.       %%
%%%%%%%%%%%%%%%%%%%%%%%%%%%%%%%%%%%%%%%%%%%%%%
\begin{appendix}
\stitle{}

% \subsection{Proof of Theorem \ref{thm:empirical_BdryMat_Robin}}

\section{Proof of Theorem \ref{thm:smoothdir}}
\label{pf:smoothdir}
\begin{proof}[Proof of ${(a)}$]
Given our assumption that $\CalX$ is compact and connected with boundary set twice-differentiable almost everywhere, the GP with Mat\'ern covariance $k_{\nu}$ admits an eigendecomposition with respect to the Lebesgue measure on $\CalX$ of the form:
\begin{equation}
    \label{eq:smoothdir_pf_1}
    \CalW_\nu(\BFx)=\sum_i \sqrt{\lambda_i}\xi_i(\BFx)Z_i,
\end{equation}
where $\{Z_i\}_i$ are i.i.d. standard normal random variables, $\{\lambda_i\}_i$ and $\{\xi\}_i$ are the eigenvalues and eigenfunctions of $k_{\nu}$ under the Lebesgue measure. Moreover, since:
\[\int_\CalX k_\nu(\Bx,\Bx)d\Bx=\sum_i\lambda_i^2\int_\CalX\xi^2_i(\Bx)d\Bx\leq k_\nu(0) {\rm Vol}(\CalX),\]
it follows that $\sqrt{\lambda_i}\xi_i\in L^2(\CalX)$ for any $i$. Let $g_i$ be the solution to the following elliptical PDE:
\begin{equation}
    \label{eq:smoothdir_pf_2}
    (\kappa^2-\bigtriangleup)g_i=\sqrt{\lambda_i}\xi_i,  \quad g_i(\BFx)=0, \quad \forall\BFx \in\partial \CalX.
\end{equation}
Thus, given that $\partial \mathcal{X}$ is twice-differentiable almost everywhere and \cite[Theorem 4, Chapter 6]{Evans15}, we have that $g_i$ is $L^2$ twice-differentiable and $g_i=0$ on the boundary.

For finite $i^*$, define:
\[f_{i^*}=\sum_{i=1}^{i^*} g_iZ_i,\quad \CalW_{\nu,i^*}= \sum_{i=1}^{i^*} \sqrt{\lambda_i}\xi_i(\BFx)Z_i.\]
By the superposition principle of linear PDEs, we have:
\[   (\kappa^2-\bigtriangleup)\CalW_{\nu,i^*}=\sum_{i=1}^{i^*} Z_i   (\kappa^2-\bigtriangleup)\sqrt{\lambda_i}\xi_i=f_{i^*},\]
so $f_{i^*}$ satisfies the boundary condition almost surely for any $i^*$. Then by Theorem 12.1 in \cite{lifshits2012lectures}, we have:
\[f=\lim_{i^*\to\infty}f_{i^*},\quad {\rm a.s.,}\]
where $f$ is the solution to \eqref{eq:bdry_SPDE}. So $f$ satisfies the boundary condition almost surely.
\end{proof}

\begin{proof}[Proof of ${(b)}$] 
    We shall prove a stronger version of the theorem. We first define the concept of $\nu$-almost H\"older continuity as follows:
    \begin{definition}
        Let $C^\gamma(\CalX)$ be the $\gamma$-H\"older space. The $\nu$-almost H\"older space on $\CalX$ is defined as
        \[C^{-\nu}(\CalX)=\cap_{\gamma<\nu}C^{\gamma}(\CalX).\]
    \end{definition}
    From \cite[Proposition 10]{da2023sample}, we know that for any Mat\'ern-$\nu$ GP $\CalW_\nu$, its sample paths are $\nu$-almost H\"older. So the solution to \eqref{eq:bdry_SPDE} satisfies:
    \[\bigtriangleup f\in C^{-\nu}(\CalX),\quad {\rm a.s.}\]
    According to the definition of $\bigtriangleup=\sum_{l=1}^d\partial_{x_lx_l}$, we can conclude that $f\in C^{-(\nu+2)}(\CalX)$ almost surely. Hence, for any non-integer $\nu$, $f$ is $\lfloor\nu+2\rfloor$-differentiable and for any integer $\nu$, $f$ is $\nu+1$-differentiable.
\end{proof}


\section{Proof of Theorem \ref{thm:smooth}}
\label{pf:smoothrobin}
\begin{proof}
    The proof is identical to the proof of Theorem \ref{thm:smooth} except that \eqref{eq:smoothdir_pf_2} is replaced by
    \begin{equation}
    \label{eq:smooth_pf_1}
    (\kappa^2-\bigtriangleup)g_i=\sqrt{\lambda_i}\xi_i,  \quad \frac{\partial g_i(\Bx)}{\partial \bold{n}}+c(\mathbf{x})g_i(\Bx)=0, \quad \forall\BFx\in\CalX. 
\end{equation}
However, the change in boundary conditions does not affect our conclusions; all the theorems remain valid under Robin boundary conditions.
    
\end{proof}

%\section{Proof of Corollary \ref{coro:posotive_hat_K}}  Corollary \ref{coro:posotive_hat_K} is obvious, we skip the proof.


\section{Proof of Theorem \ref{thm:convergence_general}}
\label{pf:convergence_general}

To prove the theorem, we first start with some useful lemmas. The first Lemma can be found in Section 5 of \cite{tropp2012user}:

\begin{lemma}[Chernoff inequality of random matrix]
\label{lem:matrix_chernoff}
    Let $p\in\NatInt$ and $\delta\in[0,1)$ be arbitrary and let $\bold{M}_1,\cdots, \bold{M}_m\in\Real^{p\times p}$ be independent, symmetric and positive semidefinite matrices with random entries. Let $R>0$ be such that $\lambda_{\rm max}(M_i)\leq R$ holds for all $i=1,...,m$. Then,it holds that:
    \begin{align*}
        &\pr\left(\lambda_{\rm min}(\sum_{i=1}M_i)\leq (1-\delta)\underline{\lambda}\right)\leq p\left(\frac{e^{-\delta}}{(1-\delta)^{1-\delta}}\right)^{\underline{\lambda}/R}\\
        &\pr\left(\lambda_{\rm max}(\sum_{i=1}M_i)\geq (1+\delta)\overline{\lambda}\right)\leq p\left(\frac{e^{\delta}}{(1+\delta)^{1+\delta}}\right)^{\overline{\lambda}/R},
    \end{align*}
    where $\lambda_{\max}(M)$ and $\lambda_{\min}(M)$ denote the maximum and minimum eigenvalues of a matrix $M$, respectively, and $\overline{\lambda}=\lambda_{\rm max}(\E \sum_{i=1}M_i)$ and $\underline{\lambda}=\lambda_{\rm min}(\E \sum_{i=1}M_i)$.
\end{lemma}

From Lemma \ref{lem:matrix_chernoff}, we can derive the following lemma in a straightforward way.
\begin{lemma}
\label{lem:spectral_empirical_gram}
    Let $M=[\int_{{\CalX}}\phi_i(\BFx)\phi_j(\BFx)d\BFx]$ be the Gram matrix of basis functions $\{\phi_i\}_{i=1}^p$ and let $\underline{\lambda}=\lambda_{\rm min}(M)$ and $\overline{\lambda}=\lambda_{\rm max}(M)$. Then:
    \begin{align}
        &\pr\left(\lambda_{\rm min}(\frac{1}{m}\sum_{i=1}\BPhi(\Bu_i)\BPhi(\Bu_i)^T)\leq \frac{\underline{\lambda}}{2} \quad \textrm{or} \quad \lambda_{\rm max}(\frac{1}{m}\sum_{i=1}\BPhi(\Bu_i)\BPhi(\Bu_i)^T)\geq \frac{3\overline{\lambda}}{2}\right)\nonumber\\
        \leq & p(\sqrt{2}e^{-1/2})^{m\underline{\lambda}/S_\phi}+p(\frac{e^{1/2}}{(3/2)^{3/2}})^{m\overline{\lambda}/S_\phi}. \label{eq:spectral_empirical_gram}
    \end{align}
\end{lemma}
\proof
Let $M_i$ denote the i.i.d. sample of Gram matrix $\frac{1}{m}\BPhi(\Bu_i)\BPhi^T(\Bu_i)$. Note that $\lambda_{\rm max}(M_1)\leq {S_\phi}/{m}$ almost surely because:
\[\lambda_{\rm max}(M_i)=\frac{1}{m}\BPhi(\Bu_i)^T\BPhi(\Bu_i)=\frac{1}{m}\sum_{j=1}\phi_j(\Bu_i)\phi_j(\Bu_i)\leq \frac{S_\phi}{m}.\]
Applying Lemma \ref{lem:matrix_chernoff} with $R=S_\phi/m$ and $\delta=1/2$, we can obtain the result.
\endproof





\proof[Proof of Theorem \ref{thm:convergence_general}]
Let $\CalP$ denote the $L_2$ projection onto the $\CalF_p=\text{span}\{\phi_i\}_{i=1}^p$. Similarly,
 the operator $\BPhi(\cdot)^T\BPhi(\BU)^\dagger$ can then be viewed as an empirical projection onto  $\CalF_p$ induced by $\BU$. Let $\CalP_m$ denote this empirical projection. For brevity, let $k_{\nu}$ denote the Mat\'ern-$\nu$ kernel in the following content.
 
 Based on the eigenfunction decomposition \eqref{lem:eigen_mat} and the path integral representation of BdryMat\'ern kernel (\eqref{eq:bdryMatern_path_integral_dirichlet} for Dirichlet setting; \eqref{eq:bdryMatern_path_integral_robin} for Robin setting), the BdryMat\'ern kernel and its regression kernel can be written as:
 \begin{align}
      k_{\nu,\mathcal{B}}=&\sum_i\lambda_ie_i(\BFx)e_i(\BFx')-\sum_i\lambda_ie_i(\Bx)\E[G_i(\BFx')]\nonumber\\
     &-\sum_i\lambda_ie_i(\Bx')\E[G_i(\Bx)]+\sum_i\lambda_i\E[G_i(\Bx)]E[G_i(\Bx')],\label{eq_pf_general_convergence_1}
 \end{align}
 and:
 \begin{align}
     \hat{k}_{Q,m}(\Bx,\Bx')=&\sum_i\lambda_i\CalP_m[e_i](\BFx)\CalP_m[e_i](\BFx')-\sum_i\lambda_i\CalP_m[e_i](\Bx)\CalP_m[G_i](\BFx')\nonumber\\
     &-\sum_i\lambda_i\CalP_m[e_i](\Bx')\CalP_m[G_i](\BFx)\nonumber\\
     &+\sum_i\lambda_i\CalP_m[G_i](\Bx)\CalP_m[G_i](\BFx'),\label{eq_pf_general_convergence_2}
 \end{align}
where:
\begin{equation*}
    G_i(\BFx)\coloneqq\begin{cases}
        &e^{-\kappa^2\tau} e_i\circ \BB_\tau\quad \quad \quad \quad \quad \quad \quad \quad \ \text{(Dirichlet)}\\
        & \int_0^\infty e^{-\kappa^2t-\int_0^tc(\BB_s)dLs}e_i\circ\BB_t dL_t\quad \text{(Robin)} 
    \end{cases}
\end{equation*}
 can be viewed as  random functions with $\BB_0=\BFx$. To construct our estimator $\hat{k}_{Q,m}$, we have one realizations $G_{i}(\Bu_j)$ for each $\Bu_j$. %Notice that our estimator is biased. This can be seem from the term $\sum_i\lambda_i\CalP_m[G_i](\Bx)\CalP_m[G_i](\BFx')$ of $\hat{K}$:
%\begin{align*}
%     \sum_i\lambda_i\CalP_m[G_i](\Bx)\CalP_m[G_i](\BFx')=\CalP_m \left(\sum_i\lambda_i G_iG_i\right)\CalP_m^T(\BFx,\Bx').
%\end{align*}
%So  the diagonal $(\Bu_i,\Bu_i)$ is biased
%\begin{align*}
%&\sum_i\lambda_i\E[G_i(\Bu_i)G_i(\Bu_i)]-\sum_i\lambda_i\E[G_i(\Bu_i)]\E[G_i(\Bu_i)]\\
%    =& k(0)\E[\BB_\tau e^{-2\kappa^2\tau}|\BB_0=\Bu_i]-\E[k(\BB_\tau,\BB'_\gamma)e^{-\kappa^2(\tau+\gamma)}|\BB_0=\BB'_0=\Bu_i]>0
%\end{align*}
%where $\BB_t$ and $\BB'_t$  are independent Brownian motions with initial conditions $\BB_0=\BB'_0=\Bu_i$, $\tau$ and $\gamma$ are the hitting times of $\BB_t$ and $\BB'_t$  to the boundary $\partial\CalX$, respectively. Nevertheless, we will show later in this proof that the bias will shrink to 0 if we can properly choose inducing points and basis functions.

Let $\CalU_*$ denote the following event: 
\[\CalU_*=\left\{\lambda_{\rm min}(\frac{1}{m}\sum_{i=1}M_i)\geq \frac{\underline{\lambda}}{2} \quad \textrm{and} \quad \lambda_{\rm max}(M_i)\leq \frac{3\overline{\lambda}}{2}\right\},\]
where, following the notation from Lemma \ref{lem:spectral_empirical_gram},  $M_i=\BPhi(\Bu_i)\BPhi(\Bu_i)^T$ denotes the sample Gram matrix, and $\underline{\lambda}=\lambda_{\min}(M)$ and $\overline{\lambda}=\lambda_{\max}(M)$ denote the minimum and maximum eigenvalues of the Gram matrix $[M]_{i,j}=\langle\phi_i,\phi_j\rangle_{L^2}$, respectively.


From Lemma \ref{lem:spectral_empirical_gram}, it follows that the probability of $\CalU_*$ is at least:
\[\pr\left(\CalU_*\right)\geq 1-p(\sqrt{2}e^{-1/2})^{m\underline{\lambda}/S_\phi}-p(\frac{e^{1/2}}{(3/2)^{3/2}})^{m\overline{\lambda}/S_\phi}.\]


We first use \eqref{eq_pf_general_convergence_1} and \eqref{eq_pf_general_convergence_2} to have the following decomposition:
\begin{align}
    &\hat{k}_{m,Q}(\BFx,\BFx')- k_{\nu,\mathcal{B}}(\BFx,\BFx')=\CalK_1(\BFx,\BFx')-\CalK_2(\BFx,\BFx')-\CalK_2(\BFx',\BFx)+\CalK_3(\BFx,\BFx'),\label{eq_pf_general_convergence_4}
\end{align}   
where:
\begin{align*}
    &\CalK_1(\BFx,\BFx')= \sum_i\lambda_i\left(\CalP_m[e_i](\Bx)\CalP_m[e_i](\BFx')-e_i(\BFx)e_i(\BFx')\right),\\
    &\CalK_2(\BFx,\BFx')=\sum_i\lambda_i\left(\CalP_m[e_i](\BFx)\CalP_m[G_i](\BFx')-e_i(\BFx)\E[G_i(\BFx')]\right),\\
     &\CalK_3(\BFx,\BFx')=\sum_i\lambda_i\left(\CalP_m[G_i](\BFx)\CalP_m[G_i](\BFx')-\E[G_i(\BFx)]\E [G_i(\Bx')]\right).
     %&\CalK_4(\BFx,\BFx')= \CalP_m\text{diag}\left[\sum_i\lambda_i\E[G_i(\Bu_i)G_i(\Bu_i)]-\E[G_i(\Bu_i)]\E[G_i(\Bu_i)]\right]\CalP_m^T(\BFx,\BFx').
\end{align*}



Write $k=k_{\nu,\mathcal{B}}$ for notation simplicity. Substitute \eqref{eq_pf_general_convergence_4} into the $L^2$ error between $\hat{k}_{m,Q}$ and $k$, we have on the event $U_*$
\begin{align}
   \|\hat{k}_{m,Q}-k\|_{L^2(\CalX\times\CalX)}^2 /4 \leq& \underbrace{\|\CalK_1\|_{L^2(\CalX\times\CalX)}^2}_{A_1}+\underbrace{2\|\CalK_2\|_{L^2(\CalX\times\CalX)}^2}_{A_2}+\underbrace{\|\CalK_3\|_{L^2(\CalX\times\CalX)}^2}_{A_3} .\label{eq_pf_general_convergence_5}
\end{align}

{\large \textbf{Estimation of $A_1$}}: $A_1$ is independent of boundary condition. On the set $\CalU_*$,
\begin{align}
\|\CalK_1\|_{L^2(\CalX\times\CalX)}^2=&\int_{\CalX\times\CalX}\bigg|\sum_i\lambda_i\big[\CalP_m[e_i](\Bx)\CalP_m[e_i](\BFx')-\CalP_m[e_i](\Bx)e_i(\Bx')\nonumber\\
    &+\CalP_m[e_i](\Bx)e_i(\Bx')-e_i(\BFx)e_i(\BFx')\big]\bigg|^2d\BFx d\Bx'\nonumber\\
    \leq &2\int_{\CalX\times\CalX}\bigg|\sum_i\lambda_i\CalP_m[e_i](\Bx)\big[\CalP_m[e_i](\BFx')-e_i(\Bx')\big]\bigg|^2d\Bx d\Bx'\nonumber\\
    &+2\int_{\CalX\times\CalX}\bigg|\sum_i\lambda_ie_i(\Bx')\big[\CalP_m[e_i](\BFx)-e_i(\Bx)\big]\bigg|^2d\Bx d\Bx'\nonumber\\
    \leq &2\underbrace{\int_\CalX\sum_i\lambda_i\big|\CalP_m[e_i](\BFx)\big|^2d\BFx }_{A_{11}}\underbrace{\int_\CalX\sum_i\lambda_i\big|\CalP_m[e_i](\BFx')-e_i(\Bx')\big|^2d\BFx'}_{A_{12}}\nonumber\\
    &+2\underbrace{\int_\CalX\sum_i\lambda_i\big|e_i(\BFx)\big|^2d\BFx }_{A_{13}}\underbrace{\int_\CalX\sum_i\lambda_i\big|\CalP_m[e_i](\BFx)-e_i(\Bx)\big|^2d\BFx }_{A_{12}}.\label{eq_pf_general_convergence_6}
\end{align}
For the term $A_{11}$, transform it back to the kernel representation, we can obtain:
\begin{align*}
    A_{11}=&\int_\CalX \BPhi(\BFx)^T\left(\frac{1}{m}\sum_iM_i\right)^{-1}\frac{\BPhi(\BU)}{m}k_\nu(\BU,\BU)\frac{\BPhi(\BU)^T}{m}\left(\frac{1}{m}\sum_{i=1}M_i\right)^{-1}\BPhi(\Bx)d\BFx\\
    =&\text{Tr}\left[\frac{1}{m}{\BPhi(\BU)^T}\left(\frac{1}{m}\sum_{i=1}M_i\right)^{-1}M\left(\frac{1}{m}\sum_iM_i\right)^{-1}{\BPhi(\BU)}\frac{k_\nu(\BU,\BU)}{m}\right]\\
    \leq & \text{Tr}[\frac{k_\nu(\BU,\BU)}{m}]\lambda_{max}\left(\frac{1}{m}{\BPhi(\BU)^T}\left(\frac{1}{m}\sum_{i=1}M_i\right)^{-1}M\left(\frac{1}{m}\sum_iM_i\right)^{-1}{\BPhi(\BU)}\right)\\
    \leq & C_1 \frac{\overline{\lambda}^2}{\underline{\lambda}^2},
\end{align*}
where the last line is because on event $U_*$, $\lambda_{\rm min}(\frac{1}{m}\sum_{i=1}M_i)\geq \frac{\underline{\lambda}}{2}$,  $ \lambda_{\rm max}(M_i)\geq \frac{3\overline{\lambda}}{2}$, and $\text{Tr}[\frac{k_\nu(\BU,\BU)}{m}]=k(0)$ for the Mat\'ern kernel.

For the term $A_{12}$, notice that the empirical projection $\CalP_m$ is invariant under the $L^2$ projection $\CalP$, i.e., $\CalP_m\CalP=\CalP$. We can obtain on $\CalU_*$:
\begin{align*}
   A_{12}\leq &\int_\CalX\sum_{i}\lambda_i\left(2|\CalP_m[e_i-\CalP e_i](\BFx)|^2+2|\CalP[e_i](\BFx)-e_i(\BFx)|^2\right)d\BFx\\
    =&\sum_i2\lambda_i \int_\CalX \|\BPhi^T(\BFx)(\frac{1}{m}\sum_{l=1}M_l)^{-1}\frac{1}{m}\BPhi(\BU)[e_i(\BU)-\CalP[e_i](\BU)]\|^2d\Bx\\
    &+2\sum_i\lambda_i\int_\CalX \left|\CalP[e_i](\BFx)-e_i(\BFx)\right|^2d\BFx\\
    \leq & \sum_iC_2\lambda_i \frac{\overline{\lambda}}{\underline{\lambda}^2}\|\frac{1}{m}\BPhi(\BU)[e_i(\BU)-\CalP[e_i](\BU)]\|^2+2\sum_i\lambda_i\int_\CalX \left|\CalP[e_i](\BFx)-e_i(\BFx)\right|^2d\BFx.
\end{align*}
 On $\CalU_*$, we have:
\begin{align*}
  \|\frac{1}{m}\BPhi(\BU)[e_i(\BU)-\CalP[e_i](\BU)]\|^2= & \sum_{j=1}^p\left(\frac{1}{m}\sum_{l=1}\phi_j(\Bu_l)\left[e_i(\Bu_l)-\CalP[e_i](\Bu_l)\right]\right)^2\\
  \leq & \|e_i-\CalP e_i\|_n^2\frac{1}{m}\sum_{j=1}^p\sum_{l=1}\phi_j(\Bu_l)^2\\
  \leq & \frac{S_\phi}{m}\left(\frac{1}{m}\sum_{j=1}\left|e_i(\Bu_j)-\CalP e_i(\Bu_j)\right|^2\right),
\end{align*}
where the second line is from H\"older inequality. Also, notice that the infinite sum of $|e_i(\BFx)-\CalP[e_i](\BFx)|^2$ can be represented as operators acting on kernel $k$:
\begin{align*}
    \sum_i\lambda_i |e_i(\BFx)-\CalP[e_i](\BFx)|^2=& \sum_i\lambda_i \left(\bold{I}-\CalP\right)e_ie_i\left(\bold{I}-\CalP\right)^T(\BFx,\BFx)\\
    =&\left(\bold{I}-\CalP\right)\left( \sum_i\lambda_i e_ie_i\right)\left(\bold{I}-\CalP\right)^T(\BFx,\BFx)\\
    =&\left(\bold{I}-\CalP\right)k_\nu(\cdot,\cdot)\left(\bold{I}-\CalP\right)^T(\BFx,\BFx).
\end{align*}
So $A_{11}\cdot A_{12}$ is bounded as:
\begin{align*}
    & \int_\CalX\sum_i\lambda_i\big|\CalP_m[e_i](\BFx)\big|^2d\BFx \int_\CalX\sum_i\lambda_i\big|\CalP_m[e_i](\BFx')-e_i(\Bx')\big|^2d\BFx' \\
    \leq&C_3(\frac{\overline{\lambda}}{\underline{\lambda}})^2\bigg[\frac{\overline{\lambda}S_\phi}{\underline{\lambda}^2m}\left(\frac{1}{m}\sum_{j=1}\left(\bold{I}-\CalP\right)k_\nu(\cdot,\cdot)\left(\bold{I}-\CalP\right)^T(\Bu_j,\Bu_j)\right)\\
    &+\int_\CalX\left(\bold{I}-\CalP\right)k_\nu(\cdot,\cdot)\left(\bold{I}-\CalP\right)^T(\BFx,\BFx)d\Bx\bigg].
\end{align*}


For $A_{13}$, it is straightforward to check that it is:
\[\int_\CalX d\BFx k(0)=k(0)\text{vol}(\CalX).\]
Substitute the above identities of $A_{11}$, $A_{12}$, and $A_{13}$ into \eqref{eq_pf_general_convergence_6}, we then have:
\begin{align}
\|\CalK_1\|_{L^2(\CalX\times\CalX)}^2\leq&C_4(\frac{\overline{\lambda}}{\underline{\lambda}})^2\bigg[\frac{\overline{\lambda}S_\phi}{\underline{\lambda}^2m}\left(\frac{1}{m}\sum_{j=1}\left(\bold{I}-\CalP\right)k_\nu(\cdot,\cdot)\left(\bold{I}-\CalP\right)^T(\Bu_j,\Bu_j)\right)\label{eq_pf_general_convergence_7}\\
    &+\int_\CalX\left(\bold{I}-\CalP\right)k_\nu(\cdot,\cdot)\left(\bold{I}-\CalP\right)^T(\BFx,\BFx)d\Bx\bigg].\nonumber
\end{align}

{\large \textbf{Estimation of $A_2$}}: For $A_2$, first decompose $\CalK_2$ as follows:
\begin{align*}
    \CalK_2(\Bx,\Bx')=\CalK_{21}(\Bx,\Bx')+\CalK_{22}(\Bx,\Bx')+\CalK_{23}(\Bx,\Bx'),
\end{align*}
where:
\begin{align*}
    &\CalK_{21}(\Bx,\Bx')=\sum_i\lambda_i\left(\CalP_m[e_i](\Bx)\CalP_m[G_i](\BFx')-\CalP_m[e_i](\BFx)\CalP_m\E[G_i](\BFx')\right)\\
    &\CalK_{22}(\Bx,\Bx')=\sum_i\lambda_i\left(\CalP_m[e_i](\BFx)\CalP_m\E[G_i](\BFx')-\CalP_m[e_i](\BFx)\E[G_i(\BFx')]\right)\\
    &\CalK_{23}(\Bx,\Bx')=\sum_i\lambda_i\left(\CalP_m[e_i](\BFx)\E[G_i(\BFx')]-e_i(\BFx)\E[G_i(\BFx')]\right).
\end{align*}
So $A_2$ can be further decomposed as:
\begin{equation}
    A_2/4\leq \underbrace{\|\CalK_{21}\|_{L^2(\CalX\times\CalX)}^2}_{A_{21}}+\underbrace{\|\CalK_{22}\|_{L^2(\CalX\times\CalX)}^2}_{A_{22}}+\underbrace{\|\CalK_{23}\|_{L^2(\CalX\times\CalX)}^2}_{A_{23}}. \label{eq_pf_general_convergence_8}
\end{equation}

For $A_{21}$:
\begin{align*}
    \|\CalK_{21}\|_{L^2(\CalX\times\CalX)}^2\leq &\int_{\CalX}\sum_i\lambda_i\left|\CalP_m[e_i](\BFx)\right|^2d\BFx\int_{\CalX}\sum_i\lambda_i\left|\CalP_m[ \E[G_i]-G_i](\BFx)\right|^2d\BFx\\
    = & A_{11} \int_{\CalX} \sum_i\lambda_i \|\BPhi(\Bx)^T(\frac{1}{m}\sum_{i=1}M_i)^{-1}\BPhi(\BU)\boldsymbol{\varepsilon}_i\|^2d\Bx,
\end{align*}
where $\boldsymbol{\varepsilon}$ are independent distributed zero-mean observation noise $\boldsymbol{\varepsilon}_{ij}=\E_{G_i}[G_i(\Bu_j)]-G_i(\Bu_j)$. Analogously to the estimation of $A_{12}$, 
\begin{align*}
    \int_{\CalX} \sum_i\lambda_i \|\BPhi(\Bx)^T(\frac{1}{m}\sum_{i=1}M_i)^{-1}\BPhi(\BU)\boldsymbol{\varepsilon}_i\|^2d\Bx\leq C_2\frac{\overline{\lambda}}{\underline{\lambda}^2}\sum_i\lambda_i\|\frac{1}{m}\BPhi(\BU)\boldsymbol{\varepsilon}_i\|^2.
\end{align*}
Because noise $\boldsymbol{\varepsilon}$ are zero-mean, we have:
\begin{align*}
    \sum_i\lambda_i\ \|\frac{1}{m}\BPhi(\BU)\boldsymbol{\varepsilon}_i\|^2 
    =& \sum_i \lambda_i \sum_{j=1}^p\left(\frac{1}{m}\sum_{l=1}\phi_j(\Bu_l)\boldsymbol{\varepsilon}_{il}\right)^2\\
    \leq & \sum_i\lambda_i \sum_{j=1}^p\left(\frac{1}{m}\sum_{l=1}\phi_j(\Bu_l)^2\right)\left(\frac{1}{m}\sum_{l=1}\boldsymbol{\varepsilon}_{il}^2\right)\\
    \leq & \frac{S_\phi}{m}\frac{1}{m}\sum_{l=1}\sum_i\lambda_i \boldsymbol{\varepsilon}_{il}^2.
\end{align*}

For the Dirichlet boundary setting:
\begin{align*}
    \sum_i\lambda_i \boldsymbol{\varepsilon}_{il}^2 = &  \sum_i\lambda_i\left(\E[e_i(\BB_{\tau})e^{-\kappa^2\tau}|\BB_0=\Bu_l]- e_i(\BB_{\tau})e^{-\kappa^2\tau}|\BB_0=\Bu_l\right)^2\\
    =&  \E_{\BB_t,\BB'_t,\tau,\gamma}[k_\nu(\BB_\tau,\BB'_\gamma)e^{-\kappa^2(\tau+\gamma)}|\BB_0=\BB'_0=\Bu_l]\\
    &+k_\nu(\BB_\tau,\BB_\tau)e^{-2\kappa\tau}|\BB_0=\Bu_l\\
    &-2\E_{\BB_t,\tau}[k_\nu(\BB_\tau,\BB'_\gamma)e^{-\kappa^2(\tau+\gamma)}|\BB_0=\BB'_0=\Bu_l]
    \leq 4k(0),
\end{align*}
where the last line is from the fact that Mat\'ern kernel satisfies $k_\nu(0)\geq k_\nu(\Bx,\Bx')$ for any $\Bx,\Bx'$.

For the Robin boundary setting:
\begin{align*}
   \sum_i\lambda_i \boldsymbol{\varepsilon}_{il}^2 = & \E \int_0^\infty \int_0^\infty e^{-\kappa^2(t+\tau)-\int_0^tc(\BB_s)dL_s-\int_0^\tau c(\BB'_s)dL'_s}k_\nu(\BB_t,\BB'_\tau)dL_tdL_\tau\\
   &+\int_0^\infty \int_0^\infty e^{-\kappa^2(t+\tau)-\int_0^tc(\BB_s)dL_s-\int_0^\tau c(\BB'_s)dL'_s}k_\nu(\BB_t,\BB'_\tau)dL_tdL_\tau\\
    &-2\E_{\BB_t,L_t}\int_0^\infty \int_0^\infty e^{-\kappa^2(t+\tau)-\int_0^tc(\BB_s)dL_s-\int_0^\tau c(\BB'_s)dL'_s}k_\nu(\BB_t,\BB'_\tau)dL_tdL_\tau\\
    &\leq 4k_\nu(0),
\end{align*}
where the last line is from \eqref{eq:pf_empirical_BdryMat_robin_1}.

Therefore, in both cases,  we always have:
\begin{align*}
     \sum_i\lambda_i\ \|\frac{1}{m}\BPhi(\BU)\boldsymbol{\varepsilon}_i\|^2
    \leq &\frac{4k_\nu(0)S_\phi}{m}.
\end{align*}


So we have the following estimation of $A_{21}$
\[A_{21}\leq A_{11}\frac{S_\phi\overline{\lambda}}{m\underline{\lambda}^2}4k_\nu(0)\leq C_5\frac{S_\phi\overline{\lambda}^3}{m\underline{\lambda}^4}.\]

For $A_{22}$, 
\begin{align*}
    \|\CalK_{22}\|^2_{L^2(\CalX\times \CalX)}\leq &\int_\CalX\sum_i\lambda_i\left|\CalP_m[e_i](\Bx)\right|^2d\BFx\int_\CalX\sum_i\lambda_i\left|\CalP_m\E[G_i](\BFx)-\E[G_i](\BFx)\right|^2d\BFx.
\end{align*}
So we can notice that its estimation is almost the same as that of $A_{11}\cdot A_{12} $. We only need to replace the eigenfunctions $\{e_i\}$ by the expectation  $\{\E G_i(\cdot)\}$ and calculate the infinite sum of $\left|\CalP\E[G_i](\BFx)-\E[G_i](\BFx)\right|^2$ as follows:
\begin{align*}
    &\sum_i\lambda_i \left|\CalP\E[G_i](\BFx)-\E[G_i](\BFx)\right|^2=\sum_i\lambda_i(\bold{I}-\CalP)\E[G_i]\E[G_i](\bold{I}-\CalP)^T(\BFx,\BFx)\\
     \coloneqq &(\bold{I}-\CalP)\bar{k}(\cdot,\cdot)(\bold{I}-\CalP)^T(\BFx,\BFx),
     \end{align*}
     where:
\begin{align*}     
    &\bar{k}(\cdot,\cdot)\\
    =& \begin{cases}
        &\E\left[\sum_i\lambda_ie_i(\BB_\tau)e_i(\BB_\gamma')e^{-\kappa^2(\tau+\gamma)}\bigg|\BB_0=\cdot,\BB'_0=\cdot\right]\quad\text{(Dirichlet)}\\
        & \E\left[\sum_i\lambda_i\int_{\Real_+^2}e^{-\kappa^2(t+\tau)-\int_0^tc(\BB_s)dL_s-\int_0^\tau c(\BB'_s)dL'_s}e_i(\BB_t)e_i(\BB'_\tau)dL_tdL'_\tau\bigg|\BB_0=\cdot,\BB'_0=\cdot\right]
    \end{cases}\\
    & {\quad\quad\quad\quad\quad\quad\quad\quad\quad\quad\quad\quad\quad\quad\quad\quad\quad\quad\quad\quad\quad\quad\quad\quad\quad\quad\quad\quad\quad\quad\quad\quad\quad \  \text{(Robin)}}\\
    =&\begin{cases}
        &\E\left[k_\nu(\BB_\tau,\BB'\gamma)\bigg|\BB_0=\cdot,\BB'_0=\cdot\right]\quad\text{(Dirichlet)}\\
        &\E\left[\int_{\Real_+^2}e^{-\kappa^2(t+\tau)-\int_0^tc(\BB_s)dL_s-\int_0^\tau c(\BB'_s)dL'_s}k_\nu(\BB_t,\BB'_\tau)dL_tdL'_\tau\bigg|\BB_0=\cdot,\BB'_0=\cdot\right]\  \text{(Robin)}
    \end{cases}
   % =& (\bold{I}-\CalP)\E\left[k(\BB_\tau,\BB_\gamma')e^{-\kappa^2(\tau+\gamma)}\bigg|\BB_0=\cdot,\BB'_0=\cdot\right](\bold{I}-\CalP)^T(\BFx,\BFx)\\
\end{align*}
So $A_{22}$ is bounded as:
\begin{align*}
 A_{22}
    \leq&C_3(\frac{\overline{\lambda}}{\underline{\lambda}})^2\bigg[\frac{\overline{\lambda}S_\phi}{\underline{\lambda}^2m}\left(\frac{1}{m}\sum_{j=1}\left(\bold{I}-\CalP\right)\bar{k}(\cdot,\cdot)\left(\bold{I}-\CalP\right)^T(\Bu_j,\Bu_j)\right)\\
    &+\int_\CalX\left(\bold{I}-\CalP\right) \bar{k}(\cdot,\cdot)\left(\bold{I}-\CalP\right)^T(\BFx,\BFx)d\Bx\bigg].
\end{align*}

For $A_{23}$,
\begin{align*}
    &\|\CalK_{23}\|^2_{L^2(\CalX\times\CalX)}\\ \leq &\int_\CalX \sum_i\lambda_i\left|\E[G_i(\BFx)]\right|^2d \BFx \int_\CalX \sum_i\lambda_i\left|e_i(\BFx)-\CalP_m[e_i](\BFx)\right|^2d\BFx\\
    =& \int_\CalX \sum_i\lambda_i\left|\E[G_i(\BFx)]\right|^2d \BFx  A_{12}.\\
    \end{align*}
From our previous calculations, the infinite sum is bounded as follows for both the Dirichlet and Robin boundary settings:
\begin{align*}
    & \sum_i\lambda_i\left|\E[G_i(\BFx)]\right|^2 \\
    = &\begin{cases}
        &\E\left[k_\nu(\BB_\tau,\BB_\gamma')e^{-\kappa^2(\tau+\gamma)}\bigg|\BB_0=\Bx,\BB'_0=\Bx\right]\\
        &\E\left[\int_{\Real_+^2}e^{-\kappa^2(t+\tau)-\int_0^tc(\BB_s)dL_s-\int_0^\tau c(\BB'_s)dL'_s}k_\nu(\BB_t,\BB'_\tau)dL_tdL'_\tau\bigg|\BB_0=\BFx,\BB'_0=\BFx\right]
    \end{cases}\\
    \leq &k_\nu(0).
\end{align*}
So:
\begin{align*}
    \|\CalK_{23}\|^2_{L^2(\CalX\times\CalX)}\leq & \text{vol}(\CalX)k_\nu(0) A_{12}\\
    \leq& C_6 \bigg[\frac{\overline{\lambda}S_\phi}{\underline{\lambda}^2m}\left(\frac{1}{m}\sum_{j=1}\left(\bold{I}-\CalP\right)k_\nu(\cdot,\cdot)\left(\bold{I}-\CalP\right)^T(\Bu_j,\Bu_j)\right)\\
    &+\int_\CalX\left(\bold{I}-\CalP\right)k_\nu(\cdot,\cdot)\left(\bold{I}-\CalP\right)^T(\BFx,\BFx)d\Bx\bigg],
\end{align*}
where the last line is from our precious calculations for $A_{12}$.

Substitute the identities of $A_{21}$, $A_{22}$, and $A_{23}$ into \eqref{eq_pf_general_convergence_8}, we then have:
\begin{align}
    &\|\CalK_2\|^2_{L^2(\CalX\times\CalX)}\label{eq_pf_general_convergence_9}\\
    \leq & C_7\bigg[\frac{S_\phi\overline{\lambda}^3}{m\underline{\lambda}^4}\nonumber\\
    &+(\frac{\overline{\lambda}}{\underline{\lambda}})^2\left(\frac{\overline{\lambda}S_\phi}{\underline{\lambda}^2m}\left[\frac{1}{m}\sum_{l=1}\left(\bold{I}-\CalP\right)\bar{k}\left(\bold{I}-\CalP\right)^T(\Bu_l,\Bu_l)\right]+\int_\CalX\left(\bold{I}-\CalP\right)\bar{k}\left(\bold{I}-\CalP\right)^T(\BFx,\BFx)d\Bx\right)\nonumber\\
    +& \left(\frac{\overline{\lambda}S_\phi}{\underline{\lambda}^2m}\left[\frac{1}{m}\sum_{l=1}\left(\bold{I}-\CalP\right)k_\nu(\cdot,\cdot)\left(\bold{I}-\CalP\right)^T(\Bu_l,\Bu_l)\right]+\int_\CalX\left(\bold{I}-\CalP\right)k_\nu(\cdot,\cdot)\left(\bold{I}-\CalP\right)^T(\BFx,\BFx)d\Bx\right)\bigg]. \nonumber
\end{align}

{\large \textbf{Estimation of $A_3$}}: For $A_3$, first decompose $\CalK_3$ as follows:
\begin{align*}
    \CalK_3(\Bx,\Bx')=\CalK_{31}(\Bx,\Bx')+\CalK_{32}(\Bx,\Bx')+\CalK_{33}(\Bx,\Bx')+\CalK_{34}(\Bx,\Bx'),
\end{align*}
where:
\begin{align*}
    &\CalK_{31}(\Bx,\Bx')=\sum_i\lambda_i\left(\CalP_m[G_i](\Bx)\CalP_m[G_i](\BFx')-\CalP_m[G_i](\BFx)\CalP_m\E[G_i](\BFx')\right)\\
    &\CalK_{32}(\Bx,\Bx')=\sum_i\lambda_i\left(\CalP_m[G_i](\BFx)\CalP_m\E[G_i](\BFx')-\CalP_m\E[G_i](\BFx)\CalP_m\E[G_i](\Bx')\right)\\
    &\CalK_{33}(\Bx,\Bx')=\sum_i\lambda_i\left(\CalP_m\E[G_i](\BFx)\CalP_m\E[G_i](\BFx')-\E[G_i(\BFx)]\CalP_m\E[G_i](\BFx')\right)\\
    &\CalK_{34}(\Bx,\Bx')=\sum_i\lambda_i\left(\E[G_i(\BFx)]\CalP_m\E[G_i](\BFx')-\E[G_i(\BFx)]\E[G_i](\BFx')\right).
\end{align*}
So $A_3$ can be further decomposed as:
\begin{align}
    A_3/8\leq &\underbrace{ \|\CalK_{31}\|_{L^2(\CalX\times\CalX)}^2 }_{A_{31}}+\underbrace{ \|\CalK_{32}\|_{L^2(\CalX\times\CalX)}^2 }_{A_{32}} +\underbrace{ \|\CalK_{33}\|_{L^2(\CalX\times\CalX)}^2 }_{A_{33}}+\underbrace{ \|\CalK_{34}\|_{L^2(\CalX\times\CalX)}^2 }_{A_{34}}.\label{eq_pf_general_convergence_10}
\end{align}
For $A_{31}$:
\begin{align*}
    \|\CalK_{31}\|_{L^2(\CalX\times\CalX)}^2\leq &\int_{\CalX}\sum_i\lambda_i\left|\CalP_m[G_i](\BFx)\right|^2d\BFx\int_{\CalX}\sum_i\lambda_i\left|\CalP_m[ \E[G_i]-G_i](\BFx)\right|^2d\BFx.\\
\end{align*}
Notice that the second multiplier on the right hand side has been estimated in the calculations for $A_{21}$:
\begin{equation}
    \int_{\CalX}\sum_i\lambda_i\left|\CalP_m[ \E[G_i]-G_i](\BFx)\right|^2d\BFx \leq C_8\frac{S_\phi\overline{\lambda}}{m\underline{\lambda}^2}. \label{eq_pf_general_convergence_11}
\end{equation}
We only need to estimate the first multiplier:
\begin{align*}
   &\int_{\CalX}\sum_i\lambda_i\left|\CalP_m[G_i](\BFx)\right|^2d\BFx\\
   =&\int_\CalX \BPhi(\BFx)^T\left(\frac{1}{m}\sum_iM_i\right)^{-1}\frac{\BPhi(\BU)}{m}\bold{H}\frac{\BPhi(\BU)^T}{m}\left(\frac{1}{m}\sum_{i=1}M_i\right)^{-1}\BPhi(\Bx)d\BFx\\
    =&\text{Tr}\left[\frac{1}{m}{\BPhi(\BU)^T}\left(\frac{1}{m}\sum_{i=1}M_i\right)^{-1}M\left(\frac{1}{m}\sum_iM_i\right)^{-1}{\BPhi(\BU)}\frac{\bold{H}}{m}\right]\\
    \leq & \text{Tr}\left[\frac{\bold{H}}{m}\right]\lambda_{max}\left(\frac{1}{m}{\BPhi(\BU)^T}\left(\frac{1}{m}\sum_{i=1}M_i\right)^{-1}M\left(\frac{1}{m}\sum_iM_i\right)^{-1}{\BPhi(\BU)}\right)\\
    \leq & \text{Tr}\left[\frac{\bold{H}}{m}\right] C_1\left(\frac{\overline{\lambda}}{\underline{\lambda}}\right)^2,
\end{align*}
where for the Dirichlet boundary setting:
\begin{align*}
    \bold{H}_{ij}=&\sum_i\lambda_ie_i(\BB^{(i)}_{\tau_i})e_i(\BB^{(j)}_{\tau_j})e^{-\kappa^2(\tau_i+\tau_j)}\bigg|\BB^{(i)}_0=\Bu_i,\BB^{(j)}_0=\Bu_j\\
    =& k_\nu(\BB^{(i)}_{\tau_i},\BB^{(j)}_{\tau_j})e^{-\kappa^2(\tau_i+\tau_j)}\bigg|\BB^{(i)}_0=\Bu_i,\BB^{(j)}_0=\Bu_j\leq k(0).
\end{align*}
Similarly, for the Robin boundary setting:
\begin{align*}
    \bold{H}_{ij}
    =&\int_{\Real_+^2}e^{-\kappa^2(t+\tau)-\int_0^tc(\BB_s)dL_s-\int_0^\tau c(\BB'_s)dL'_s}k_\nu(\BB_t,\BB'_\tau)dL_tdL'_\tau\bigg|\BB_0=\Bu_i,\BB'_0=\Bu_j\leq k(0).
\end{align*}

In both cases, we have $\text{Tr}\left[\frac{\bold{H}}{m}\right]\leq k(0)$ so we can obtain:
\begin{align}
   &\int_{\CalX}\sum_i\lambda_i\left|\CalP_m[G_i](\BFx)\right|^2d\BFx
    \leq  C_1\left(\frac{\overline{\lambda}}{\underline{\lambda}}\right)^2.\label{eq_pf_general_convergence_12}
\end{align}
So $A_{31}$ is bounded as:
\[A_{31}\leq C_9\frac{S_\phi\overline{\lambda}^3}{m\underline{\lambda}^4}.\]

For $A_{32}$,
\begin{align*}
   &\|\CalK_{32}\|^2_{L^2(\CalX\times\CalX)}\\ \leq&\int_\CalX\sum_i\lambda_i\left|\CalP_m\E[G_i](\BFx)\right|^2d\BFx\int_\CalX\sum_i\lambda_i\left(\CalP_m[G_i](\BFx)-\CalP_m\E[G_i](\BFx)\cal\right)^2d\Bx \\
  \leq&\E_{\{G_i\}}\left[\int_\CalX\sum_i\lambda_i\left|\CalP_m[G_i](\BFx)\right|^2d\BFx\right]\int_\CalX\sum_i\lambda_i\left(\CalP_m[G_i](\BFx)-\CalP_m\E[G_i](\BFx)\cal\right)^2d\Bx \\
  \leq & C_1\left(\frac{\overline{\lambda}}{\underline{\lambda}}\right)^2\int_\CalX\sum_i\lambda_i\left(\CalP_m[G_i](\BFx)-\CalP_m\E[G_i](\BFx)\cal\right)^2d\Bx \\
  \leq & A_{31},
\end{align*}
where the fourth line is from \eqref{eq_pf_general_convergence_12} and the last line is from \eqref{eq_pf_general_convergence_11}.

For $A_{33}$:
\begin{align*}
     &\|\CalK_{33}\|^2_{L^2(\CalX\times\CalX)}\\ 
     \leq &    \int_\CalX\sum_i\lambda_i\left|\CalP_m\E[G_i](\Bx)\right|^2d\Bx \int_\CalX\sum_i\lambda_i\left(\CalP_m\E[G_i](\Bx)-\E[G_i](\Bx)\right)^2d\Bx  \\
     \leq &  C_1\left(\frac{\overline{\lambda}}{\underline{\lambda}}\right)^2 \int_\CalX\sum_i\lambda_i\left(\CalP_m\E[G_i](\Bx)-\E[G_i](\Bx)\right)^2d\Bx  \\
     \leq & C_1 \left(\frac{\overline{\lambda}}{\underline{\lambda}}\right)^2 \left(\frac{\overline{\lambda}S_\phi}{\underline{\lambda}^2m}\left[\frac{1}{m}\sum_{l=1}\left(\bold{I}-\CalP\right)\bar{k}\left(\bold{I}-\CalP\right)^T(\Bu_l,\Bu_l)\right]+\int_\CalX\left(\bold{I}-\CalP\right)\bar{k}\left(\bold{I}-\CalP\right)^T(\BFx,\BFx)d\Bx\right),
\end{align*}
where the last line is from the estimation of $A_{22}$.

For $A_{44}$, calculations similar to what we did for $A_{22}$ and $A_{23}$ gives:
\begin{align*}
    &\|\CalK_{33}\|^2_{L^2(\CalX\times\CalX)}\\ 
     \leq &  \int_{\CalX}\sum_i\lambda_i \left|\E [G_i(\BFx)]\right|^2d\BFx \int_\CalX\sum_i\lambda_i\left(\CalP_m\E[G_i](\Bx)-\E[G_i](\Bx)\right)^2d\Bx \\
     %= & \int_\CalX \E \left[k(\BB_\tau,\BB'_\gamma)|\BB_0=\BB'_0=\BFx\right]d\BFx  \int_\CalX\sum_i\lambda_i\left(\CalP_m\E[G_i](\Bx)-\E[G_i](\Bx)\right)^2d\Bx  \\
     \leq & k_\nu(0)\text{vol}(\CalX)\left(\frac{\overline{\lambda}S_\phi}{\underline{\lambda}^2m}\left[\frac{1}{m}\sum_{l=1}\left(\bold{I}-\CalP\right)\bar{k}\left(\bold{I}-\CalP\right)^T(\Bu_l,\Bu_l)\right]+\int_\CalX\left(\bold{I}-\CalP\right)\bar{k}\left(\bold{I}-\CalP\right)^T(\BFx,\BFx)d\Bx\right).
\end{align*}
Substitute the identities of $A_{31}$, $A_{32}$, $A_{33}$ and $A_{34}$ into \eqref{eq_pf_general_convergence_10}, we then have:
\begin{align}\label{eq_pf_general_convergence_13}
    &\|\CalK_3\|^2_{L^2(\CalX\times\CalX)}\leq  C_{10}\bigg[\frac{S_\phi\overline{\lambda}^3}{m\underline{\lambda}^4}+(\frac{\overline{\lambda}}{\underline{\lambda}})^2\big(\frac{\overline{\lambda}S_\phi}{\underline{\lambda}^2m}\left[\frac{1}{m}\sum_{l=1}\left(\bold{I}-\CalP\right)\bar{k}\left(\bold{I}-\CalP\right)^T(\Bu_l,\Bu_l)\right]\\
    &+\int_\CalX\left(\bold{I}-\CalP\right)\bar{k}\left(\bold{I}-\CalP\right)^T(\BFx,\BFx)d\Bx\big)\bigg]\nonumber.
\end{align}
Substitute   \eqref{eq_pf_general_convergence_7}, \eqref{eq_pf_general_convergence_9}, and \eqref{eq_pf_general_convergence_13} into \eqref{eq_pf_general_convergence_5}, we have:
\begin{align}
   &\|\hat{k}_{m,Q}-k\|_{L^2(\CalX\times\CalX)}^2\leq C\left(\frac{\overline{\lambda}}{\underline{\lambda}}\right)^2\bigg[\frac{\overline{\lambda}S_\phi}{\underline{\lambda}^2m}\left(1+\frac{1}{m}\sum_{l=1} (\bold{I}-\CalP)[k_\nu+\bar{k}](\bold{I}-\CalP)^T(\Bu_l,\Bu_l) \right)\nonumber\\
    &+\int_\CalX(\bold{I}-\CalP)[k_\nu+\bar{k}](\bold{I}-\CalP)^T(\Bx,\Bx)d\Bx\bigg].
\end{align}
By applying Chebyshev's inequality on the averaged $\frac{1}{m}\sum_{l=1} (\bold{I}-\CalP)[k_\nu+\bar{k}](\bold{I}-\CalP)^T(\Bu_l,\Bu_l)$, together with the probability of the event $\CalU_*$, we have the following bound with probability at least $1-v_*-p(\sqrt{2}e^{-1/2})^{m\underline{\lambda}/S_\phi}-p(\frac{e^{1/2}}{(3/2)^{3/2}})^{m\overline{\lambda}/S_\phi}$:
\begin{equation*}
    \|\hat{k}_{m,Q}-k\|_{L^2(\CalX\times\CalX)}^2\leq C \left(\frac{\overline{\lambda}}{\underline{\lambda}}\right)^2 \left[\left(\frac{\overline{\lambda}S_\phi}{\underline{\lambda}^2m}+1\right)\int_\CalX(\bold{I}-\CalP)[k_\nu+\bar{k}](\bold{I}-\CalP)^T(\Bx,\Bx)d\Bx+\frac{\overline{\lambda}S_\phi}{\underline{\lambda}^2m}\right],
\end{equation*}
where $v_*=\int_\CalX\left|(\bold{I}-\CalP)[k+\bar{k}](\bold{I}-\CalP)^T(\Bx,\Bx)\right|^2d\Bx$. In our analysis, the statistical error $\frac{\overline{\lambda}S_\phi}{\underline{\lambda}^2m}=o(1)$, otherwise the overall error $ \|\hat{k}_{m,Q}-k\|_{L^2(\CalX\times\CalX)}^2$ blows up. We thus have the final result \eqref{eq:convergence_general}.
\endproof



\iffalse
\subsubsection{Proof of Theorem \ref{thm:mat_eigen_regression}}

\proof[Proof of Theorem \eqref{thm:mat_eigen_regression}]
The Gram matrix $M=[\langle e_i,e_j\rangle_{L^2}]=\delta_{ij}$ is the identity matrix so $\underline{\lambda}=\overline{\lambda}=1$.  Also, $S_\phi\leq pk(0)^2$ obviously. We now prove the approximation error. 

For the approximation error  of $k=k_\nu$, we have
\begin{align*}
    \int_\CalX\left(\BI-\CalP\right)k\left(\BI-\CalP\right)(\Bx,\Bx)d\Bx
    =&\int_\CalX\sum_i\lambda_i\left|\left(\BI-\CalP\right)[e_i](\Bx)\right|^2d\Bx
    =\sum_{i>p}\lambda_i =\CalO(p^{-2\nu/d})
\end{align*}
where the last equality is from Lemma \ref{lem:eigen_mat}.

For the approximation error of $\bar{k}$, we have
\begin{align*}
    &\int_\CalX\left(\BI-\CalP\right)\bar{k}\left(\BI-\CalP\right)^T(\Bx,\Bx)d\Bx\\
    =&\int_\CalX\left(\BI-\CalP\right)\left[\sum_i\lambda_i\E[e_i(\BB_\tau)e^{-\kappa^2\tau}|\BB_0=\BFx]\E[e_i(\BB_\tau)e^{-\kappa^2\tau}||\BB_0=\BFx]\right]\left(\BI-\CalP\right)^T(\Bx,\Bx)d\Bx\\
    =& \sum_i\lambda_i \int_\CalX\left|\sum_{j>p} \langle e_j, \E[e_i(\BB_\tau)e^{-\kappa^2\tau}||\BB_0=\cdot]\rangle_{L^2}e_j(\BFx)\right|^2d\BFx\\
    \leq &\sum_i\lambda_i \sum_{j>p} \langle e_j, \E[e_i(\BB_\tau)e^{-\kappa^2\tau}||\BB_0=\cdot]\rangle_{L^2}^2\\
    =& \sum_{j>p} \int_{\CalX\times\CalX} e_j(\BFx)\bar{k}(\BFx,\BFx')e_j(\BFx')d\BFx\BFx'\\
    =& \sum_{j>p} \int_{\CalX\times\CalX} e_j(\BFx)\int_{\CalX\times\CalX}G(\BFx,\Bs)k(\Bs,\Bu)G(\Bu,\Bx')d\Bs d\Bu e_j(\BFx')d\BFx\BFx'\\
    \leq & C\sum_{j>p}\lambda_j\leq \CalO(p^{-2\nu/d})
\end{align*}
where $G$ is the Green function of the operator $(\kappa^2-\bigtriangleup)$, the second-to-last line is because $\bar{k}$ is the solution to 
\[(\kappa^2-\bigtriangleup)\bar{k}(\kappa^2-\bigtriangleup)^T=k, \quad \bar{k}(\partial\CalX,\cdot)=0,\]
and the last line is because the eigenvalues of $(\kappa^2-\bigtriangleup)$ on $\CalX$ is bounded below so
\begin{align}
    &\int_{\CalX\times\CalX} e_j(\BFx)\int_{\CalX\times\CalX}G(\BFx,\Bs)k(\Bs,\Bu)G(\Bu,\Bx')d\Bs d\Bu e_j(\BFx')d\BFx\BFx'\nonumber\\
    \leq & \lambda_{min}(\kappa^2-\bigtriangleup) ^2\int_{\CalX\times\CalX} e_j(\BFx)\left[(\kappa^2-\bigtriangleup)\int_{\CalX\times\CalX}G(\BFx,\Bs)k(\Bs,\Bu)G(\Bu,\Bx')d\Bs d\Bu(\kappa^2-\bigtriangleup)^T\right]  e_j(\BFx')d\BFx\BFx'\nonumber\\
    = & \lambda_{min}(\kappa^2-\bigtriangleup)^2\int_{\CalX\times\CalX} e_j(\BFx) k(\Bx,\Bx') e_j(\BFx')d\BFx\BFx'\leq C\lambda_j.\label{eq:mat_eigen_regression_1}
\end{align}
Where the second line is from the definition of Green function $(\kappa^2-\bigtriangleup)G(\BFx,\Bs)=\delta_\Bs$ and the last line is because $e_j$ is the eigenfunction of the Mat\'ern kernel $k$.

For the probability lower bound. Firstly,
\begin{align*}
    v
    _*=&\int_\CalX \left| (\BI-\CalP)[k+\bar{k}](\BI-\CalP)^T(\Bx,\Bx)\right|^2d\Bx\\
    =&\int_\CalX\left|\sum_{j>p}e_j(\Bx)^2\left(\lambda_j+\int_{\CalX\times\CalX} e_j(\BFx)\int_{\CalX\times\CalX}G(\BFx,\Bs)k(\Bs,\Bu)G(\Bu,\Bx')d\Bs d\Bu e_j(\BFx')d\BFx\BFx' \right)\right|^2d\Bx\\
    \leq & C_1 \int_\CalX \left|\sum_{j>p}\lambda_j e_j(\Bx)^2\right|^2d\Bx \leq  C_1\left(\sum_{j>p}\lambda_j\right)\sum_{j>p}\lambda_j\int_\CalX e_j^4(\BFx)d\BFx\\
    \leq & C_2 p^{-2\nu/d}\sum_{j>p}\lambda_j\|e_j\|^{d/\alpha}_{\CalH_\nu}\|e_j\|_{L^2(\CalX)}^{4-d/\alpha}
    = C_2p^{-2\nu/d}\sum_{j>p}j^{-2\nu/d}=C_3 p^{-4\nu/d+1}
\end{align*}
where the third line is from \eqref{eq:mat_eigen_regression_1} and Jensen's inequality and the last line is from Lemma \ref{lem:RKHS_interpolation inequality} and our assumption that $\nu>d/2$.

Secondly,
\begin{align*}
    &p(\sqrt{2}e^{-1/2})^{m\underline{\lambda}/S_\phi} +p\left(\frac{e^{1/2}}{(3/2)^{3/2}}\right)^{m\overline{\lambda}/S_\phi}\leq 2p\left(\frac{e^{1/2}}{(3/2)^{3/2}}\right)^{m\overline{\lambda}/S_\phi}\\
    \leq & 2p \left(e^{-0.1}\right)^{m\overline{\lambda}/S_\phi}\leq \exp\left[-\frac{m}{10p}+\log 2p\right]\leq \exp\left[-C_2(\frac{m}{p}+\log p)\right]
\end{align*}


Let $p=\CalO(m^{\frac{1}{2\nu/d+1}})$, we can have the results.

\endproof

\fi



\section{Proof of Theorem \ref{thm:b_spline_regression}}
\label{pf:b_spline_regression}

We first present the  following lemma on important properties of the Mat\'ern kernel. These are well-known results and can be found in standard references (e.g., \cite[Sec 1.2.1]{yang2022graph}). 
\begin{lemma}
    \label{lem:eigen_mat}
    The Mat\'ern kernel $k_\nu$ has the following eigenfunction decomposition on the compact domain $\CalX$:
    \begin{equation}
        \label{eq:eign_mat}
        k_\nu(\BFx,\BFx')=\sum_i\lambda_ie_i(\BFx)e_i(\BFx'),\quad \Bx,\Bx'\in\CalX.
    \end{equation}
    Here, the eigenvalues $\lambda_i=\CalO(i^{-(2\nu/d+1)})$, the eigenfunctions $\{e_i\}$ are mutually orthogonal under the $L^2$ inner product $\langle\cdot,\cdot\rangle_2$ (see below), and the reproducing kernel Hilbert space (RKHS) inner product $\langle\cdot,\cdot\rangle_{k_\nu}$ induced by the Mat\'ern kernel $k_\nu$ is given by:
    \begin{align*}
        &\langle e_i,e_j\rangle_{L^2}\coloneqq\int_\CalX e_i(\Bs)e_j(\Bs)d\Bs=\delta_{ij},\ \langle e_i,e_j\rangle_{k_\nu}=\lambda_i^{-1}\delta_{ij},
    \end{align*}
   where $\delta_{ij}$ is the Kronecker delta. Moreover, the RKHS induced by $k_\nu$ is equivalent to the Sobolev space with smoothness order  $\alpha=\nu+d/2$.
\end{lemma}
\noindent \cite{Wendland10,paulsen2016introduction} provides further details on this RKHS connection.

The following Lemma is a Sobolev inequality for the RKHS induced by Mat\'ern-$\nu$ kernel, which is also a Sobolev space with fractional order of smoothness $\alpha=\nu+d/2$.
\begin{lemma}
    \label{lem:RKHS_interpolation inequality}
    For any function $f$ in the RKHS induced by Mat\'ern-$\nu$ kernel $k_\nu$ on $\CalX$, we have:
    \[\|f\|_{L^4(\CalX)}\leq C\|f\|_{\CalH_\nu(\CalX)}^{\frac{d}{4\alpha}}\|f\|_{L^2(\CalX)}^{1-\frac{d}{4\alpha}},\quad \alpha=\nu+d/2,\]
    where $C$ is some constant independent of $f$.
\end{lemma}
\proof
Let $\alpha^*=\lfloor\alpha\rfloor$ be the largest integer no greater than $\alpha$. We can use the Gagliardo-Nirenberg inequality (Sobolev inequality for Sobolev space with integer order of smoothness, see \cite{nirenberg1959elliptic,gagliardo1959ulteriori}) to obtain:
\[\|f\|_{L^4(\CalX)}\leq C_1  \|f\|_{\CalH^{\alpha^*}(\CalX)}^{\frac{d}{4\alpha^*}}\|f\|_{L^2(\CalX)}^{1-\frac{d}{4\alpha^*}},\]
where  $\|f\|_{\CalH^{\alpha^*}(\CalX)}$ is the Sobolev norm of $f$ with smoothness order $\alpha^*$. It  can be defined via the following RKHS defined on the whole space $\Real^d$:
\begin{align*}
    &\|f\|_{\CalH^{\alpha^*}(\CalX)}=\min_{h\in\CalH^{\alpha^*}(\Real^d)}\|h\|_{\CalH^{\alpha^*}(\Real^d)},\quad h=f,\quad \text{on}\ \CalX\\
    \text{where}\ &\|h\|_{\CalH^{\alpha^*}(\Real^d)}=\int_{\Real^d}\left|\CalF[h](\Bomega)\right|^2(\kappa+\|\Bomega\|^2)^{\alpha^*} d\Bomega.
\end{align*}
Let $h^*=\arg \min_{h\in\CalH^{\alpha^*}(\Real^d)}\|h\|_{\CalH^{\alpha^*}(\Real^d)}$ such that $h^*=f$ on $\CalX$. From the norm equivalence between RKHS and Sobolev spaces (see \cite{Wendland10,Evans15}), we have:
\begin{align*}
    &C_2\|f\|_{{\CalH^{\alpha^*}(\CalX)}}\leq  \|h^*\|_{\CalH^{\alpha^*}(\Real^d)} \leq C_3\|f\|_{{\CalH^{\alpha^*}(\CalX)}}\\
    &C_4\|f\|_{{L^2(\CalX)}}\leq  \|h^*\|_{L^2(\Real^d)} \leq C_5\|f\|_{{L^2(\CalX)}},
\end{align*}
where $C_2-C_5$ are some constants independent of $f$. Then the Gagliardo-Nirenberg inequality becomes:
\begin{equation}
    \|f\|_{L^4(\CalX)}\leq C_6  \|h^*\|_{\CalH^{\alpha^*}(\Real^d)}^{\frac{d}{4\alpha^*}}\|f\|_{L^2(\CalX)}^{1-\frac{d}{4\alpha^*}}.\label{eq:RKHS_interpolation_inequality_1}
\end{equation}
Let $p=\alpha/\alpha^*\geq 1$. The RKHS norm in the above equation can be written as:
\begin{align}
    \|h^*\|_{\CalH^{\alpha^*}(\Real^d)}^2=&\int_{\Real^d}(\kappa+\|\Bomega\|^2)^{\alpha^*}|\CalF[h^*](\Bomega)|^2 d\Bomega\nonumber\\
    =& \int_{\Real^d}(\kappa+\|\Bomega\|^2)^{\alpha^*}|\CalF[h^*](\Bomega)|^{2/p}|\CalF[h^*](\Bomega)|^{2-2/p} d\Bomega\nonumber\\
    \leq & \left(\int_{\Real^d}(\kappa+\|\Bomega\|^2)^{\alpha}||\CalF[h^*](\Bomega)|^{2} d\Bomega\right)^{\alpha^*/\alpha}\left(\int_{\Real^d}|\CalF[h^*](\Bomega)|^{2}d\Bomega\right)^{1-\alpha^*/\alpha}\nonumber\\
    \leq & C_7 \left(\int_{\Real^d}(\kappa+\|\Bomega\|^2)^{\alpha}||\CalF[h^*](\Bomega)|^{2} d\Bomega\right)^{\alpha^*/\alpha}\|f\|_{L^2(\CalX)}^{2-2\alpha^*/\alpha} \label{eq:RKHS_interpolation_inequality_2},
\end{align}
where the third line is from H\"older's inequality. Substituting \eqref{eq:RKHS_interpolation_inequality_2} into \eqref{eq:RKHS_interpolation_inequality_1} and applying the RKHS norm equivalence identities, we can have the final result.
\endproof




The following lemma is a restatement of  \cite[Theorem 1.2]{dung2011optimal} or \cite[Lemma 3.1]{dung2011adaptive} using the fact that the RKHS $\CalH_\nu(\CalX)$ induced by Mat\'ern-$\nu$ kernel $k_\nu$  on $\CalX$ is a Sobolev space with (non-integer) smoothness order $\alpha=\nu+d/2$.
\begin{lemma}
    \label{lem:spline_approx_err}
     Let $\CalP$ be the projection onto the set of B-spline $F_l$ with $|F_l|=p=\CalO(2^{ld})$ and smoothness $s-1/2>\alpha$. Then for any function $f\in\CalH_\nu$, we have the following $L^r$ error for any $r<\infty$:
    \[\| f-\CalP f\|_{L^r}\leq C \min\{p^{-\frac{\alpha}{d}+\max\{1/2-1/r,0\}}\|f\|_{\CalH_\nu},\|f\|_r\},\]
    where $C$ is some universal constant independent of $f$ and $F_l$.
\end{lemma}
The follow lemma provides estimation for the eigenvalues of the Gram matrix induced by B-splines. It is a special case of the general de Boor's conjecture proved in \cite{shadrin2001norm}.
\begin{lemma}
\label{lem:de_Boor}
    Let B-splines with smoothness $s<\infty$ be the set of basis function $\{\phi_i\}_{i=1}^p=F_l$. Let $M=[\langle \phi_i,\phi_j\rangle_{L^2}]_{ij}$ be the Gram matrix. Then:
    \[\frac{\lambda_{\rm max}(M)}{\lambda_{\rm min}(M)}=\CalO(1)\quad\text{and}\quad \lambda_{\rm min}(M)\geq \CalO(p^{-1}).\]
\end{lemma}
\proof
The entries of Gram matrix can be labeled using the indices $(\bold{j},\bold{u})$:
\[M_{\bold{j},\bold{u}}=\langle M^d_{\bold{l},\bold{j}},M^d_{\bold{l},\bold{u}}\rangle_{L^2}.\]
Because the B-spline is in product form $M^d_{\bold{l},\bold{j}}=\prod_{i=1}^dN_s(2^{l_i}x_i-j_i)$, each entry of $M$ is also the product of $d$ one-dimensional $L^2$ inner products:
\begin{align*}
    M_{\bold{j},\bold{u}}=\prod_{i=1}^d \langle N_s(2^{l}\cdot -j_i),N_s(2^{l}\cdot -u_i)\rangle_{L^2}.
\end{align*}
Therefore, the Gram matrix $M$ is the Kronecker product of $d$ Gram matrices $M^{(i)}$, i.e., $M=\bigotimes_{i=1}^dM^{(i)}$, with:
\begin{align*}
    M^{(i)}_{ju}=&\langle N_s(2^{l}\cdot -j),N_s(2^{l}\cdot -u) \rangle_{L^2},\quad,j,u\in\{-s,-s+1,\cdots,2^l-1,2^l\}.
\end{align*}
Also, the value of $M^{(i)}_{ju}$ only depends on $|j-u|$ and $M^{(i)}_{ju}=0$ for any $|j-u|>s$:
\begin{align*}
    \int_0^1  N_s(2^{l}x -j)N_s(2^{l}x -u)dx= &2^{-l}\int_{-j}^{2^l-j}N_s(x)N_s(x-(u-j))dx.
\end{align*}

The matrix $\int_{-j}^{2^l-j}N_s(x)N_s(x-(u-j))dx$ is  de Boor $L^2$ projection matrix and its maximum and minimum eigenvalues are all independent of the number of grid points $\{-s,-s+1,\cdots,2^l-1,2^l\}$ according to \cite{shadrin2001norm}. 


Therefore, we have:
\[\frac{\lambda_{\max}(M^{(i)})}{\lambda_{\min}(M^{(i)})}=\CalO(1),\quad \lambda_{\min}(M(i))\geq 2^{-l}.\]
This leads to the final result:
\[\frac{\lambda_{\max}(M)}{\lambda_{\min}(M)}=\prod_{i=1}^d \frac{\lambda_{\max}(M^{(i)})}{\lambda_{\min}(M^{(i)})}=\CalO(1)  \quad\text{and}\quad \lambda_{\min}(M)=\prod_{i=1}^d \lambda_{\min}(M^{(i)})\geq \CalO(2^{-dl})=\CalO(p^{-1} ).\]

\proof[Proof of Theorem \ref{thm:b_spline_regression}]

The conditional number $\overline{\lambda}/\underline{\lambda}$ of the Gram matrix $M$ is $\CalO(1)$ and the minimum eigenvalue $\underline{\lambda}^{-1}=\CalO(p)$ according to Lemma \ref{lem:de_Boor}. Also, the quantity $S_\phi=\CalO(1)$ because almost all $\phi_j\in F_l$ are disjoint. We now prove the approximation error.

From Lemma \ref{lem:spline_approx_err} for the projection error for functions in the RKHS  $\CalH_\nu$ induced by the Mat\'ern-$\nu$ kernel $k$, we have the following $L^2$ (i.e., $r=2$) approximation error for eigenfunction $e_i$:
\begin{align*}
    \int_\CalX\left|\left(\BI-\CalP\right)[e_i](\Bx)\right|^2d\Bx\leq C\min\{p^{-{2\nu}/{d}-1}\|e_i\|_{\CalH_\nu}^2,1\}\leq C_1  \min\{p^{-{2\nu}/{d}-1}\lambda_i^{-1},1\},
\end{align*}
where the last equality is from the properties of eigenvalues and eigenfunctions for the Mat\'ern kernel in Lemma \ref{lem:eigen_mat}. Therefore,  the approximation error of $k=k_\nu$ is:
\begin{align*}
    \int_\CalX\left(\BI-\CalP\right)k\left(\BI-\CalP\right)(\Bx,\Bx)d\Bx
    =&\sum_i\lambda_i\int_\CalX\left|\left(\BI-\CalP\right)[e_i](\Bx)\right|^2d\Bx\\
    \leq & C_1\sum_i \lambda_i  \min\{p^{-{2\nu}/{d}-1}\lambda_i^{-1},1\}\\
    \leq & C_2\left(\sum_{i\leq p}  p^{-{2\nu}/{d}-1}+\sum_{i> p}  \lambda_i\right)
    \leq  C_3 p^{-2\nu/d}.
\end{align*}
Define:
\begin{equation*}
    g_i(\BFx)\coloneqq\begin{cases}
        &\E[e_i(\BB_\tau)e^{-\kappa^2\tau}|\BB_0=\Bx]\quad (\text{Dirichlet})\\
        &\E\left[\int_{0}^{\infty}e^{-\kappa^2t-\int_0^tc(\BB_s)dL_s}e_i(\BB_t)dL_t\bigg|\BB_0=\BFx\right]\ ( \text{Robin})
    \end{cases}
\end{equation*}



Because $g_i(\Bx)$ is also the solution to the PDE from the path integral identity, we have the following identity:
\begin{align*}
    (\kappa^2-\bigtriangleup)g_i(\Bx)=0,
\end{align*}
corresponding to the Dirichlet or Robin boundary conditions. Therefor, the RKHS norm and $L^2$ norm of $g_i$ can be bounded as follows:
\begin{align*}
    \|g_i\|_2^2=&\int_\CalX \int_{\CalX\times\CalX}G(\BFx,\Bu)e_i(\Bu)e_i(\bold{v})G(\BFx,\bold{v})d\Bu d\bold{v} d\BFx\\
    \leq & \lambda^2_{min}(\kappa^2-\bigtriangleup)\|e_i\|_2^2=\CalO(1),
\end{align*}
\begin{align*}
    \langle g_i,g_i\rangle_{\CalH_\nu}=& \langle \int_\CalX G(\cdot,\Bu)e_i(\Bu)d\Bu,\int_\CalX G(\cdot,\Bu)e_i(\Bu)d\Bu \rangle_{\CalH_\nu}\\
    \leq &  \lambda^2_{min}(\kappa^2-\bigtriangleup)\langle e_i,e_i\rangle_{\CalH_\nu}=\CalO(\lambda_i^{-1}),
\end{align*}
where  $G(\BFx,\BFx')$ is the Green's function of operator $(\kappa^2-\bigtriangleup)$ corresponding to different boundary conditions such that, i.e., $(\kappa^2-\bigtriangleup)G(\BFx,\BFx')=\delta_{\BFx-\BFx'}$ 
and $ \lambda_{min}(\kappa^2-\bigtriangleup)$ denotes the minimum eigenvalues of the operator $(\kappa^2-\bigtriangleup)$ and it is bounded for both the Dirichlet and Robin boundary conditions.



%So calculations similar to \eqref{eq:mat_eigen_regression_1} shows $\|g_i\|_{\CalH_\nu}^2\leq \CalO(\lambda_i^{-1})$ abd $\|g_i\|_{L^2}^2\leq \CalO(1)$. 
Therefore, the $L^2$ approximation error of $\CalP$ for $\bar{k}$ can be computed as follows:
\begin{align*}
    &\int_\CalX\left(\BI-\CalP\right)\bar{k}\left(\BI-\CalP\right)^T(\Bx,\Bx)d\Bx\\
    %=&\int_\CalX\left(\BI-\CalP\right)\left[\sum_i\lambda_i\E[e_i(\BB_\tau)e^{-\kappa^2\tau}|\BB_0=\BFx]\E[e_i(\BB_\tau)e^{-\kappa^2\tau}||\BB_0=\BFx]\right]\left(\BI-\CalP\right)^T(\Bx,\Bx)d\Bx\\
    =& \sum_i\lambda_i \int_\CalX\left|\left(\BI-\CalP\right)[g_i](\Bx) \right|^2d\BFx\\
    \leq & C_3 \sum_i \lambda_i \min\{p^{-{2\nu}/{d}-1}\|g_i\|^2_{\CalH_\nu},1\}\leq  C_4\left(\sum_{i\leq p}  p^{-{2\nu}/{d}-1}+\sum_{i> p}  \lambda_i\right)
    \leq  C_5 p^{-2\nu/d}.
\end{align*}
For the  $v_*$ in the probability lower bound:
\begin{align}
    v
    _*=&\int_\CalX \left| (\BI-\CalP)[k+\bar{k}](\BI-\CalP)^T(\Bx,\Bx)\right|^2d\Bx\nonumber\\
    =& \int_\CalX\left|(\BI-\CalP)[\sum_i\lambda_i(e_ie_i+g_ig_i)](\BI-\CalP)^T(\Bx,\Bx)\right|^2d\Bx\nonumber\\
    =& \int_\CalX\left|\sum_i \lambda_i \left[|(\BI-\CalP)[e_i](\Bx)|^2+|(\BI-\CalP)[g_i](\Bx)|^2\right]\right|^2d\Bx\nonumber\\
    \leq & 2\int_\CalX \left(\sum_i \lambda_i\right)\left(\sum_i\lambda_i\left[|(\BI-\CalP)[e_i](\Bx)|^2+|(\BI-\CalP)[g_i](\Bx)|^2\right]^2\right)d\Bx\nonumber\\
    \leq & C_4 \sum_i\lambda_i\left(\int_\CalX  |(\BI-\CalP)[e_i](\Bx)|^4d\Bx+\int_\CalX  |(\BI-\CalP)[g_i](\Bx)|^4d\Bx\right)\nonumber\\
    \leq & C_5\sum_i\lambda_i \left(\min\{p^{-4\alpha/d+1}\|e_i\|_{\CalH_v}^4,\|e_i\|_{L^4}^4\}+\min\{p^{-4\alpha/d+1}\|g_i\|_{\CalH_v}^4,\|g_i\|_{L^4}^4\}\right) \nonumber\\
    \leq & C_5\left(\sum_{i\leq p} \lambda_i^{-1}p^{-4\alpha/d+1}+\sum_{i>p}\lambda_i^{1-\frac{d}{2\alpha }}\right)\nonumber\\
    \leq & C_6 p^{-\frac{2\nu}{d}+1}\nonumber,
\end{align}
where the second-to-last line is from Lemma \ref{lem:RKHS_interpolation inequality}, our assumption $\nu>d/2$,  and the fact that $\|e_i\|^2_{\CalH_\nu}=\lambda_i^{-1}$, $\|g_i\|^2_{\CalH_\nu}\leq \CalO(\lambda_i^{-1})$, $\|e_i\|_{L^2}=1$, and $\|g_i\|_{L^2}\leq \CalO(1)$. For the remaining two terms in our probability bound, we have:
\begin{align*}
    &p(\sqrt{2}e^{-1/2})^{m\underline{\lambda}/S_\phi} +p\left(\frac{e^{1/2}}{(3/2)^{3/2}}\right)^{m\overline{\lambda}/S_\phi}\leq 2p\left(\frac{e^{1/2}}{(3/2)^{3/2}}\right)^{m\overline{\lambda}/S_\phi}\\
    \leq & 2p \left(e^{-0.1}\right)^{m\overline{\lambda}/S_\phi}\leq \exp\left[-\frac{m}{10p}+\log 2p\right]\leq \exp\left[-C_2(\frac{m}{p}+\log p)\right].
\end{align*}
Finally, setting $p=\CalO(m^{\frac{1}{2\nu/d+1}})$, the final result is proven.
\endproof


% Notice that the closed-form solution \eqref{eq:kernel_green_1d_closed} aligns precisely with \eqref{eq:bdryMatern_path_integral_dirichlet} and \eqref{eq:bdryMatern_path_integral_robin}, where \(k_{hp}\) and \(k_{hh}\) represent the exact expectations with respect to a single Brownian motion and two Brownian motions, respectively. Since \(k^{\rm BM}_{\nu}\) incorporates boundary conditions at \(0\) and \(1\), it can be extended to construct a multi-dimensional tensor BdryMatérn kernel:
% \[
% k^{\rm TBM}_{\nu}(\BFx,\BFx')=\prod_{j=1}^d k_\nu^{\rm BM}(x_j,x_j'),
% \]
% ensuring that \(k^{\rm TBM}_{\nu}(\BFx,\BFx')\) satisfies the general boundary condition \eqref{eq:boundar_mixed_1D} on the boundaries of the cube \([0,1]^d\). 



\section{Computational Time Complexity of  Kernel Regression by B-Spline}
\label{sec:B_spline_time_complexity}
Recall that the kernel regression is written as follows:
\begin{equation}
\label{eq:kernel_reg2_appendix}
 \hat{k}_{m,Q}(\BFx, \BFx') = \sum_{q,q'=1}^Q m_{q,q'}{\phi}_q(\BFx) {\phi}_{q'}(\BFx'), \; \mathbf{M} = [m_{q,q'}]_{q,q'=1}^Q = [\BPhi(\BU)]^\dagger \hat{\BK}^{\rm C}(\BU, \BU) [\BPhi(\BU)^T]^\dagger.
\end{equation}
If the basis function $\{\phi_q\}_{q=1}^Q=F_\zeta=\{M_{\zeta, \mathbf{j}}^d, \mathbf{j}\in \mathcal{J}(\zeta)\}$ are B-splines defined as \eqref{eq:B_spline} and the support of each $M_{\zeta, \mathbf{j}}^d$ is:
\begin{equation}
\label{eq:B_spline_support_appendix}
    {\rm supp}\{M_{\zeta, \mathbf{j}}^d\}=[0,2^{-\zeta}(s+1)]^d+\bold{j}2^{-\zeta},
\end{equation}
we can have the following lemma regarding the number of non-zero elements of the vector $\{M_{\zeta, \mathbf{j}}^d(\BFx), \mathbf{j}\in \mathcal{J}(\zeta)\}$ for any $\BFx\in\CalX$:
\begin{lemma}
    \label{lem:nnz_B_spline}
     The number of non-zero element of $[\phi_1(\BFx),\cdots,\phi_Q(\BFx)]$ is at most $(2s)^d$ for any $\BFx\in\CalX$.
\end{lemma}
\begin{proof}
    Let $\CalX$ be partitioned into small, non-overlapping rectangles of the form ${[0, 2^{-\zeta}(s+1)]^d + \boldsymbol{j} 2^{-\zeta}}$ for integer vectors $\boldsymbol{j}$. Then, every point $\BFx \in \CalX$ lies in exactly one such rectangle. Without loss of generality, assume $\BFx \in [0, 2^{-\zeta}(s+1)]^d$. We then count the basis functions $\phi_q$ that are nonzero on this rectangle. According to \eqref{eq:B_spline_support_appendix}, there are at most $(2s)^d$ such functions.
\end{proof}
By Lemma \ref{lem:nnz_B_spline}, to evaluate \eqref{eq:kernel_reg2_appendix} at any given $\BFx$, it suffices to compute only those entries ${m_{q,q'}}$ for which the basis functions $\phi_q$ and $\phi_{q'}$ have overlapping supports. As a result, we only need to compute at most $(2s)^d Q$ entries of ${m_{q,q'}}$, rather than the full inverse of the matrix $\boldsymbol{M}$. This leads to the following lemma:
\begin{lemma}
    \label{lem:time_complexity_sparse_M}
     Let $\mathbf{M} = [m_{q,q'}]_{q,q'=1}^Q = [\BPhi(\BU)]^\dagger \hat{\BK}^{\rm C}(\BU, \BU) [\BPhi(\BU)^T]^\dagger.$ Define the set:
     \[\mathbf{M}^{>0}=\{m_{q,q'}:  {\rm supp}\{\phi_q\}\bigcap   {\rm supp}\{\phi_q'\}\neq \emptyset\}.\]
     Then the size $\bold{M}^{>0}$ is at most $(2s)^dQ$ and it can be computed in $\CalO((2s)^dQ^3)$ time.
\end{lemma}
\begin{proof}
    We first count the size of $\bold{M}^{>0}.$ According to Lemma \ref{lem:nnz_B_spline}, there are at most $(2s)^d$ B-spline basis functions with overlapping support with each $\phi_q$. So we can directly conclude that $|\bold{M}^{>0}|=\CalO((2s)^d)Q$. 
    
    For computational time complexity, we directly provide an algorithm to solve for $\bold{M}^{>0}$ in $\CalO((2s)^dQ^3)$ time. Define:
    \begin{align*}
        \bold{S}=\left[\BPhi(\BU)\BPhi(\BU)^T\right]^{-1},\quad \bold{R}=\BPhi(\BU) \hat{\BK}^{\rm C}(\BU, \BU) \BPhi(\BU)^T,
    \end{align*}
    so that $\forall m_{q,q'}$, $ m_{q,q'}= \bold{S}_{q,:}\bold{R}\bold{S}_{:,q'}$. Note that $\bold{S}\in\Real^{Q\times Q}$, so it can be computed in $\CalO(Q^3)$ time. Also, because of the sparsity of $ \BPhi(\BU)$, $\bold{R}$ can be computed in $\CalO(Q^3)$ time. 
    
    After we have computed $\bold{R}$, the computational time complexity for matrix-vector multiplication $m_{q,q'}=\bold{S}_{q,:}\bold{R}\bold{S}_{:,q'}$ is $\CalO(Q^2)$. In total, we only need to compute those $m_{q,q'}\in\bold{M}^{>0}$, so we need to repeat this process $(2s)^dQ$ times.

    Therefore, the total time complexity for all $m_{q,q'}\in\bold{M}^{>0}$ is $\CalO((2s)^dQ^3)$.
\end{proof}
After we can have the values of $\bold{M}_0$, we have the following efficient computation for kernel values:

\begin{theorem}
    \label{thm:time_complexity_kernel}Given $\bold{M}^{>0}$,  $\hat{k}_{m,Q}(\BFx, \BFx') $ can be computed in $(2s)^{2d}$ steps at most for any pair of $\BFx$ and $\BFx'$.
\end{theorem}
\begin{proof}
    By Lemma \ref{lem:nnz_B_spline}, for any $\BFx$ and $\BFx'$, the vectors $\BPhi(\BFx)$ and $\BPhi(\BFx')$ each contain at most $n^* = (2s)^d$ nonzero entries. Denote these nonzero basis functions as $[\phi_{q_1}, \ldots, \phi_{q_{n^*}}]$ and $[\phi_{q_1'}, \ldots, \phi_{q_{n^*}'}]$, respectively. Then:
    \[\hat{k}_{m,Q}(\BFx, \BFx')=\sum_{i=1}^{n^*}\sum_{j=1}^{n^*}m_{q_i,q_j}\phi_i(\BFx)\phi_j(\BFx).\]
    There are thus at most $(2s)^{2d}$ computations in total.
\end{proof}
\end{appendix}

\end{document}
